% ============================================================
% M3 S4 拓展册（下）· 拓区收编（115 题）
% 依据：_tmpM3S4预备0913/S4波次方案.md §3.1（221 题次口径）＋canonical键名总表.md＋
%   定稿/定稿汇总-第2章.md（拓区账）＋_tmpM3S4题源闸0914/首跑钉码.md。
% 版式：练习线 qp-m3.sty（md5 7c3930362be8a0a2bdf21bbf8ac16573，本地挂载副本）；题后紧跟答案制；
%   全线括线模（\ansblockgrayfalse 导言一次置定——「>8 行→括线模」在全括线下自然满足）。
% 组织：节域序（canonical 01→…→17）→题型族组→组内席序；组名行只给 ≥2 成员族打。
% 跳号对号档：撤席（10-拓1／11-拓2／12-拓27／拓14-13／拓14-31／16-18）＝印面号空缺不重编号，
%   注记随席（tex 注释＋值台账抄件，不入印面）；章末 19 席缺料未收编（残余项，见汇编报告）。
% ============================================================
\documentclass[fontset=none]{ctexart}
\usepackage{qp-m3}
% —— 印面逐字符号路由补钉（承练习件母版；− 走 TNR、▱ NSC 子块；禁改 sty 保 md5） ——
\xeCJKDeclareCharClass{Default}{"2212}%
\setCJKfamilyfont{nscblk}[Path=C:/提示词/工作区/字替对照-0909/variantF/fonts/,BoldFont=NSC-w600.ttf]{NSC-w500.ttf}%
\xeCJKDeclareSubCJKBlock{nscblk}{"25B1}%
\newcommand{\pxparallelogram}{{\CJKfamily{nscblk}▱}}%
\xeCJKsetup{CJKglue={\hskip 0.10em plus 0.08em minus 0.05em}}%
\ansblockgrayfalse
\renewcommand{\qpzhangming}{第二章\quad 平面解析几何}%
\renewcommand{\qpjianming}{拓展册（下）}%
% —— 拓展册局部宏（组名行＝M2 拓展册档独立行；悬挂题号＝席号字面自适应宽，撤席跳号照印） ——
\newsavebox{\thbox}%
\newcommand{\tuhao}[1]{\par\glueguard{1}%
  \sbox{\thbox}{\textbf{#1.}}%
  \setlength{\lxhang}{\dimexpr\wd\thbox+1.4mm\relax}%
  \noindent\hangindent=\lxhang\hangafter=0
  \llap{\hbox to\lxhang{\textbf{#1.}\hss}}}%
\definecolor{gray102}{HTML}{666666}%
\newcommand{\zux}[1]{\par\glueguard{2}\addvspace{2.4mm}\noindent
  {\fontsize{10.5pt}{18pt}\selectfont\heibf 【#1】}\par\addvspace{1.2mm}}%

\begin{document}

\zhangtitle{拓展册（下）}

\begin{multicols}{2}
\emergencystretch=1em

\jietitle{2.6.2　双曲线的性质（课时14·拓区）}
% 【拓14-13｜让位席（教材练A2 判重让位）不占键；【拓14-31｜双收4 撤】不设块——两席跳号留痕，注记随席（台账抄件）。】
\zux{拓展·通径直角求e}
\tuhao{01} \tieside{难}过双曲线 \(x^{2}/a^{2}-y^{2}/b^{2}=1\) 的右焦点 \(F_{2}\) 作垂直于实轴的弦 \(PQ\)，\(F_{1}\) 是左焦点，若 \(\angle PF_{1}Q=90^{\circ}\)，则双曲线的离心率是（　　　）
\bindopt {\fontsize{10.5pt}{\qpTJgLeadLiti}\selectfont \makebox[\dimexpr0.25\linewidth-0.25\lxhang-0.25em\relax][l]{\(A\)．\(\sqrt{2}\)}\makebox[\dimexpr0.25\linewidth-0.25\lxhang-0.25em\relax][l]{\(B\)．\(1+\sqrt{2}\)}\makebox[\dimexpr0.25\linewidth-0.25\lxhang-0.25em\relax][l]{\(C\)．\(2+\sqrt{2}\)}\makebox[\dimexpr0.25\linewidth-0.25\lxhang-0.25em\relax][l]{\(D\)．\(3-\sqrt{2}\)}\par}
\begin{ansblock}[2章-拓-课时14-01]
% ans:2章-拓-课时14-01
\ansitem{01}{\(B\)}
\ansnote{详解}{\(PQ\) 为通径，\(P(c,\pm b^{2}/a)\)，\(|PQ|=2b^{2}/a\)。 \(\triangle PF_{1}Q\) 中 \(\angle PF_{1}Q=90^{\circ}\)，\(F_{2}\) 为斜边 \(PQ\) 的中点，由直角三角形斜边中线定理：\(|F_{1}F_{2}|=|PQ|/2\)，即 \(2c=b^{2}/a\)，\(2ac=b^{2}=c^{2}-a^{2}\)。 两边除以 \(a^{2}\)：\(e^{2}-2e-1=0\)，\(e>1\)，\(e=1+\sqrt{2}\)，故选\(B\)。 验算：\(e=1+\sqrt{2}\) 时 \(e^{2}=3+2\sqrt{2}\)，\(e^{2}-2e-1=(3+2\sqrt{2})-2(1+\sqrt{2})-1=0\) 。}
\end{ansblock}

\zux{拓展·渐近线交点中点垂直求e}
\tuhao{02} \tieside{难}已知双曲线 \(C\)：\(x^{2}/a^{2}-y^{2}/b^{2}=1\)（\(a>0\)，\(b>0\)）的左、右焦点分别为 \(F_{1}\)，\(F_{2}\)，过 \(F_{1}\) 的直线与 \(C\) 的两条渐近线分别交于 \(A\)，\(B\) 两点．若 \(|F_{1}A|=|AB|\)，\(F_{1}B\cdot F_{2}B=0\)，则 \(C\) 的离心率为\kongda{\(2\)}．
\begin{ansblock}[2章-拓-课时14-02]
% ans:2章-拓-课时14-02
\ansitem{02}{\(2\)}
\ansnote{详解}{\(A\) 在线段 \(F_{1}B\) 上且 \(|F_{1}A|=|AB|\)，故 \(A\) 为 \(F_{1}B\) 的中点；\(O\) 为 \(F_{1}F_{2}\) 的中点，故 \(OA\) 是 \(\triangle F_{1}F_{2}B\) 的中位线：\(OA\parallel F_{2}B\)，\(|F_{2}B|=2|OA|\)。 由 \(F_{1}B\cdot F_{2}B=0\) 得 \(F_{1}B\perp F_{2}B\)，结合 \(OA\parallel F_{2}B\) 得 \(OA\perp F_{1}B\)。 于是 \(\triangle F_{1}OB\) 中，\(OA\) 既是边 \(F_{1}B\) 上的中线又是高，故 \(\triangle F_{1}OB\) 等腰：\(|OB|=|OF_{1}|=c\)，且 \(\angle AOF_{1}=\angle AOB\)。 又 \(A\)、\(B\) 分别在两条渐近线上，两渐近线关于 \(x\) 轴对称，故 \(\angle BOF_{2}=\angle AOF_{1}\)；而 \(\angle AOF_{1}+\angle AOB+\angle BOF_{2}=180^{\circ}\)，三者在上述等量关系下相等，各为 \(60^{\circ}\)。 故 \(b/a=tan60^{\circ}=\sqrt{3}\)，\(e=\sqrt{}(1+(b/a)^{2})=\sqrt{1+3}=2\)。 验算：三内角和 \(3\times 60^{\circ}=180^{\circ}\) ；\(b/a=\sqrt{3}\) 时 \(e=2\) 。}
\end{ansblock}

\zux{拓展·正六边形模型求e}
\tuhao{03} \tieside{难}已知 \(F_{1}\)，\(F_{2}\) 是双曲线 \(x^{2}/a^{2}-y^{2}/b^{2}=1\)（\(a>0\)，\(b>0\)）的左、右焦点，以 \(F_{1}F_{2}\) 为直径的圆与双曲线交于不同的四点，顺次连接焦点和这四点恰好组成一个正六边形，则该双曲线的离心率为（　　　）
\bindopt {\fontsize{10.5pt}{\qpTJgLeadLiti}\selectfont \makebox[\dimexpr0.5\linewidth-0.5\lxhang-0.25em\relax][l]{\(A\)．\(\sqrt{2}\)}\makebox[\dimexpr0.5\linewidth-0.5\lxhang-0.25em\relax][l]{\(B\)．\(\sqrt{3}\)}\par}
\bindopt {\fontsize{10.5pt}{\qpTJgLeadLiti}\selectfont \makebox[\dimexpr0.5\linewidth-0.5\lxhang-0.25em\relax][l]{\(C\)．\(\sqrt{3}+1\)}\makebox[\dimexpr0.5\linewidth-0.5\lxhang-0.25em\relax][l]{\(D\)．\((\sqrt{5}+1)/2\)}\par}
\begin{ansblock}[2章-拓-课时14-03]
% ans:2章-拓-课时14-03
\ansitem{03}{\(C\)}
\ansnote{详解}{六个顶点内接于以 \(F_{1}F_{2}\) 为直径的圆，由直径所对圆周角 \(\angle F_{1}PF_{2}=90^{\circ}\)；正六边形各边对圆心角 \(60^{\circ}\)，得 \(\angle PF_{2}F_{1}=30^{\circ}\)。 解三角形：\(|PF_{2}|=2c\cdot cos30^{\circ}=\sqrt{3}c\)，\(|PF_{1}|=2c\cdot sin30^{\circ}=c\)。 \(P\) 在双曲线上，\(|PF_{2}|-|PF_{1}|=2a \rightarrow  \sqrt{3}c-c=2a \rightarrow  e=c/a=2/(\sqrt{3}-1)=\sqrt{3}+1\)，故选\(C\)。 验算：\(2/(\sqrt{3}-1)\) 有理化 \(=2(\sqrt{3}+1)/2=\sqrt{3}+1\) ；\(P\) 在右支，长焦半径减短焦半径 。 （卷③自带详解照录，验算一致。）}
\end{ansblock}

\zux{拓展·蒙日圆交点＋焦半径比例求e}
% 图债注：「如图」指源件配图，印面无图（冻片【注】裁定：去装饰图/条件全在文字/尺寸随注——图债登记汇编报告）
\tuhao{04} \tieside{难}如图，\(F_{1}\)，\(F_{2}\) 分别是双曲线 \(x^{2}/a^{2}-y^{2}/b^{2}=1\)（\(a>0\)，\(b>0\)）的左、右焦点，点 \(P\) 是双曲线与圆 \(x^{2}+y^{2}=a^{2}+b^{2}\) 在第二象限的一个交点，点 \(Q\) 在双曲线上，且 \(F_{1}P=(1/2)F_{2}Q\)，则双曲线的离心率为（　　　）
\bindopt {\fontsize{10.5pt}{\qpTJgLeadLiti}\selectfont \makebox[\dimexpr0.5\linewidth-0.5\lxhang-0.25em\relax][l]{\(A\)．\(\sqrt{10}/2\)}\makebox[\dimexpr0.5\linewidth-0.5\lxhang-0.25em\relax][l]{\(B\)．\(\sqrt{2}\)}\par}
\bindopt {\fontsize{10.5pt}{\qpTJgLeadLiti}\selectfont \makebox[\dimexpr0.5\linewidth-0.5\lxhang-0.25em\relax][l]{\(C\)．\(\sqrt{3}\)}\makebox[\dimexpr0.5\linewidth-0.5\lxhang-0.25em\relax][l]{\(D\)．\(\sqrt{17}/3\)}\par}
\begin{ansblock}[2章-拓-课时14-04]
% ans:2章-拓-课时14-04
\ansitem{04}{\(D\)}
\ansnote{详解}{圆 \(x^{2}+y^{2}=a^{2}+b^{2}\) 的半径为 \(c\)，其直径恰为 \(F_{1}F_{2}\)，故 \(P\) 在圆上 \(\Rightarrow  \angle F_{1}PF_{2}=90^{\circ}\)（直径所对圆周角）。 设 \(|PF_{1}|=n\)，由定义 \(|PF_{2}|=2a+n\)。\(Rt\triangle PF_{1}F_{2}\) 中：\(n^{2}+(2a+n)^{2}=4c^{2} \rightarrow  2n^{2}+4an+4a^{2}=4(a^{2}+b^{2}) \rightarrow  n^{2}=2b^{2}-2an\)。 由 \(F_{1}P=(1/2)F_{2}Q\)（同向向量）得 \(|F_{2}Q|=2n\) 且 \(F_{1}P\parallel F_{2}Q\)，故 \(\angle QF_{2}F_{1}=180^{\circ}-\angle PF_{1}F_{2}\)，\(cos\angle QF_{2}F_{1}=-n/(2c)\)（\(Rt\triangle \) 中 \(cos\angle PF_{1}F_{2}=n/(2c)\)）。 由定义 \(|F_{1}Q|=2a+2n\)。\(\triangle F_{1}F_{2}Q\) 中余弦定理： \((2a+2n)^{2}=(2n)^{2}+(2c)^{2}+4n^{2} \rightarrow  4a^{2}+8an+4n^{2}=4n^{2}+4c^{2}+4n^{2} \rightarrow  4a^{2}+8an=4c^{2}+4n^{2} \rightarrow  a^{2}+2an=c^{2}+n^{2}=(a^{2}+b^{2})+n^{2} \rightarrow  2an=b^{2}+n^{2}\)，即 \(b^{2}=2an-n^{2}\)。 与 \(n^{2}=2b^{2}-2an\) 联立：\(n^{2}=2(2an-n^{2})-2an=2an-2n^{2} \rightarrow  3n^{2}=2an \rightarrow  a=(3/2)n\)； \(b^{2}=2an-n^{2}=3n^{2}-n^{2}=2n^{2}\)；\(c^{2}=a^{2}+b^{2}=(9/4+2)n^{2}=(17/4)n^{2}\)。 \(e=c/a=[(\sqrt{17}/2)n]/[(3/2)n]=\sqrt{17}/3\)，故选\(D\)。 验算：\(n=1\) 时 \(a=1.5\)、\(b=\sqrt{2}\)、\(c=\sqrt{17}/2\approx 2.06\)，\(c^{2}=a^{2}+b^{2}=2.25+2=4.25\) 。 （卷③自带详解照录校订，读数一致。）}
\end{ansblock}

\zux{拓展·圆交点斜率差求e}
\tuhao{05} \tieside{难}已知双曲线 \(C\)：\(x^{2}/a^{2}-y^{2}/b^{2}=1\)（\(a>0\)，\(b>0\)）的左、右顶点分别为 \(A\)，\(B\)，圆 \(P\)：\(x^{2}+(y-a)^{2}=2a^{2}\) 与双曲线 \(C\) 在第一象限的交点为 \(M\)，记直线 \(MA\)，\(MB\) 的斜率分别为 \(m\)，\(n\)，且 \(n-m=3\)，则双曲线的离心率为\kongda{\(\sqrt{3}\)}．
\begin{ansblock}[2章-拓-课时14-05]
% ans:2章-拓-课时14-05
\ansitem{05}{\(\sqrt{3}\)}
\ansnote{详解}{\(O(0,0)\)，\(A(-a,0)\)，\(B(a,0)\)，圆 \(P\) 的圆心为 \(P(0,a)\)，半径 \(\sqrt{2}a\)。 \(|OA|=|OB|=|OP|=a\)，故 \(O\)、\(A\)、\(P\)、\(B\) 四点共圆（圆心 \(O\)、半径 \(a\)），其中 \(AB\) 为该圆的直径（\(A\)、\(B\) 关于 \(O\) 对称），\(P\) 在圆上，由直径所对圆周角得 \(\angle APB=90^{\circ}\)。 \(M\) 在圆上且在第一象限，位于弦 \(AB\) 所对优弧，故圆周角 \(\angle AMB=(1/2)\angle APB=45^{\circ}\)。 夹角公式：\(tan\angle AMB=(n-m)/(1+nm)=1\)，\(n-m=3 \rightarrow  nm=2\)。 设 \(M(x,y)\)：\(m=y/(x+a)\)，\(n=y/(x-a)\)，\(nm=y^{2}/(x^{2}-a^{2})\)；由 \(y^{2}=b^{2}(x^{2}/a^{2}-1)\) 得 \(nm=b^{2}/a^{2}=2\)。 \(e^{2}=1+b^{2}/a^{2}=3\)，\(e=\sqrt{3}\)。 验算：\(mn=(y^{2})/((x-a)(x+a))\)，\((x-a)(x+a)=x^{2}-a^{2}=a^{2}(y^{2}/b^{2})-a^{2}=(a^{2}/b^{2})y^{2} \rightarrow  mn=b^{2}/a^{2}\) 。 （卷③自带详解照录校订：\(|OA|=|OB|=|OP|\) 与圆心角推导等价，此处按向量验。）}
\end{ansblock}

\zux{拓展·通径上张角60°求e最大值}
\tuhao{06} \tieside{难}\(l\) 是经过双曲线 \(C\)：\(x^{2}/a^{2}-y^{2}/b^{2}=1\)（\(a>0\)，\(b>0\)）焦点 \(F\) 且与实轴垂直的直线，\(A\)，\(B\) 是双曲线 \(C\) 的两个顶点，若在 \(l\) 上存在一点 \(P\)，使 \(\angle APB=60^{\circ}\)，则双曲线离心率的最大值为\kongda{\(2\sqrt{3}/3\)}．
\begin{ansblock}[2章-拓-课时14-06]
% ans:2章-拓-课时14-06
\ansitem{06}{\(2\sqrt{3}/3\)}
\ansnote{详解}{设 \(A(a,0)\)，\(B(-a,0)\)，\(P(c,y)\)（\(l\)：\(x=c\)）。\(k\_PA=y/(c-a)\)，\(k\_PB=y/(c+a)\)。 由夹角公式 \(tan\angle APB=|[k\_PA-k\_PB]/[1+k\_PA\cdot k\_PB]|=\sqrt{3}\)： \(|(2ay)/(c^{2}-a^{2})| / |1+y^{2}/(c^{2}-a^{2})|=\sqrt{3} \rightarrow \) 整理为关于 \(y\) 的方程 \(\sqrt{3}y^{2}-2ay+\sqrt{3}(c^{2}-a^{2})=0\) 有实根。 判别式 \(\Delta =4a^{2}-12(c^{2}-a^{2})=4(4a^{2}-3c^{2})\geqslant 0 \rightarrow  4a^{2}\geqslant 3c^{2} \rightarrow  e^{2}\leqslant 4/3\)，\(e\leqslant 2\sqrt{3}/3\)。 离心率最大值 \(2\sqrt{3}/3\)。 验算：\(e=2\sqrt{3}/3\) 时 \(\Delta =0\)，\(P\) 取对称轴位置 \(y=a/\sqrt{3}\) 。 （卷③自带详解照录，验算一致。）}
\end{ansblock}

\zux{拓展·共渐近线＋焦点求方程}
\tuhao{07} \tieside{难}与双曲线 \(x^{2}-2y^{2}=2\) 共渐近线且一个焦点为 \((0,-6)\) 的双曲线方程为（　　　）
\bindopt {\fontsize{10.5pt}{\qpTJgLeadLiti}\selectfont \makebox[\dimexpr0.5\linewidth-0.5\lxhang-0.25em\relax][l]{\(A\)．\(x^{2}/12-y^{2}/24=1\)}\makebox[\dimexpr0.5\linewidth-0.5\lxhang-0.25em\relax][l]{\(B\)．\(y^{2}/12-x^{2}/24=1\)}\par}
\bindopt {\fontsize{10.5pt}{\qpTJgLeadLiti}\selectfont \makebox[\dimexpr0.5\linewidth-0.5\lxhang-0.25em\relax][l]{\(C\)．\(x^{2}/24-y^{2}/12=1\)}\makebox[\dimexpr0.5\linewidth-0.5\lxhang-0.25em\relax][l]{\(D\)．\(y^{2}/24-x^{2}/12=1\)}\par}
\begin{ansblock}[2章-拓-课时14-07]
% ans:2章-拓-课时14-07
\ansitem{07}{\(B\)}
\ansnote{详解}{共渐近线设 \(x^{2}-2y^{2}=m\)（\(m\neq 0\)）。焦点 \((0,-6)\) 在 \(y\) 轴上，故焦点在 \(y\) 轴，\(m<0\)。 标准形 \(y^{2}/(-m/2)-x^{2}/(-m)=1\)：\(a^{2}=-m/2\)，\(b^{2}=-m\)，\(c^{2}=a^{2}+b^{2}=-3m/2=36 \rightarrow  m=-24\)。 方程 \(y^{2}/12-x^{2}/24=1\)，选\(B\)。 验算：\(c^{2}=12+24=36\)，\(c=\pm 6\) ，焦点含 \((0,-6)\) 。 （卷③自带详解照录，验算一致。）}
\end{ansblock}

\zux{拓展·等轴双曲线|PO|²恒等证明}
\tuhao{08} \tieside{难}设等轴双曲线 \(C\) 的中心为 \(O\)，焦点为 \(F_{1}\)，\(F_{2}\)，\(P\) 为 \(C\) 上任意一点，求证：\(|PO|^{2}=|PF_{1}|\cdot |PF_{2}|\)．
\liubai[16mm]{解答书写区}
\begin{ansblock}[2章-拓-课时14-08]
% ans:2章-拓-课时14-08
\ansitem{08}{证明见解析。}
\ansnote{详解}{证明：设 \(C\)：\(x^{2}/m^{2}-y^{2}/m^{2}=1\)（\(m>0\)），则 \(F_{1}(-\sqrt{2}m,0)\)，\(F_{2}(\sqrt{2}m,0)\)。设 \(P(x,y)\)，满足 \(x^{2}-y^{2}=m^{2}\)。 \(|PO|^{2}=x^{2}+y^{2}\)。 \(|PF_{1}|\cdot |PF_{2}|=\sqrt{(x+\sqrt\{2\}m)^{2}+y^{2}}\cdot \sqrt{(x-\sqrt\{2\}m)^{2}+y^{2}} \allowbreak =\sqrt{}\{[(x^{2}+2m^{2}+y^{2})+2\sqrt{2}mx]\cdot [(x^{2}+2m^{2}+y^{2})-2\sqrt{2}mx]\} \allowbreak =\sqrt{(x^{2}+2m^{2}+y^{2})^{2}-8m^{2}x^{2}}\)。 根号内用 \(x^{2}-y^{2}=m^{2}\)（即 \(y^{2}=x^{2}-m^{2}\)）化简： \(x^{2}+2m^{2}+y^{2}=2x^{2}+m^{2}\)，\((2x^{2}+m^{2})^{2}-8m^{2}x^{2}=4x^{4}+4m^{2}x^{2}+m^{4}-8m^{2}x^{2}=4x^{4}-4m^{2}x^{2}+m^{4}=(2x^{2}-m^{2})^{2}\)。 故 \(|PF_{1}|\cdot |PF_{2}|=|2x^{2}-m^{2}|=2x^{2}-m^{2}\)（\(P\) 在双曲线上 \(|x|\geqslant m\)，\(2x^{2}-m^{2}\geqslant m^{2}>0\)）\(=2x^{2}-(x^{2}-y^{2})=x^{2}+y^{2}=|PO|^{2}\)。证毕。 验算：\(m=1\)，\(P(\sqrt{2},1)\)：\(|PO|^{2}=3\)；\(|PF_{1}|=\sqrt{(\sqrt\{2\}+\sqrt\{2\})^{2}+1}=\sqrt{9}=3\)，\(|PF_{2}|=\sqrt{(\sqrt\{2\}-\sqrt\{2\})^{2}+1}=1\)，积\(=3\) 。 （卷③自带详解照录校订：源详解根号内配方路径等价，此处以 \(y^{2}=x^{2}-m^{2}\) 直接配方；关键步 \(2m^{2}=x^{2}-y^{2}\) 使 \(x^{2}+2m^{2}+y^{2}=3x^{2}-y^{2}\)，积\(=(x^{2}+y^{2})^{2}\) 与源一致。）}
\end{ansblock}

\zux{拓展·中心对称弦矩形面积}
\tuhao{09} \tieside{难}已知 \(F_{1}\)，\(F_{2}\) 为双曲线 \(C\)：\(x^{2}/16-y^{2}/9=1\) 的两个焦点，\(P\)，\(Q\) 为 \(C\) 上关于坐标原点对称的两点，且 \(|PQ|=|F_{1}F_{2}|\)，则四边形 \(PF_{1}QF_{2}\) 的面积为\kongda{\(18\)}。
\begin{ansblock}[2章-拓-课时14-09]
% ans:2章-拓-课时14-09
\ansitem{09}{\(18\)}
\ansnote{详解}{\(PQ\) 与 \(F_{1}F_{2}\) 互相平分（\(P\)、\(Q\) 关于 \(O\) 对称，\(F_{1}\)、\(F_{2}\) 关于 \(O\) 对称）且 \(|PQ|=|F_{1}F_{2}|\)，对角线相等又互相平分 \(\rightarrow \) 四边形 \(PF_{1}QF_{2}\) 为矩形，其面积\(=|PF_{1}|\cdot |PF_{2}|\)。 矩形中 \(|PF_{1}|^{2}+|PF_{2}|^{2}=|F_{1}F_{2}|^{2}=4c^{2}=100\)；又 \(||PF_{1}|-|PF_{2}||=2a=8\)。 \(2|PF_{1}||PF_{2}|=100-64=36 \rightarrow  |PF_{1}||PF_{2}|=18\)，面积\(18\)。 验算：\(a=4\)，\(b=3\)，\(c=5\)，\(4c^{2}=100\) ；\(2|PF_{1}||PF_{2}|=4c^{2}-(2a)^{2}=100-64=36 \rightarrow  18\) 。 （卷③自带详解照录，验算一致。）}
\end{ansblock}

\zux{拓展·OP·FP 最小值}
\tuhao{10} \tieside{难}若坐标原点 \(O\) 和点 \(F(-2,0)\) 分别为双曲线 \(x^{2}/a^{2}-y^{2}=1\)（\(a>0\)）的中心和左焦点，点 \(P\) 为双曲线右支上的任意一点，则 \(OP\cdot FP\) 的最小值为\kongda{\(3+2\sqrt{3}\)}。
\begin{ansblock}[2章-拓-课时14-10]
% ans:2章-拓-课时14-10
\ansitem{10}{\(3+2\sqrt{3}\)}
\ansnote{详解}{左焦点 \((-\sqrt{a^{2}+1},0)=(-2,0) \rightarrow  a^{2}=3\)。 设 \(P(x_{0},y_{0})\)（右支 \(x_{0}\geqslant \sqrt{3}\)），\(y_{0}^{2}=x_{0}^{2}/3-1\)。 \(OP\cdot FP=(x_{0},y_{0})\cdot (x_{0}+2,y_{0})=x_{0}^{2}+2x_{0}+y_{0}^{2}=(4/3)x_{0}^{2}+2x_{0}-1\)。 在 \([\sqrt{3},+\infty )\) 上导数 \((8/3)x_{0}+2>0\)，单调递增，\(x_{0}=\sqrt{3}\) 时取最小： \((4/3)\times 3+2\sqrt{3}-1=3+2\sqrt{3}\)。 验算：\(4+2\sqrt{3}-1=3+2\sqrt{3}\) ；对称轴 \(x_{0}=-3/4\) 在区间外 。 （卷③自带详解照录，验算一致。）}
\end{ansblock}

\zux{拓展·野生动物监测定位}
% 图债注：「如图」指源件配图，印面无图（冻片【注】裁定：去装饰图/条件全在文字/尺寸随注——图债登记汇编报告）
\tuhao{11} \tieside{难}如图，某野生保护区监测中心设置在点 \(O\) 处，正西、正东、正北处有\(3\)个监测点 \(A\)，\(B\)，\(C\)，且 \(|OA|=|OB|=|OC|=30km\)，一名野生动物观察员在保护区遇险，发出求救信号，\(3\)个监测点均收到求救信号，\(A\) 点接收到信号的时间比 \(B\) 点接收到信号的时间早 \((40/V_{0})s\)（注：信号每秒传播 \(V_{0}km\)）
\((1)\) 求观察员所有可能出现的位置的轨迹方程；
\((2)\) 若 \(C\) 点信号失灵，现立即以 \(C\) 为圆心进行\("\)圆形\("\)红外扫描，为保证有救援希望，扫描半径 \(r\) 至少是多少千米？
\liubai[24mm]{解答书写区}
\begin{ansblock}[2章-拓-课时14-11]
% ans:2章-拓-课时14-11
\ansitem{11}{\((1) x^{2}/400-y^{2}/500=1\)（\(x<0\)）；\((2) 20\sqrt{2} km\)。}
\ansnote{详解}{\((1)\) 以 \(O\) 为原点、\(AB\) 所在直线为 \(x\) 轴建系：\(A(-30,0)\)，\(B(30,0)\)，\(C(0,30)\)。 设观察员位置 \(P(x,y)\)：\(|PB|-|PA|=(40/V_{0})\cdot V_{0}=40<|AB|=60\)，轨迹为双曲线左支。 \(2a=40\)，\(2c=60\)，\(a=20\)，\(c=30\)，\(b^{2}=c^{2}-a^{2}=500\)。 轨迹方程 \(x^{2}/400-y^{2}/500=1\)（\(x<0\)）。 \((2)\) 设轨迹上一点 \(M(x,y)\)：\(|MC|^{2}=x^{2}+(y-30)^{2}=x^{2}+y^{2}-60y+900\)。 由 \(x^{2}=(4/5)y^{2}+400\)：\(|MC|^{2}=(9/5)y^{2}-60y+1300=(9/5)(y-50/3)^{2}+800\geqslant 800\)。 当 \(y=50/3\) 时 \(|MC|min=20\sqrt{2}\)，扫描半径 \(r\) 至少 \(20\sqrt{2} km\)。 验算：\(y=50/3\) 时 \((9/5)y^{2}-60y+1300=(9/5)(2500/9)-1000+1300=500-1000+1300=800\) 。 （卷③自带详解照录，验算一致。）}
\end{ansblock}

\zux{拓展·共轭双曲线新定义}
\tuhao{12}\duoxuan\hspace{0.6em} \tieside{难}定义：以双曲线的实轴为虚轴，虚轴为实轴的双曲线与原双曲线互为共轭双曲线．以下关于共轭双曲线的结论正确的是（　　　）
\lxopt{\(A\)．与 \(x^{2}/a^{2}-y^{2}/b^{2}=1\)（\(a>0\)，\(b>0\)）共轭的双曲线是 \(y^{2}/a^{2}-x^{2}/b^{2}=1\)（\(a>0\)，\(b>0\)）；}
\lxopt{\(B\)．互为共轭的双曲线渐近线不相同；}
\lxopt{\(C\)．互为共轭的双曲线的离心率为 \(e_{1}\)、\(e_{2}\)，则 \(e_{1}e_{2}\geqslant 2\)；}
\lxopt{\(D\)．互为共轭的双曲线的\(4\)个焦点在同一圆上}
\begin{ansblock}[2章-拓-课时14-12]
% ans:2章-拓-课时14-12
\ansitem{12}{\(CD\)}
\ansnote{详解}{\(A\)：按定义实虚轴互换，共轭双曲线应为 \(y^{2}/b^{2}-x^{2}/a^{2}=1\)，\(A\) 错； \(B\)：原曲线渐近线 \(y=\pm (b/a)x\)，共轭曲线渐近线 \(y=\pm (b/a)x\)（\(a/b\) 与 \(b/a\) 互换后方程相同），渐近线相同，\(B\) 错； \(C\)：\(e_{1}=c/a\)，\(e_{2}=c/b\)，\(e_{1}e_{2}=c^{2}/(ab)=(a^{2}+b^{2})/(ab)=a/b+b/a\geqslant 2\)（当且仅当 \(a=b\) 取等），\(C\) 对； \(D\)：焦点 \((\pm c,0)\) 与 \((0,\pm c)\) 均在圆 \(x^{2}+y^{2}=c^{2}\) 上，\(D\) 对。故选\(CD\)。 验算：\(C\) 取等条件 \(a=b\)（等轴双曲线 \(e_{1}=e_{2}=\sqrt{2}\)，积\(2\)）；\(D\)：\(c^{2}=a^{2}+b^{2}\) 。 （卷③自带详解照录，验算一致。）}
\end{ansblock}

\zux{拓展·基本量直读}
\tuhao{14} \tieside{难}写出双曲线 \(x^{2}-y^{2}/24=1\) 的实轴长、虚轴长、焦点坐标、渐近线方程．
\liubai[16mm]{解答书写区}
\begin{ansblock}[2章-拓-课时14-14]
% ans:2章-拓-课时14-14
\ansitem{14}{实轴长\(2\)；虚轴长\(4\sqrt{6}\)；焦点\((\pm 5,0)\)；渐近线方程 \(y=\pm 2\sqrt{6}x\)。}
\ansnote{详解}{\(a^{2}=1\)、\(b^{2}=24\)，\(a=1\)、\(b=2\sqrt{6}\)，\(c^{2}=1+24=25\)，\(c=5\)。 实轴长\(2a=2\)；虚轴长\(2b=4\sqrt{6}\)；焦点\((\pm 5,0)\)；渐近线 \(y=\pm (b/a)x=\pm 2\sqrt{6}x\)。 （本块自撰亲算：\(25=1+24\) ；教材页底答案条乱码，读数自验。）}
\end{ansblock}

\zux{拓展·焦点＋渐近线求方程与e}
\tuhao{15} \tieside{难}已知双曲线的一个焦点是 \((5,0)\)，一条渐近线方程为 \(3x-4y=0\)，求这个双曲线的标准方程和离心率．
\liubai[16mm]{解答书写区}
\begin{ansblock}[2章-拓-课时14-15]
% ans:2章-拓-课时14-15
\ansitem{15}{\(x^{2}/16-y^{2}/9=1\)；\(e=5/4\)。}
\ansnote{详解}{焦点 \((5,0)\) 在 \(x\) 轴：\(b/a=3/4\)（渐近线 \(y=\pm (b/a)x\)）。 \(c=5\)，\(a^{2}+b^{2}=25\)，\(b^{2}=(9/16)a^{2} \rightarrow  a^{2}(1+9/16)=25 \rightarrow  a^{2}=16\)，\(b^{2}=9\)。 标准方程 \(x^{2}/16-y^{2}/9=1\)，\(e=c/a=5/4\)。 （本块自撰亲算：\(16+9=25\) ；渐近线 \(y=\pm (3/4)x\) 。）}
\end{ansblock}

\zux{拓展·λ系共同点}
\tuhao{16} \tieside{难}当实数 \(\lambda \neq 0\) 时，方程 \(x^{2}/a^{2}-y^{2}/b^{2}=\lambda \) 表示的都是双曲线，这些双曲线的共同点是什么？
\liubai[16mm]{解答书写区}
\begin{ansblock}[2章-拓-课时14-16]
% ans:2章-拓-课时14-16
\ansitem{16}{它们的渐近线相同，均为 \(y=\pm (b/a)x\)。}
\ansnote{详解}{\(\lambda >0\) 时方程即 \(x^{2}/(a^{2}\lambda )-y^{2}/(b^{2}\lambda )=1\)，焦点在 \(x\) 轴，渐近线由 \(x^{2}/(a^{2}\lambda )-y^{2}/(b^{2}\lambda )=0\) 得 \(y=\pm (b/a)x\)； \(\lambda <0\) 时方程即 \(y^{2}/(-b^{2}\lambda )-x^{2}/(-a^{2}\lambda )=1\)，焦点在 \(y\) 轴，渐近线由同式得 \(y=\pm (b/a)x\)。 故所有曲线渐近线相同：\(y=\pm (b/a)x\)。 （本块自撰亲算；渐近线只取决于 \(a:b\) 之比，与 \(\lambda \) 无关。）}
\end{ansblock}

\zux{拓展·顶点焦距四等分求方程}
\tuhao{17} \tieside{难}已知双曲线两顶点间的距离是\(6\)，两焦点的连线被两顶点和中心四等分，求双曲线的标准方程．
\liubai[16mm]{解答书写区}
\begin{ansblock}[2章-拓-课时14-17]
% ans:2章-拓-课时14-17
\ansitem{17}{\(x^{2}/9-y^{2}/27=1\)（焦点在 \(x\) 轴；焦点在 \(y\) 轴时为 \(y^{2}/9-x^{2}/27=1\)）。}
\ansnote{详解}{顶点距 \(6 \rightarrow  2a=6\)，\(a=3\)。 焦点段 \([-c,c]\) 上的分点 \(-a,0,a\) 将其四等分，等分距 \(a=c/2 \rightarrow  c=6\)。 \(b^{2}=c^{2}-a^{2}=36-9=27\)。 焦点在 \(x\) 轴：\(x^{2}/9-y^{2}/27=1\)。 （本块自撰亲算：题面未指明焦点位置，教材按 \(x\) 轴书写；\(y\) 轴情形同法。验算 \(3\)、\(6\) 间距均为\(3\) 。）}
\end{ansblock}

\zux{拓展·同心圆网格作图探究}
% 图债注：「如图」指源件配图，印面无图（冻片【注】裁定：去装饰图/条件全在文字/尺寸随注——图债登记汇编报告）
\tuhao{18} \tieside{难}如图，\(A\)，\(B\) 是平面上的两点，且 \(|AB|=10\)，图中的一系列圆是圆心分别为 \(A\)，\(B\) 的两组同心圆，每组同心圆的半径分别是 \(1\)，\(2\)，\(3\)，…．在两组同心圆的交点中，找出与 \(A\)，\(B\) 两点的距离之差的绝对值等于 \(6\) 的点，并把这些点用光滑的曲线顺次连接起来，观察所得曲线的形状．
\liubai[16mm]{解答书写区}
\begin{ansblock}[2章-拓-课时14-18]
% ans:2章-拓-课时14-18
\ansitem{18}{曲线是双曲线（两支）：以 \(AB\) 所在直线为 \(x\) 轴、\(AB\) 中点为原点，其方程为 \(x^{2}/9-y^{2}/16=1\)。}
\ansnote{详解}{曲线上点 \(P\) 满足 \(||PA|-|PB||=6<|AB|=10\)，由定义轨迹为双曲线，\(2a=6\)、\(2c=10\)，\(a=3\)、\(c=5\)，\(b^{2}=25-9=16\)。 以 \(AB\) 中点为原点、\(AB\) 为 \(x\) 轴建系，方程 \(x^{2}/9-y^{2}/16=1\)。 作图观察：交点连成两支光滑曲线，与双曲线吻合。 （本块自撰亲算；作图探究题，图形观察与定义互证。）}
\end{ansblock}

\zux{拓展·渐近线＋椭圆端点求方程}
\tuhao{19} \tieside{难}已知双曲线的渐近线方程为 \(3x\pm 4y=0\)，它的焦点是椭圆 \(x^{2}/10+y^{2}/5=1\) 的长轴的端点，求此双曲线的标准方程．
\liubai[16mm]{解答书写区}
\begin{ansblock}[2章-拓-课时14-19]
% ans:2章-拓-课时14-19
\ansitem{19}{\(x^{2}/(32/5)-y^{2}/(18/5)=1\)。}
\ansnote{详解}{椭圆 \(a^{2}=10\)、\(b^{2}=5\)，\(c^{2}=5\)，长轴端点 \((\pm \sqrt{10},0)\)，即双曲线焦点在 \(x\) 轴，\(c=\sqrt{10}\)。 渐近线 \(y=\pm (3/4)x \rightarrow  b/a=3/4\)。 \(a^{2}+b^{2}=c^{2}=10\)，\(a^{2}(1+9/16)=10 \rightarrow  a^{2}\times 25/16=10 \rightarrow  a^{2}=32/5\)，\(b^{2}=10-32/5=18/5\)。 标准方程 \(x^{2}/(32/5)-y^{2}/(18/5)=1\)。 （本块自撰亲算：\(32/5+18/5=10\) ；\(b/a=\sqrt{18/5}\div \sqrt{32/5}=\sqrt{18/32}=3/4\) 。）}
\end{ansblock}

\zux{拓展·焦点＋渐近线求方程}
\tuhao{20} \tieside{难}已知中心在原点的双曲线的一个焦点是 \(F_{1}(-4,0)\)，一条渐近线方程是 \(3x-2y=0\)，求此双曲线的标准方程．
\liubai[16mm]{解答书写区}
\begin{ansblock}[2章-拓-课时14-20]
% ans:2章-拓-课时14-20
\ansitem{20}{\(x^{2}/(64/13)-y^{2}/(144/13)=1\)。}
\ansnote{详解}{焦点在 \(x\) 轴，\(c=4\)；渐近线 \(y=(3/2)x \rightarrow  b/a=3/2\)。 \(a^{2}+b^{2}=16\)，\(b^{2}=(9/4)a^{2} \rightarrow  a^{2}(13/4)=16 \rightarrow  a^{2}=64/13\)，\(b^{2}=16-64/13=144/13\)。 标准方程 \(x^{2}/(64/13)-y^{2}/(144/13)=1\)。 （本块自撰亲算：\(64/13+144/13=208/13=16\) ；\((144/64)=9/4\) 。）}
\end{ansblock}

\zux{拓展·半实轴＋共焦点求方程}
\tuhao{21} \tieside{难}求半实轴长为 \(2\sqrt{3}\)，且与双曲线 \(x^{2}/16-y^{2}/4=1\) 有公共焦点的双曲线的标准方程．
\liubai[16mm]{解答书写区}
\begin{ansblock}[2章-拓-课时14-21]
% ans:2章-拓-课时14-21
\ansitem{21}{\(x^{2}/12-y^{2}/8=1\)。}
\ansnote{详解}{原双曲线 \(c^{2}=16+4=20\)，公共焦点 \(\rightarrow \) 所求双曲线 \(c^{2}=20\)，焦点在 \(x\) 轴。 半实轴长 \(a=2\sqrt{3}\)，\(a^{2}=12\)；\(b^{2}=c^{2}-a^{2}=20-12=8\)。 标准方程 \(x^{2}/12-y^{2}/8=1\)。 （本块自撰亲算：\(12+8=20\) 。）}
\end{ansblock}

\zux{拓展·双曲线上点到定点距离最小值}
\tuhao{22} \tieside{难}已知 \(P\) 是双曲线 \(x^{2}-y^{2}/3=1\) 上一点，\(A(3,0)\)，求 \(|PA|\) 的最小值．
\liubai[16mm]{解答书写区}
\begin{ansblock}[2章-拓-课时14-22]
% ans:2章-拓-课时14-22
\ansitem{22}{\(2\)}
\ansnote{详解}{设 \(P(x,y)\)（\(|x|\geqslant 1\)）：\(|PA|^{2}=(x-3)^{2}+y^{2}=(x-3)^{2}+3(x^{2}-1)=4x^{2}-6x+6\)。 对称轴 \(x=3/4<1\)，故 \(x=1\)（右顶点）时取最小：\(4-6+6=4\)，\(|PA|min=2\)。 （本块自撰亲算：\(y^{2}=3(x^{2}-1)\) 代入消元 ；对称轴在 \(|x|\geqslant 1\) 左侧，端点取值。）}
\end{ansblock}

\zux{拓展·距离比定值求轨迹（第二定义引题）}
\tuhao{23} \tieside{难}设动点 \(M\) 到定点 \(F(3,0)\) 的距离与它到直线 \(l\)：\(x=4/3\) 的距离之比为 \(3/2\)，求点 \(M\) 的轨迹方程．
\liubai[16mm]{解答书写区}
\begin{ansblock}[2章-拓-课时14-23]
% ans:2章-拓-课时14-23
\ansitem{23}{\(x^{2}/4-y^{2}/5=1\)。}
\ansnote{详解}{（直接法，不引第二定义）设 \(M(x,y)\)： \(\sqrt{(x-3)^{2}+y^{2}}=(3/2)|x-4/3|\)。 两边平方：\(x^{2}-6x+9+y^{2}=(9/4)(x^{2}-(8/3)x+16/9)=(9/4)x^{2}-6x+4\)。 \(-6x\) 抵消：\(y^{2}=(9/4)x^{2}+4-x^{2}-9=(5/4)x^{2}-5\)，即 \((5/4)x^{2}-y^{2}=5\)， 标准形 \(x^{2}/4-y^{2}/5=1\)。 验算：\(F(3,0)\) 与准线位置 \(x=a^{2}/c=4/3\)、比值 \(c/a=3/2\) 相容（\(a=2\)，\(c=3\) ）。}
\end{ansblock}

\zux{拓展·焦点与准线距离之比（第二定义引题）}
\tuhao{24} \tieside{难}已知双曲线 \(x^{2}/a^{2}-y^{2}/b^{2}=1\)（\(a>0\)，\(b>0\)）的左焦点为 \(F(-c,0)\)，直线 \(l\)：\(x=-a^{2}/c\)．设 \(P\) 是双曲线上的一点，求 \(P\) 到 \(F\) 的距离与 \(P\) 到直线 \(l\) 的距离之比．
\liubai[16mm]{解答书写区}
\begin{ansblock}[2章-拓-课时14-24]
% ans:2章-拓-课时14-24
\ansitem{24}{比值等于 \(c/a\)（离心率 \(e\)）。}
\ansnote{详解}{（直接法）\(P(x,y)\) 在双曲线上：\(y^{2}=b^{2}(x^{2}/a^{2}-1)\)。 \(d_{1}=|PF|=\sqrt{(x+c)^{2}+y^{2}}=\sqrt{x^{2}+2cx+c^{2}+b^{2}x^{2}/a^{2}-b^{2}}\)，以 \(b^{2}=c^{2}-a^{2}\) 代入： \(=x^{2}+2cx+c^{2}+(c^{2}-a^{2})(x^{2}/a^{2})-(c^{2}-a^{2})=(c^{2}/a^{2})x^{2}+2cx+a^{2}=[(c/a)x+a]^{2}\)。 （\(x\geqslant a\) 时 \((c/a)x+a>0\)；\(x\leqslant -a\) 时亦有 \((c/a)x+a<0\)，平方去根号后 \(|PF|=(c/a)|x+a^{2}/c|\)。） \(d_{2}=|x+a^{2}/c|\)（\(P\) 到 \(l\)：\(x=-a^{2}/c\) 的距离）。 故 \(d_{1}/d_{2}=(c/a)|x+a^{2}/c|/|x+a^{2}/c|=c/a\)。 （本块自撰亲算：配方 \([(c/a)x+a]^{2}=c^{2}x^{2}/a^{2}+2cx+a^{2}\) 展开核对 。）}
\end{ansblock}

\zux{拓展·焦点三角形面积}
\tuhao{25} \tieside{难}已知 \(F_{1}\)，\(F_{2}\) 是双曲线 \(x^{2}/9-y^{2}/16=1\) 的两个焦点，点 \(M\) 在双曲线上，如果 \(MF_{1}\perp MF_{2}\)，求 \(\triangle MF_{1}F_{2}\) 的面积．
\liubai[16mm]{解答书写区}
\begin{ansblock}[2章-拓-课时14-25]
% ans:2章-拓-课时14-25
\ansitem{25}{\(16\)}
\ansnote{详解}{\(a=3\)，\(b=4\)，\(c=5\)。设两焦半径 \(r_{1}=|MF_{1}|\)，\(r_{2}=|MF_{2}|\)（\(r_{1}>r_{2}\)）。 \(MF_{1}\perp MF_{2}\)：\(r_{1}^{2}+r_{2}^{2}=|F_{1}F_{2}|^{2}=100\)；定义：\(r_{1}-r_{2}=2a=6\)。 \((r_{1}-r_{2})^{2}=100-2r_{1}r_{2}=36 \rightarrow  r_{1}r_{2}=32\)。 \(S=(1/2)r_{1}r_{2}\cdot sin90^{\circ}=16\)。 （本块自撰亲算：勾股 \(100-36=64=2r_{1}r_{2} \rightarrow  r_{1}r_{2}=32\) 。）}
\end{ansblock}

\zux{拓展·斜率积轨迹与分类}
\tuhao{26} \tieside{难}如果过点 \((6,0)\) 的直线与过点 \((-6,0)\) 的直线相交于点 \(M\)，而且两直线斜率的乘积为 \(a\)，其中 \(a\neq 0\)．
\((1)\) 求点 \(M\) 的轨迹方程；
\((2)\) 讨论 \(M\) 的轨迹是何种曲线．
\liubai[24mm]{解答书写区}
\begin{ansblock}[2章-拓-课时14-26]
% ans:2章-拓-课时14-26
\ansitem{26}{\((1) ax^{2}-y^{2}=36a\)（\(y\neq 0\)）；\((2) a>0\) 时为焦点在 \(x\) 轴的双曲线（除去顶点）；\(a<0\) 时为椭圆：其中 \(a=-1\) 为圆 \(x^{2}+y^{2}=36\)（\(y\neq 0\)），\(-1<a<0\) 为焦点在 \(x\) 轴的椭圆，\(a<-1\) 为焦点在 \(y\) 轴的椭圆（均除去长轴端点——方程中 \(y\neq 0\)）。}
\ansnote{详解}{\((1)\) 过 \((6,0)\) 斜率 \(k_{1}\)、过 \((-6,0)\) 斜率 \(k_{2}\)，\(k_{1}k_{2}=a\)。 \(M(x,y)\)：\(k_{1}=y/(x-6)\)，\(k_{2}=y/(x+6)\)（\(x\neq \pm 6\)，且 \(M\) 不与 \((\pm 6,0)\) 重合故 \(y\neq 0\)）。 \([y/(x-6)]\cdot [y/(x+6)]=a \rightarrow  y^{2}=a(x^{2}-36) \rightarrow  ax^{2}-y^{2}=36a\)（\(y\neq 0\)）。 \((2) a>0\)：\(x^{2}/36-y^{2}/(36a)=1\)，焦点在 \(x\) 轴的双曲线（除去 \((\pm 6,0)\)）； \(a<0\)：记 \(a=-|a|\)：\(x^{2}/36+y^{2}/(36|a|)=1\)，椭圆； 其中 \(a=-1\)：\(x^{2}+y^{2}=36\)，圆；\(-1<a<0\)：\(36>36|a|\)，焦点在 \(x\) 轴的椭圆；\(a<-1\)：\(36<36|a|\)，焦点在 \(y\) 轴的椭圆（均除去 \(x\) 轴上的点，即 \(y\neq 0\)）。 （本块自撰亲算；分类讨论按 \(|a|\) 与 \(1\) 比较。）}
\end{ansblock}

\zux{拓展·第二定义轨迹与例3解释}
\tuhao{27} \tieside{难}设动点 \(M\) 到定点 \(F(-c,0)\) 的距离与它到直线 \(l\)：\(x=-a^{2}/c\) 的距离之比为 \(c/a\)，其中 \(c>a>0\)，求点 \(M\) 的轨迹方程，并用得到的轨迹方程解释 \(2.6.2\) 中例 \(3\) 得到的结果．
\liubai[16mm]{解答书写区}
\begin{ansblock}[2章-拓-课时14-27]
% ans:2章-拓-课时14-27
\ansitem{27}{\(x^{2}/a^{2}-y^{2}/b^{2}=1\)，其中 \(b^{2}=c^{2}-a^{2}\)。}
\ansnote{详解}{（直接法）设 \(M(x,y)\)： \(\sqrt{(x+c)^{2}+y^{2}}=(c/a)|x+a^{2}/c|\)。 平方：\(x^{2}+2cx+c^{2}+y^{2}=(c^{2}/a^{2})(x^{2}+2a^{2}x/c+a^{4}/c^{2})=(c^{2}/a^{2})x^{2}+2cx+a^{2}\)。 \(2cx\) 抵消：\(y^{2}=(c^{2}/a^{2})x^{2}-x^{2}+a^{2}-c^{2}=(c^{2}/a^{2}-1)x^{2}-(c^{2}-a^{2})=(b^{2}/a^{2})x^{2}-b^{2}\)（\(b^{2}=c^{2}-a^{2}\)）。 即 \((b^{2}/a^{2})x^{2}-y^{2}=b^{2} \rightarrow  x^{2}/a^{2}-y^{2}/b^{2}=1\)。 解释：\(2.6.2\) 例 \(3\) 中所得轨迹方程正是 \(x^{2}/a^{2}-y^{2}/b^{2}=1\)（\(b^{2}=c^{2}-a^{2}\)），本题表明该双曲线可由\("\)到定点与到定直线距离之比为常数 \(c/a"\)直接刻画。 （本块自撰亲算。）}
\end{ansblock}

\zux{拓展·双垂线恒过定点}
\tuhao{28} \tieside{难}双曲线 \(C\)：\(x^{2}-y^{2}/2=1\)，过定点 \(A(-1,0)\) 的两条互相垂直的直线分别交双曲线于 \(P\)、\(Q\) 两点，直线 \(PQ\) 恒过定点（　　）
\bindopt {\fontsize{10.5pt}{\qpTJgLeadLiti}\selectfont \makebox[\dimexpr0.25\linewidth-0.25\lxhang-0.25em\relax][l]{\(A\)．\((1,0)\)}\makebox[\dimexpr0.25\linewidth-0.25\lxhang-0.25em\relax][l]{\(B\)．\((2,0)\)}\makebox[\dimexpr0.25\linewidth-0.25\lxhang-0.25em\relax][l]{\(C\)．\((3,0)\)}\makebox[\dimexpr0.25\linewidth-0.25\lxhang-0.25em\relax][l]{\(D\)．\((4,0)\)}\par}
\begin{ansblock}[2章-拓-课时14-28]
% ans:2章-拓-课时14-28
\ansitem{28}{\(C\)}
\ansnote{详解}{设 \(PQ\)：\(x=my+b\)，联立 \(x^{2}-y^{2}/2=1\)：\((m^{2}-1/2)y^{2}+2bmy+b^{2}-1=0\)。 设 \(P(x_{1},y_{1})\)，\(Q(x_{2},y_{2})\)：\(y_{1}+y_{2}=-2bm/(m^{2}-1/2)\)，\(y_{1}y_{2}=(b^{2}-1)/(m^{2}-1/2)\)。 \(AP\perp AQ\)（\(A(-1,0)\)）：\((x_{1}+1)(x_{2}+1)+y_{1}y_{2}=0\)，以 \(x_{i}=my_{i}+b\) 代入： \((1+m^{2})y_{1}y_{2}+(b+1)m(y_{1}+y_{2})+(b+1)^{2}=0\)。 通分（乘 \(m^{2}-1/2\)）： \(2(b^{2}-1)(1+m^{2})-4bm^{2}(b+1)+(2m^{2}-1)(b+1)^{2}=0\)。 按 \(m\) 整理——\(m^{2}\) 系数：\(2(b^{2}-1)-4b(b+1)+2(b+1)^{2}=2b^{2}-2-4b^{2}-4b+2b^{2}+4b+2=0\)（恒为\(0\)）； 常数项：\(2(b^{2}-1)-(b+1)^{2}=b^{2}-2b-3=0 \rightarrow  b=3\) 或 \(b=-1\)。 \(b=-1\) 时 \(PQ\) 过 \(A(-1,0)\)（\(P\)、\(Q\) 之一与 \(A\) 重合或直线退化），舍去；\(b=3\)，\(PQ\) 恒过 \((3,0)\)。故选\(C\)。 验算（特例）：\(PQ\perp x\) 轴即 \(m=0\)、\(x=3\)：与 \(x^{2}-y^{2}/2=1\) 交于 \((3,\pm 4)\)；\(AP=(4,4)\)、\(AQ=(4,-4)\)，\(AP\cdot AQ=16-16=0\) 。}
\end{ansblock}

\zux{拓展·e=2＋倍角定点存在}
\tuhao{29} \tieside{难}已知双曲线 \(C\)：\(x^{2}/a^{2}-y^{2}/b^{2}=1\)（\(a>0\)，\(b>0\)）的右焦点为 \(F(c,0)\)，离心率为 \(2\)，直线 \(x=a^{2}/c\) 与双曲线 \(C\) 的一条渐近线交于点 \(P\)，且 \(|PF|=\sqrt{3}\)．
\((1)\) 求双曲线 \(C\) 的标准方程；
\((2)\) 设 \(Q\) 为双曲线右支上的一个动点，证明：在 \(x\) 轴的负半轴上存在定点 \(M\)，使得 \(\angle QFM=2\angle QMF\)．
\liubai[24mm]{解答书写区}
\begin{ansblock}[2章-拓-课时14-29]
% ans:2章-拓-课时14-29
\ansitem{29}{\((1) x^{2}-y^{2}/3=1\)；\((2)\) 证明见解析，\(M(-1,0)\)。}
\ansnote{详解}{\((1)\) 取渐近线 \(y=(b/a)x\)，与 \(x=a^{2}/c\) 交于 \(P(a^{2}/c, ab/c)\)。 \(|PF|^{2}=(c-a^{2}/c)^{2}+(ab/c)^{2}\)，其中 \(c-a^{2}/c=(c^{2}-a^{2})/c=b^{2}/c\)，故 \(|PF|^{2}=b^{4}/c^{2}+a^{2}b^{2}/c^{2}=b^{2}(a^{2}+b^{2})/c^{2}=b^{2}c^{2}/c^{2}=b^{2}\)。 由 \(|PF|=\sqrt{3}\) 得 \(b^{2}=3\)。 \(e^{2}=c^{2}/a^{2}=4=(a^{2}+b^{2})/a^{2} \rightarrow  a^{2}=1\)。 \(C\)：\(x^{2}-y^{2}/3=1\)。 \((2) F(2,0)\)。设 \(Q(x_{0},y_{0})\)（\(x_{0}\geqslant 1\)），\(3x_{0}^{2}-y_{0}^{2}=3\)。 ① \(x_{0}=2\) 时 \(y_{0}=\pm 3\)：\(\angle QFM=90^{\circ}\) 需 \(\angle QMF=45^{\circ}\)，\(|MF|=|QF|=|y_{0}|=3 \rightarrow  M(-1,0)\)，符合。 ② \(x_{0}\neq 2\) 时设 \(M(t,0)\)（\(t<0\)）。有向正切：\(tan\angle QFM=-y_{0}/(x_{0}-2)\)，\(tan\angle QMF=y_{0}/(x_{0}-t)\)。 \(\angle QFM=2\angle QMF \Rightarrow  -y_{0}/(x_{0}-2)=2[y_{0}/(x_{0}-t)]/(1-[y_{0}/(x_{0}-t)]^{2})\)（倍角公式）。 \(y_{0}\neq 0\) 时约去 \(y_{0}\)、交叉相乘：\(-[(x_{0}-t)^{2}-y_{0}^{2}]=2(x_{0}-2)(x_{0}-t)\)， 展开整理：\(3x_{0}^{2}-y_{0}^{2}-(4+4t)x_{0}+4t+t^{2}=0\)。 以 \(3x_{0}^{2}-y_{0}^{2}=3\) 代入：\(-(4+4t)x_{0}+(t^{2}+4t+3)=0\)，对任意 \(x_{0}\geqslant 1\)（\(x_{0}\neq 2\)）恒成立， 故 \(-(4+4t)=0\) 且 \(t^{2}+4t+3=0 \rightarrow  t=-1\)（两式公共解）。 \(y_{0}=0\)（\(Q\) 为右顶点）时直接验证 \(M(-1,0)\) 亦满足 \(\angle QFM=2\angle QMF\)。 综上，存在定点 \(M(-1,0)\)。 验算：\(t^{2}+4t+3=(t+1)(t+3)\) 与 \(-(4+4t)\) 公共根 \(t=-1\) ；取 \(Q(2,3)\)：\(\angle QFM=90^{\circ}\)、\(\angle QMF=45^{\circ}\)，\(2\times 45^{\circ}=90^{\circ}\) 。}
\end{ansblock}

\zux{拓展·定义轨迹＋斜率范围＋定点}
\tuhao{30} \tieside{难}已知 \(F_{1}(-2,0)\)，\(F_{2}(2,0)\)，点 \(P\) 满足 \(|PF_{1}|-|PF_{2}|=2\)，记点 \(P\) 的轨迹为 \(\Gamma \)．斜率为 \(k\) 的直线 \(l\) 过点 \(F_{2}\)，且与轨迹 \(\Gamma \) 相交于 \(A\)，\(B\) 两点．
\((1)\) 求轨迹 \(\Gamma \) 的方程；
\((2)\) 求斜率 \(k\) 的取值范围；
\((3)\) 在 \(x\) 轴上是否存在定点 \(M\)，使得无论直线 \(l\) 绕点 \(F_{2}\) 怎样转动，总有 \(MA\perp MB\) 成立？如果存在，求出定点 \(M\)；如果不存在，请说明理由．
\liubai[32mm]{解答书写区}
\begin{ansblock}[2章-拓-课时14-30]
% ans:2章-拓-课时14-30
\ansitem{30}{\((1) x^{2}-y^{2}/3=1\)（\(x>0\)）；\((2) (-\infty ,-\sqrt{3})\cup (\sqrt{3},+\infty )\)；\((3)\) 存在，\(M(-1,0)\)。}
\ansnote{详解}{\((1) |PF_{1}|-|PF_{2}|=2<|F_{1}F_{2}|=4\)，轨迹为以 \(F_{1}\)、\(F_{2}\) 为焦点、实轴长 \(2\) 的双曲线右支：\(c=2\)，\(a=1\)，\(b^{2}=3\)，\(\Gamma \)：\(x^{2}-y^{2}/3=1\)（\(x>0\)）。 \((2) l\)：\(y=k(x-2)\)，代入 \(\Gamma \)：\((3-k^{2})x^{2}+4k^{2}x-4k^{2}-3=0\)。 \(A\)、\(B\) 为两交点（右支上 \(x>0\)）：需 \(3-k^{2}\neq 0\)、\(\Delta >0\)、\(x_{1}+x_{2}=-4k^{2}/(3-k^{2})>0\)、\(x_{1}x_{2}=(-4k^{2}-3)/(3-k^{2})>0\)。 由 \(x_{1}+x_{2}>0\) 得 \(k^{2}>3\)；此时 \(x_{1}x_{2}>0\)、\(\Delta =16k^{4}-4(3-k^{2})(-4k^{2}-3)>0\) 亦成立。 故 \(k\in (-\infty ,-\sqrt{3})\cup (\sqrt{3},+\infty )\)。 \((3)\) 存在性：设 \(M(m,0)\)。\(l\) 与 \(\Gamma \) 联立同上，\((3-k^{2})x^{2}+4k^{2}x-4k^{2}-3=0\)（\(k^{2}>3\)）， \(x_{1}+x_{2}=-4k^{2}/(3-k^{2})\)，\(x_{1}x_{2}=(-4k^{2}-3)/(3-k^{2})\)。 \(MA\perp MB \Leftrightarrow  (x_{1}-m)(x_{2}-m)+y_{1}y_{2}=0\)，\(y_{i}=k(x_{i}-2)\)： \((1+k^{2})x_{1}x_{2}-(m+2k^{2})(x_{1}+x_{2})+m^{2}+4k^{2}=0\)。 乘 \((3-k^{2})\) 展开： \((1+k^{2})(-4k^{2}-3)+4k^{2}(2k^{2}+m)+(m^{2}+4k^{2})(3-k^{2})=0 \rightarrow  k^{2}(-m^{2}+4m+5)+3m^{2}-3=0\)，对一切 \(k^{2}>3\) 恒成立。 故 \(-m^{2}+4m+5=0\) 且 \(3m^{2}-3=0 \rightarrow  m=-1\)（\(m=1\) 不满足第一式）。 斜率不存在时 \(l\)：\(x=2\)，\(A(2,3)\)、\(B(2,-3)\)：\(MA\cdot MB=(3,3)\cdot (3,-3)=0\) 。 综上，存在 \(M(-1,0)\)。 验算：\(k^{2}(-m^{2}+4m+5)+3m^{2}-3\) 在 \(m=-1\) 时两项恒为\(0\) ；\((3)\) 的韦达代入式与 \((2)\) 一致 。}
\end{ansblock}

\zux{拓展·重心定离心率}
\tuhao{32-1} \tieside{难}已知双曲线 \(x^{2}/a^{2}-y^{2}=1\)（\(a>0\)）上关于原点对称的两个点 \(P\)，\(Q\)，右顶点为 \(A\)，线段 \(AP\) 的中点为 \(E\)，直线 \(QE\) 交 \(x\) 轴于 \(M(1,0)\)，则双曲线的离心率为（　　）
\bindopt {\fontsize{10.5pt}{\qpTJgLeadLiti}\selectfont \makebox[\dimexpr0.5\linewidth-0.5\lxhang-0.25em\relax][l]{\(A\)．\(\sqrt{5}\)}\makebox[\dimexpr0.5\linewidth-0.5\lxhang-0.25em\relax][l]{\(B\)．\(\sqrt{5}/3\)}\par}
\bindopt {\fontsize{10.5pt}{\qpTJgLeadLiti}\selectfont \makebox[\dimexpr0.5\linewidth-0.5\lxhang-0.25em\relax][l]{\(C\)．\(\sqrt{10}\)}\makebox[\dimexpr0.5\linewidth-0.5\lxhang-0.25em\relax][l]{\(D\)．\(\sqrt{10}/3\)}\par}
\begin{ansblock}[2章-拓-课时14-32-1]
% ans:2章-拓-课时14-32-1
\ansitem{32-1}{\(D\)}
\ansnote{详解}{\(P(x_{0},y_{0})\)、\(Q(-x_{0},-y_{0})\)，\(\triangle APQ\) 三顶点坐标和为 \((a,0)\)，其重心为 \(G(a/3,0)\)。 重心在每条中线上，故 \(G\) 在中线 \(QE\) 上；\(QE\) 与 \(x\) 轴的交点唯一，故 \(M(1,0)\) 即重心：\(a/3=1 \rightarrow  a=3\)。 \(b^{2}=1\)，\(c=\sqrt{9+1}=\sqrt{10}\)，\(e=c/a=\sqrt{10}/3\)，故选\(D\)。 验算：\(P\)、\(Q\)、\(A\) 的纵坐标和 \(y_{0}-y_{0}+0=0\)，重心在 \(x\) 轴 ；中线 \(QE\) 过 \((a/3,0)\) 由重心坐标公式 。}
\end{ansblock}

\zux{拓展·斜率积范围}
\tuhao{32-2} \tieside{难}已知双曲线 \(C\)：\(x^{2}/a^{2}-y^{2}/b^{2}=1\)（\(a>0\)，\(b>0\)）离心率为 \(2\)，\(A\)、\(B\) 分别为左、右顶点，点 \(P\) 为双曲线 \(C\) 在第一象限内的任意一点，点 \(O\) 为坐标原点，若 \(PA\)、\(PB\)、\(PO\) 的斜率分别为 \(k_{1}\)、\(k_{2}\)、\(k_{3}\)，设 \(m=k_{1}\cdot k_{2}\cdot k_{3}\)，则 \(m\) 的取值范围为\kongda{\((0, 3\sqrt{3})\)}．
\begin{ansblock}[2章-拓-课时14-32-2]
% ans:2章-拓-课时14-32-2
\ansitem{32-2}{\((0, 3\sqrt{3})\)}
\ansnote{详解}{\(e=2 \rightarrow  c^{2}=4a^{2} \rightarrow  b^{2}=c^{2}-a^{2}=3a^{2}\)，\(b/a=\sqrt{3}\)。 \(k_{1}k_{2}=[y_{0}/(x_{0}+a)]\cdot [y_{0}/(x_{0}-a)]=y_{0}^{2}/(x_{0}^{2}-a^{2})=(b^{2}/a^{2})(x_{0}^{2}-a^{2})/(x_{0}^{2}-a^{2})=b^{2}/a^{2}=3\)。 \(P\) 在第一象限右支：\(k_{3}=k\_PO=y_{0}/x_{0}\in (0, b/a)=(0,\sqrt{3})\)（\(P\) 趋近顶点时 \(k_{3}\rightarrow 0\)，沿渐近线方向 \(k_{3}\rightarrow \sqrt{3}\)）。 \(m=3k_{3}\in (0,3\sqrt{3})\)。 （本块自撰亲算：坐标点差法同 \(14-12\)，\(k_{1}k_{2}=b^{2}/a^{2}\) 。）}
\end{ansblock}

\zux{拓展·机器人轨迹与直线包络}
\tuhao{33-1} \tieside{难}平面上一台机器人在运行中始终保持到点 \(P(-2,0)\) 的距离比到点 \(Q(2,0)\) 的距离大 \(2\)，若机器人接触不到过点 \(M(\sqrt{3},3)\) 且斜率为 \(k\) 的直线，则 \(k\) 的取值范围是\kongda{\([\sqrt{3}, 2\sqrt{3})\)}。
\begin{ansblock}[2章-拓-课时14-33-1]
% ans:2章-拓-课时14-33-1
\ansitem{33-1}{\([\sqrt{3}, 2\sqrt{3})\)}
\ansnote{详解}{机器人到 \(P\) 的距离比到 \(Q\) 大 \(2\)：距离差 \(2<|PQ|=4\)，且距 \(P\) 更远，轨迹为双曲线右支：\(a=1\)，\(c=2\)，\(b^{2}=3\)，即 \(x^{2}-y^{2}/3=1\)（\(x\geqslant 1\)）。 直线 \(y-3=k(x-\sqrt{3})\)。联立得 \((3-k^{2})x^{2}+(2\sqrt{3}k^{2}-6k)x+6\sqrt{3}k-3k^{2}-12=0\)。 直线与右支无公共点：\(k=\sqrt{3}\) 时直线恰为渐近线 \(y=\sqrt{3}x\)，符合（接触不到）；\(k=-\sqrt{3}\) 不合。 \(3-k^{2}\neq 0\) 时由 \(\Delta <0\) 解得 \(\sqrt{3}<k<2\sqrt{3}\)。 综上 \(k\in [\sqrt{3},2\sqrt{3})\)。 验算：\(k=\sqrt{3}\) 时直线 \(y-3=\sqrt{3}(x-\sqrt{3})\) 即 \(y=\sqrt{3}x\)，恰为渐近线，永不相触 ；\(k=2\sqrt{3}\) 时 \(\Delta =0\)，直线与右支相切（可触），故右端开 。}
\end{ansblock}

\zux{拓展·线段动点焦半径积最值}
\tuhao{33-2} \tieside{难}已知双曲线 \(x^{2}/a^{2}-y^{2}/b^{2}=1\)（\(a>0\)，\(b>0\)）的离心率为 \(2\)，\(F_{1}\)，\(F_{2}\) 分别是双曲线的左、右焦点，点 \(M(-a,0)\)，\(N(0,b)\)，点 \(P\) 为线段 \(MN\) 上的动点，当 \(PF_{1}\cdot PF_{2}\) 取得最大值和最小值时，\(\triangle PF_{1}F_{2}\) 的面积分别为 \(S_{1}\)，\(S_{2}\)，则 \(S_{1}/S_{2}=\)（　　）
\bindopt {\fontsize{10.5pt}{\qpTJgLeadLiti}\selectfont \makebox[\dimexpr0.25\linewidth-0.25\lxhang-0.25em\relax][l]{\(A\)．\(4\)}\makebox[\dimexpr0.25\linewidth-0.25\lxhang-0.25em\relax][l]{\(B\)．\(8\)}\makebox[\dimexpr0.25\linewidth-0.25\lxhang-0.25em\relax][l]{\(C\)．\(2\sqrt{3}\)}\makebox[\dimexpr0.25\linewidth-0.25\lxhang-0.25em\relax][l]{\(D\)．\(4\sqrt{3}\)}\par}
\begin{ansblock}[2章-拓-课时14-33-2]
% ans:2章-拓-课时14-33-2
\ansitem{33-2}{\(A\)}
\ansnote{详解}{\(e=2 \rightarrow  b^{2}/a^{2}=3\)，直线 \(MN\)：\(y=\sqrt{3}(x+a)\)。设 \(P(t,\sqrt{3}t+\sqrt{3}a)\)，\(t\in [-a,0]\)。 \(PF_{1}\cdot PF_{2}=(-c-t,-\sqrt{3}t-\sqrt{3}a)\cdot (c-t,-\sqrt{3}t-\sqrt{3}a)=t^{2}-c^{2}+3(t+a)^{2}=4t^{2}+6at-a^{2}=4(t+3a/4)^{2}-13a^{2}/4\)。 \(t=-3a/4\) 取最小（\(y\_P=\sqrt{3}a/4\)），\(t=0\) 取最大（\(y\_P=\sqrt{3}a\)）。 面积比\(=\)纵坐标绝对值比：\(S_{1}/S_{2}=\sqrt{3}a/(\sqrt{3}a/4)=4\)，故选\(A\)。 验算：\(4t^{2}+6at-a^{2}\) 的导数 \(8t+6a=0 \rightarrow  t=-3a/4\)（顶点值 \(-13a^{2}/4\)）；端点 \(t=-a\)：\(-3a^{2}=-12a^{2}/4\)、\(t=0\)：\(-a^{2}=-4a^{2}/4\)，均大于 \(-13a^{2}/4\)，故最小在顶点、最大在 \(t=0\) 。}
\end{ansblock}

\zux{拓展·焦三角形面积最大}
\tuhao{33-3} \tieside{难}设 \(O\) 为坐标原点，直线 \(x=a\) 与双曲线 \(C\)：\(x^{2}/a^{2}-y^{2}/b^{2}=1\)（\(a>0\)，\(b>0\)）的两条渐近线分别交于 \(D\)，\(E\) 两点．若 \(C\) 的焦距为 \(4\)，则 \(\triangle ODE\) 面积的最大值为\kongda{\(2\)}。
\begin{ansblock}[2章-拓-课时14-33-3]
% ans:2章-拓-课时14-33-3
\ansitem{33-3}{\(2\)}
\ansnote{详解}{\(D(a,b)\)，\(E(a,-b)\)，\(|DE|=2b\)，\(O\) 到直线 \(x=a\) 距离为 \(a\)：\(S=(1/2)\cdot a\cdot 2b=ab\)。 焦距 \(4 \rightarrow  c=2\)，\(a^{2}+b^{2}=4\geqslant 2ab \rightarrow  ab\leqslant 2\)（\(a=b=\sqrt{2}\) 取等）。 \(\triangle ODE\) 面积最大值 \(2\)。}
\end{ansblock}

\zux{拓展·焦半径最小值积＋弦比最小}
\tuhao{33-4} \tieside{难}已知双曲线 \(x^{2}/a^{2}-y^{2}/b^{2}=1\)（\(a>0\)，\(b>0\)）的左、右焦点分别为 \(F_{1}\)，\(F_{2}\)，点 \(P\) 在双曲线的右支上，\(|PF_{1}|\)，\(|PF_{2}|\) 的最小值 \(m_{1}\)，\(m_{2}\)，且满足 \(m_{1}m_{2}=3a^{2}\)．
\((1)\) 求双曲线的离心率；
\((2)\) 若 \(a=2\)，过点 \(F_{1}\) 的直线交双曲线于 \(A\)，\(B\) 两点，线段 \(AB\) 的垂直平分线交 \(y\) 轴于点 \(D\)（异于坐标原点 \(O\)），求 \(|AB|/|OD|\) 的最小值．
\liubai[24mm]{解答书写区}
\begin{ansblock}[2章-拓-课时14-33-4]
% ans:2章-拓-课时14-33-4
\ansitem{33-4}{\((1) e=2\)；\((2) 3/2\)。}
\ansnote{详解}{\((1)\) 右支上 \(|PF_{1}|min=c+a\)（左顶点处），\(|PF_{2}|min=c-a\)（右顶点处）。 \(m_{1}m_{2}=c^{2}-a^{2}=b^{2}=3a^{2} \rightarrow  c^{2}=4a^{2} \rightarrow  e=2\)。 \((2) a=2\)，\(c=4\)，\(b^{2}=12\)，双曲线 \(x^{2}/4-y^{2}/12=1\)，\(F_{1}(-4,0)\)。设 \(AB\)：\(y=k(x+4)\)（\(k\neq \pm \sqrt{3}\)）。 联立：\((3-k^{2})x^{2}-8k^{2}x-16k^{2}-12=0\)，\(x_{1}+x_{2}=8k^{2}/(3-k^{2})\)，\(x_{1}x_{2}=(-16k^{2}-12)/(3-k^{2})\)。 弦长 \(|AB|=\sqrt{1+k^{2}}\cdot \sqrt{(x_{1}+x_{2})^{2}-4x_{1}x_{2}}=12(1+k^{2})/|3-k^{2}|\)。 中点 \((4k^{2}/(3-k^{2}), 12k/(3-k^{2}))\)（\(y_{1}+y_{2}=k(x_{1}+x_{2})+8k=24k/(3-k^{2})\)）。 垂直平分线：\(y-12k/(3-k^{2})=-(1/k)(x-4k^{2}/(3-k^{2}))\)，令 \(x=0\) 得 \(y=16k/(3-k^{2})\)。 \(|OD|=16|k|/|3-k^{2}|\)。 \(|AB|/|OD|=[12(1+k^{2})/|3-k^{2}|]\cdot [|3-k^{2}|/(16|k|)]=(3/4)(|k|+1/|k|)\geqslant (3/2)\)（\(|k|=1\) 取等）。 最小值 \(3/2\)。 验算：\(|AB|\) 的核：\(\Delta \) 下根差\(^{2}=(8k^{2}/(3-k^{2}))^{2}+4(16k^{2}+12)/(3-k^{2})=[64k^{4}+4(16k^{2}+12)(3-k^{2})]/(3-k^{2})^{2}=[64k^{4}+192k^{2}+144-64k^{4}-48k^{2}]/(3-k^{2})^{2}=144(k^{2}+1)/(3-k^{2})^{2} \rightarrow  |AB|=\sqrt{1+k^{2}}\cdot 12\sqrt{k^{2}+1}/|3-k^{2}|\) 。}
\end{ansblock}

\zux{拓展·过定点双曲线＋k₁+k₂定点}
\tuhao{33-5} \tieside{难}已知双曲线的离心率为 \(\sqrt{6}/2\)，且该双曲线经过点 \(P(\sqrt{3}, \sqrt{2}/2)\)．
\((1)\) 求双曲线 \(C\)：\(x^{2}/a^{2}-y^{2}/b^{2}=1\)（\(a>0\)，\(b>0\)）的方程；
\((2)\) 设斜率分别为 \(k_{1}\)，\(k_{2}\) 的两条直线 \(l_{1}\)，\(l_{2}\) 均经过点 \(Q(2,1)\)，且直线 \(l_{1}\)，\(l_{2}\) 与双曲线 \(C\) 分别交于 \(A\)，\(B\) 两点（\(A\)，\(B\) 异于点 \(Q\)），若 \(k_{1}+k_{2}=1\)，试判断直线 \(AB\) 是否经过定点，若存在定点，求出该定点坐标；若不存在，说明理由。
\liubai[24mm]{解答书写区}
\begin{ansblock}[2章-拓-课时14-33-5]
% ans:2章-拓-课时14-33-5
\ansitem{33-5}{\((1) x^{2}/2-y^{2}=1\)；\((2)\) 过定点 \((0,1)\)。}
\ansnote{详解}{\((1) e^{2}=3/2=c^{2}/a^{2}=(a^{2}+b^{2})/a^{2} \rightarrow  b^{2}/a^{2}=1/2 \rightarrow  a^{2}=2b^{2}\)，即双曲线方程为 \(x^{2}/(2b^{2})-y^{2}/b^{2}=1\)。 代入 \(P\)：\(3/(2b^{2})-(1/2)/b^{2}=1 \rightarrow  (3-1)/(2b^{2})=1 \rightarrow  b^{2}=1\)，\(a^{2}=2\)。\(C\)：\(x^{2}/2-y^{2}=1\)。 \((2) AB\) 斜率不存在时：\(A(x_{0},y_{0})\)、\(B(x_{0},-y_{0})\)，\(k_{1}+k_{2}=(y_{0}-1)/(x_{0}-2)+(-y_{0}-1)/(x_{0}-2)=-2/(x_{0}-2)=1 \rightarrow  x_{0}=0\)，不在曲线上，舍；故斜率存在。 设 \(AB\)：\(y=kx+t\)，代入：\((2k^{2}-1)x^{2}+4ktx+2t^{2}+2=0\)（\(2k^{2}-1\neq 0\)）。 \(x_{1}+x_{2}=-4kt/(2k^{2}-1)\)，\(x_{1}x_{2}=(2t^{2}+2)/(2k^{2}-1)\)。 \(k_{1}+k_{2}=1\)：\((kx_{1}+t-1)/(x_{1}-2)+(kx_{2}+t-1)/(x_{2}-2)=1 \rightarrow  (2k-1)x_{1}x_{2}+(t-2k+1)(x_{1}+x_{2})-4t=0\)。 代入：\((2k-1)(2t^{2}+2)-4kt(t-2k+1)-4t(2k^{2}-1)=0\)（通分后乘 \(2k^{2}-1\)）， 整理 \(t^{2}+(2k-2)t-1+2k=0\)，即 \((t-1)(t+2k-1)=0\)。 \(t=1\)：\(y=kx+1\) 过 \((0,1)\)；\(t=1-2k\)：\(y=k(x-2)+1\) 过 \(Q(2,1)\)，不合（\(A\)、\(B\) 异于 \(Q\) 且 \(Q\) 为弦端情形）。 故直线 \(AB\) 过定点 \((0,1)\)。 验算：\(t^{2}+(2k-2)t+2k-1=(t-1)(t+2k-1)\) 展开核对 。}
\end{ansblock}

\zux{拓展·定点弦定值}
\tuhao{33-6} \tieside{难}已知双曲线 \(C\) 过点 \((4,\sqrt{3})\)，且渐近线方程为 \(y=\pm (1/2)x\)，直线 \(l\) 与曲线 \(C\) 交于点 \(M\)、\(N\) 两点．
\((1)\) 求双曲线 \(C\) 的方程；
\((2)\) 若直线 \(l\) 过点 \((1,0)\)，问在 \(x\) 轴上是否存在定点 \(Q\)，使得 \(QM\cdot QN\) 为常数？若存在，求出点坐标及此常数的值；若不存在，说明理由。
\liubai[24mm]{解答书写区}
\begin{ansblock}[2章-拓-课时14-33-6]
% ans:2章-拓-课时14-33-6
\ansitem{33-6}{\((1) x^{2}/4-y^{2}=1\)；\((2)\) 存在，\(Q(23/8,0)\)，\(QM\cdot QN=273/64\)。}
\ansnote{详解}{\((1) b/a=1/2\) 且 \(16/a^{2}-3/b^{2}=1\)：\(b^{2}=a^{2}/4 \rightarrow  16/a^{2}-12/a^{2}=1 \rightarrow  a^{2}=4\)，\(b^{2}=1\)。\(C\)：\(x^{2}/4-y^{2}=1\)。 \((2)\) 设 \(l\)：\(x=my+1\)，联立：\((m^{2}-4)y^{2}+2my-3=0\)（\(m^{2}\neq 4\)，\(\Delta >0 \rightarrow  m^{2}>3\)）。 \(y_{1}+y_{2}=-2m/(m^{2}-4)\)，\(y_{1}y_{2}=-3/(m^{2}-4)\)； \(x_{1}+x_{2}=m(y_{1}+y_{2})+2=-8/(m^{2}-4)\)，\(x_{1}x_{2}=(my_{1}+1)(my_{2}+1)=-(4m^{2}+4)/(m^{2}-4)\)。 设 \(Q(t,0)\)：\(QM\cdot QN=(x_{1}-t)(x_{2}-t)+y_{1}y_{2}=x_{1}x_{2}-t(x_{1}+x_{2})+t^{2}+y_{1}y_{2} \allowbreak =-4+t^{2}+(8t-23)/(m^{2}-4)\)。 与 \(m\) 无关 \(\Leftrightarrow  8t-23=0 \rightarrow  t=23/8\)，常数\(=-4+t^{2}=-4+529/64=273/64\)。 验算：\(x_{1}x_{2}=(-4m^{2}-4)/(m^{2}-4)=-4-20/(m^{2}-4)\)；三项合并分子 \(-20-3+8t=8t-23 \rightarrow  QM\cdot QN=-4+t^{2}+(8t-23)/(m^{2}-4)\) 。}
\end{ansblock}

\zux{拓展·椭圆性质类比证明}
\tuhao{33-7} \tieside{难}设直线 \(2x-y+1=0\) 与椭圆 \(x^{2}/3+y^{2}/4=1\) 相交于 \(A\)、\(B\) 两点．
\((1)\) 求弦长 \(|AB|\)；
\((2)\) 已知椭圆具有性质：设 \(A\)、\(B\) 为椭圆 \(x^{2}/a^{2}+y^{2}/b^{2}=1\) 上任意两点，\(M\) 是线段 \(AB\) 的中点，若直线 \(AB\)、\(OM\) 的斜率都存在，并记为 \(k\_AB\)、\(k\_OM\)，则 \(k\_AB\cdot k\_OM\) 为定值．试对双曲线 \(x^{2}/a^{2}-y^{2}/b^{2}=1\) 写出具有类似特征的性质，并加以证明．
\liubai[24mm]{解答书写区}
\begin{ansblock}[2章-拓-课时14-33-7]
% ans:2章-拓-课时14-33-7
\ansitem{33-7}{\((1) 15/4\)；\((2)\) 双曲线中 \(k\_AB\cdot k\_OM=b^{2}/a^{2}\)（定值），证明见解析。}
\ansnote{详解}{\((1)\) 联立得 \(16x^{2}+12x-9=0\)，\(x_{1}+x_{2}=-3/4\)，\(x_{1}x_{2}=-9/16\)。 \(|AB|=\sqrt{1+2^{2}}\cdot \sqrt{(-3/4)^{2}-4\times (-9/16)}=\sqrt{5}\cdot \sqrt{9/16+9/4}=\sqrt{5}\cdot \sqrt{45/16}=\sqrt{5}\cdot (3\sqrt{5})/4=15/4\)。 \((2)\) 性质：双曲线 \(x^{2}/a^{2}-y^{2}/b^{2}=1\) 上任意两点 \(A\)、\(B\)（斜率存在时），\(M\) 为 \(AB\) 中点，则 \(k\_AB\cdot k\_OM=b^{2}/a^{2}\)。 证明（点差法）：\(A(x_{3},y_{3})\)、\(B(x_{4},y_{4})\)、\(M(x_{0},y_{0})\)。两式相减： \((x_{3}+x_{4})(x_{3}-x_{4})/a^{2}=(y_{3}+y_{4})(y_{3}-y_{4})/b^{2} \rightarrow  2x_{0}(x_{3}-x_{4})/a^{2}=2y_{0}(y_{3}-y_{4})/b^{2}\)。 \(k\_AB\cdot k\_OM=[(y_{3}-y_{4})/(x_{3}-x_{4})]\cdot (y_{0}/x_{0})=b^{2}/a^{2}\)。定值证毕。 验算：本件 \(k\_AB\cdot k\_OM=b^{2}/a^{2}\) 与 \(14-12\) 的 \(k_{1}k_{2}=b^{2}/a^{2}\) 相容（中心对称两点为其特例）。}
\end{ansblock}

\zux{拓展·切线渐近线三角形定值}
\tuhao{33-8} \tieside{难}已知双曲线 \(C\)：\(x^{2}/a^{2}-y^{2}/b^{2}=1\)（\(a>0\)，\(b>0\)）的一个焦点坐标为 \((3,0)\)，其中一条渐近线的倾斜角的正切值为 \(2\sqrt{2}\)，\(O\) 为坐标原点．
\((1)\) 求双曲线 \(C\) 的方程；
\((2)\) 直线 \(l\) 与 \(x\) 轴正半轴相交于一点 \(D\)，与双曲线 \(C\) 右支相切（切点不为右顶点），且 \(l\) 分别交双曲线 \(C\) 的两条渐近线于 \(M\)、\(N\) 两点，证明：\(\triangle MON\) 的面积为定值，并求出该定值．
\liubai[24mm]{解答书写区}
\begin{ansblock}[2章-拓-课时14-33-8]
% ans:2章-拓-课时14-33-8
\ansitem{33-8}{\((1) x^{2}-y^{2}/8=1\)；\((2)\) 定值 \(2\sqrt{2}\)。}
\ansnote{详解}{\((1) c=3\)，\(b/a=2\sqrt{2}\)，\(a^{2}+b^{2}=9 \rightarrow  a^{2}(1+8)=9 \rightarrow  a^{2}=1\)，\(b^{2}=8\)。\(C\)：\(x^{2}-y^{2}/8=1\)。 \((2) l\)：\(y=kx+m\)（\(k\neq \pm 2\sqrt{2}\)，\(m\neq 0\)，斜率存在且不为\(0\)）。\(D(-m/k,0)\)，\(|OD|=|m/k|\)。 相切：联立 \((8-k^{2})x^{2}-2kmx-m^{2}-8=0\)，\(\Delta =0 \rightarrow  4k^{2}m^{2}+4(8-k^{2})(m^{2}+8)=0 \rightarrow  m^{2}=k^{2}-8\)（\(k^{2}>8\)）。 \(M\) 与 \(y=2\sqrt{2}x\) 交：\((m/(2\sqrt{2}-k), 2\sqrt{2}m/(2\sqrt{2}-k))\)；\(N\) 与 \(y=-2\sqrt{2}x\) 交：\((-m/(2\sqrt{2}+k), 2\sqrt{2}m/(2\sqrt{2}+k))\)。 \(S\triangle MON=S\triangle MOD+S\triangle NOD=(1/2)|OD|\cdot |y\_M-y\_N|=(1/2)|m|\cdot |x\_M-x\_N| \allowbreak =(1/2)|m|\cdot |4\sqrt{2}m/(8-k^{2})|=2\sqrt{2}m^{2}/(k^{2}-8)=2\sqrt{2}\)（定值）。 验算：\(m^{2}=k^{2}-8>0\) 与 \(k^{2}>8\) 一致 ；\(S\) 只依赖 \(m^{2}/k^{2}-8\) 化简后为常数 。}
\end{ansblock}

\zux{拓展·顶点斜率比定值}
\tuhao{33-9} \tieside{难}已知双曲线 \(\Gamma \)：\(x^{2}/a^{2}-y^{2}/b^{2}=1\)（\(a>0\)，\(b>0\)）的左、右顶点分别为 \(A_{1}(-1,0)\)、\(A_{2}(1,0)\)，离心率为 \(2\)，过点 \(F(2,0)\) 斜率不为 \(0\) 的直线 \(l\) 与 \(\Gamma \) 交于 \(P\)、\(Q\) 两点．
\((1)\) 求双曲线 \(\Gamma \) 的渐近线方程；
\((2)\) 记直线 \(A_{1}P\)、\(A_{2}Q\) 的斜率分别为 \(k_{1}\)、\(k_{2}\)，求证：\(k_{1}/k_{2}\) 为定值．
\liubai[24mm]{解答书写区}
\begin{ansblock}[2章-拓-课时14-33-9]
% ans:2章-拓-课时14-33-9
\ansitem{33-9}{\((1) y=\pm \sqrt{3}x\)；\((2)\) 定值 \(k_{1}/k_{2}=-1/3\)。}
\ansnote{详解}{\((1) a=1\)，\(c=2\)，\(b^{2}=3\)：\(\Gamma \)：\(x^{2}-y^{2}/3=1\)，渐近线 \(y=\pm \sqrt{3}x\)。 \((2) l\) 斜率不存在：\(P(2,3)\)、\(Q(2,-3)\)，\(k_{1}=3/3=1\)、\(k_{2}=-3/1=-3\)，比值 \(-1/3\)。 \(l\)：\(y=k(x-2)\)（\(k\neq 0\)）：\((k^{2}-3)x^{2}-4k^{2}x+4k^{2}+3=0\)，\(x_{1}+x_{2}=4k^{2}/(k^{2}-3)\)，\(x_{1}x_{2}=(4k^{2}+3)/(k^{2}-3)\)。 \(3k_{1}+k_{2}=3y_{1}/(x_{1}+1)+y_{2}/(x_{2}-1)=k[4x_{1}x_{2}-5(x_{1}+x_{2})+4]/[(x_{1}+1)(x_{2}-1)]\)。 \(4x_{1}x_{2}-5(x_{1}+x_{2})+4=[4(4k^{2}+3)-20k^{2}+4(k^{2}-3)]/(k^{2}-3)=0\)。 故 \(3k_{1}+k_{2}=0\)，\(k_{1}/k_{2}=-1/3\)（定值）。综上得证。 验算：分子 \(16k^{2}+12-20k^{2}+4k^{2}-12=0\) 。}
\end{ansblock}

\zux{拓展·定比分点定直线}
\tuhao{33-10} \tieside{难}设双曲线 \(x^{2}/a^{2}-y^{2}/b^{2}=1\)，其虚轴长为 \(2\sqrt{2}\)，且离心率为 \(\sqrt{5}\)．
\((1)\) 求双曲线 \(C\) 的方程；
\((2)\) 过点 \(P(3,1)\) 的动直线与双曲线的左右两支曲线分别交于点 \(A\)、\(B\)，在线段 \(AB\) 上取点 \(M\) 使得 \(|AM|/|MB|=|AP|/|PB|\)，证明：点 \(M\) 落在某一定直线上．
\liubai[24mm]{解答书写区}
\begin{ansblock}[2章-拓-课时14-33-10]
% ans:2章-拓-课时14-33-10
\ansitem{33-10}{\((1) 2x^{2}-y^{2}/2=1\)；\((2)\) 定直线 \(12x-y-2=0\)。}
\ansnote{详解}{\((1) 2b=2\sqrt{2} \rightarrow  b^{2}=2\)；\(e^{2}=5=1+b^{2}/a^{2} \rightarrow  a^{2}=1/2\)。\(C\)：\(2x^{2}-y^{2}/2=1\)。 \((2)\) 设 \(M(x,y)\)、\(A(x_{1},y_{1})\)、\(B(x_{2},y_{2})\)（\(x_{1}<x_{2}<3\)）。\(|AM|/|MB|=|AP|/|PB|\) 结合共线定比： \((3-x_{1})/(x_{2}-3)=-(x-x_{1})/(x_{2}-x)\)，即 \([6-(x_{1}+x_{2})]x=3(x_{1}+x_{2})-2x_{1}x_{2}\)。① 设 \(l\)：\(y-1=k(x-3)\)，代入 \(C\)：\((4-k^{2})x^{2}+(6k^{2}-2k)x-9k^{2}+6k-3=0\)， \(x_{1}+x_{2}=(2k-6k^{2})/(4-k^{2})\)，\(x_{1}x_{2}=(-9k^{2}+6k-3)/(4-k^{2})\)。代入①整理： \(12x-3=k(x-3)\)，与 ②（\(y-1=k(x-3)\)）联立消 \(k\)：\(12x-3=y-1\)，即 \(12x-y-2=0\)。 点 \(M\) 恒落在定直线 \(12x-y-2=0\) 上。 验算：① 的展开：由定比 \(x=(3(x_{1}+x_{2})-2x_{1}x_{2})/(6-(x_{1}+x_{2}))\)（分比不变）；韦达代入化简 \((6-(x_{1}+x_{2}))\cdot 12x=36(x_{1}+x_{2})-24x_{1}x_{2}\) 与 \(12x-3=k(x-3)\) 等价核：源详解整理结果 \(12x-3=k(x-3)\)，代回 \(y-1=k(x-3)\) 消 \(k\) 得 \(y=12x-2\) 。}
\end{ansblock}

\zux{拓展·定比弦与交点定直线}
\tuhao{33-11} \tieside{难}已知 \(A\)，\(B\) 分别是双曲线 \(E\)：\(x^{2}-y^{2}/4=1\) 的左，右顶点，直线 \(l\)（不与坐标轴垂直）过点 \(N(2,0)\)，且与双曲线 \(E\) 交于 \(C\)，\(D\) 两点．
\((1)\) 若 \(CN=3ND\)，求直线 \(l\) 的方程；
\((2)\) 若直线 \(AC\) 与 \(BD\) 相交于点 \(P\)，求证：点 \(P\) 在定直线上．
\liubai[24mm]{解答书写区}
\begin{ansblock}[2章-拓-课时14-33-11]
% ans:2章-拓-课时14-33-11
\ansitem{33-11}{\((1) 2\sqrt{5}x-y-4\sqrt{5}=0\) 或 \(2\sqrt{5}x+y-4\sqrt{5}=0\)；\((2)\) 定直线 \(x=1/2\)。}
\ansnote{详解}{设 \(l\)：\(x=my+2\)，\(C(x_{1},y_{1})\)、\(D(x_{2},y_{2})\)，联立：\((4m^{2}-1)y^{2}+16my+12=0\)， \(y_{1}+y_{2}=-16m/(4m^{2}-1)\)，\(y_{1}y_{2}=12/(4m^{2}-1)\)。 \((1) CN=3ND\)（\(N\) 为 \(F_{2}\)，\(C\)、\(D\) 同向分段）\(\rightarrow  y_{1}=-3y_{2}\)。 \(y_{2}=8m/(4m^{2}-1)\)，代入 \(y_{1}y_{2}=-3y_{2}^{2}=12/(4m^{2}-1)\)： \(-3\cdot 64m^{2}/(4m^{2}-1)^{2}=12/(4m^{2}-1) \rightarrow  -16m^{2}=4m^{2}-1 \rightarrow  m^{2}=1/20\)，\(m=\pm \sqrt{5}/10\)。 直线 \(l\)：\(x=\pm (\sqrt{5}/10)y+2\)，即 \(2\sqrt{5}x-y-4\sqrt{5}=0\) 或 \(2\sqrt{5}x+y-4\sqrt{5}=0\)。 \((2) AC\)：\(y=[y_{1}/(x_{1}+1)](x+1)\)；\(BD\)：\(y=[y_{2}/(x_{2}-1)](x-1)\)。以 \(x_{1}=my_{1}+2\)、\(x_{2}=my_{2}+2\) 代入，得 \(x_{1}+1=my_{1}+3\)、\(x_{2}-1=my_{2}+1\)。 联立：\((x+1)/(x-1)=(my_{1}y_{2}+3y_{2})/(my_{1}y_{2}+y_{1})\)。 一般情形（不求 \(y_{1}\)、\(y_{2}\) 具体值）：由韦达 \(y_{1}y_{2}=12/(4m^{2}-1)\)、\(y_{1}+y_{2}=-16m/(4m^{2}-1)\)，得 \(y_{1}y_{2}=-(3/4m)(y_{1}+y_{2})\)。 于是 \((x+1)/(x-1)=[-(3/4)(y_{1}+y_{2})+3y_{2}]/[-(3/4)(y_{1}+y_{2})+y_{1}]=(-3y_{1}+9y_{2})/(y_{1}-3y_{2})=-3\)。 解得 \(x=1/2\)。点 \(P\) 恒在定直线 \(x=1/2\) 上。 验算：\((-3y_{1}+9y_{2})/(y_{1}-3y_{2})=-3(y_{1}-3y_{2})/(y_{1}-3y_{2})=-3\) ；\((1)\) 中 \(-16m^{2}=4m^{2}-1 \rightarrow  20m^{2}=1\) 。}
\end{ansblock}

\zux{拓展·青花瓷单叶双曲面}
% 图债注：「如图」指源件配图，印面无图（冻片【注】裁定：去装饰图/条件全在文字/尺寸随注——图债登记汇编报告）
\tuhao{33-12} \tieside{难}青花瓷，中华陶瓷烧制工艺的珍品，是中国瓷器的主流品种之一．如图是一个落地青花瓷，其外形称为单叶双曲面，且它的外形左右对称，可以看成是双曲线的一部分绕其虚轴旋转所形成的曲面．若该花瓶横截面圆的最小直径为 \(16cm\)，上瓶口圆的直径为 \(20cm\)，上瓶口圆与最小圆圆心间的距离为 \(12cm\)，则该双曲线的离心率为\kongda{\(\sqrt{5}\)}。
\begin{ansblock}[2章-拓-课时14-33-12]
% ans:2章-拓-课时14-33-12
\ansitem{33-12}{\(\sqrt{5}\)}
\ansnote{详解}{以最细处直径 \(A_{1}A_{2}\) 为 \(x\) 轴、中垂线为 \(y\) 轴建系。\(2a=16\)，\(a=8\)。 上瓶口点 \(B(10,12)\) 在双曲线 \(x^{2}/64-y^{2}/b^{2}=1\) 上：\(100/64-144/b^{2}=1 \rightarrow  144/b^{2}=100/64-1=36/64 \rightarrow  b^{2}=256\)。 \(c^{2}=a^{2}+b^{2}=64+256=320\)，\(e^{2}=320/64=5\)，\(e=\sqrt{5}\)。}
\end{ansblock}

\zux{拓展·舰船信号定位}
% 图债注：「如图」指源件配图，印面无图（冻片【注】裁定：去装饰图/条件全在文字/尺寸随注——图债登记汇编报告）
\tuhao{33-13} \tieside{难}为捍卫钓鱼岛及其附属岛屿的领土主权，中国派出舰船\("\)唐山号\("\)、\("\)石家庄号\("\)和\("\)邯郸号\("\)在钓鱼岛领海巡航．某日，正巡逻在 \(A\) 处的\("\)唐山号\("\)突然发现来自 \(P\) 处的疑似敌舰的某信号，发现信号时\("\)石家庄号\("\)和\("\)邯郸号\("\)正分别位于如图所示的 \(B\)、\(C\) 两处，其中 \(A\) 在 \(B\) 的正东方向相距 \(6\) 海里处，\(C\) 在 \(B\) 的北偏西 \(30^{\circ}\) 方向相距 \(4\) 海里处．由于 \(B\)、\(C\) 比 \(A\) 距 \(P\) 更远，因此，\(4\) 秒后 \(B\)、\(C\) 才同时发现这一信号（该信号的传播速度为每秒 \(1\) 海里），试确定疑似敌舰相对于 \(A\) 点\("\)唐山号\("\)的位置．
\liubai[16mm]{解答书写区}
\begin{ansblock}[2章-拓-课时14-33-13]
% ans:2章-拓-课时14-33-13
\ansitem{33-13}{在 \(A\) 的北偏东 \(30^{\circ}\) 方向，相距 \(10\) 海里处。}
\ansnote{详解}{以 \(AB\) 为 \(x\) 轴、\(AB\) 中点为原点建系：\(A(3,0)\)，\(B(-3,0)\)，\(C(-5,2\sqrt{3})\)。 \(|PB|-|PA|=4<6=|AB|\)，\(P\) 在以 \(A\)、\(B\) 为焦点、实轴长 \(4\) 的双曲线右支上：\(x^{2}/4-y^{2}/5=1\)（\(x\geqslant 2\)）。 \(|PB|=|PC|\)（\(B\)、\(C\) 同时收到），\(P\) 在 \(BC\) 中垂线上：\(BC\) 中点 \((-4,\sqrt{3})\)，\(k\_BC=2\sqrt{3}/(-5+3)=-\sqrt{3}\)， 中垂线：\(y-\sqrt{3}=(\sqrt{3}/3)(x+4)\)，即 \(x-\sqrt{3}y+7=0\)。 与 \(5x^{2}-4y^{2}=20\) 联立（\(y=(x+7)/\sqrt{3}\) 代入）：\(11x^{2}-56x-256=0\)，解得 \(x=8\) 或 \(x=-32/11\)（左支交点，不合右支 \(x\geqslant 2\)，舍），\(y=5\sqrt{3}\)。 \(P(8,5\sqrt{3})\)：\(k\_AP=5\sqrt{3}/(8-3)=\sqrt{3}\)，倾斜角 \(60^{\circ}\)，\(|PA|=\sqrt{25+75}=10\)。 故 \(P\) 在 \(A\) 的北偏东 \(30^{\circ}\) 方向，相距 \(10\) 海里处。 验算：\(|PB|=\sqrt{(8+3)^{2}+(5\sqrt\{3\})^{2}}=\sqrt{196}=14\)，\(|PB|-|PA|=14-10=4\) ；\(|PC|^{2}=13^{2}+(3\sqrt{3})^{2}=196=|PB|^{2}\) 。}
\end{ansblock}

\zux{拓展·椭圆双曲线共焦点}
\tuhao{34-1} \tieside{难}若椭圆 \(x^{2}/9+y^{2}/k^{2}=1\) 与双曲线 \(x^{2}/k-y^{2}/3=1\) 有相同的焦点，则 \(k\) 的值为\kongda{\(2\)}。
\begin{ansblock}[2章-拓-课时14-34-1]
% ans:2章-拓-课时14-34-1
\ansitem{34-1}{\(2\)}
\ansnote{详解}{两者焦点均在 \(x\) 轴：\(9-k^{2}=k+3\)（\(k>0\)）\(\rightarrow  k^{2}+k-6=0 \rightarrow  (k+3)(k-2)=0 \rightarrow  k=2\)（\(k=-3\) 舍）。 验算：\(k=2\)：椭圆 \(c^{2}=9-4=5\)，双曲线 \(c^{2}=2+3=5\) 。}
\end{ansblock}

\zux{拓展·篮球截面双曲线}
% 图债注：「如图」指源件配图，印面无图（冻片【注】裁定：去装饰图/条件全在文字/尺寸随注——图债登记汇编报告）
\tuhao{34-2} \tieside{难}从某个角度观察篮球（如图\(1\)），可以得到一个对称的平面图形，如图\(2\)所示，篮球的外轮形状为圆 \(O\)，将篮球表面的粘合线看成坐标轴和双曲线，若坐标轴和双曲线与圆 \(O\) 的交点将圆 \(O\) 的周长八等分，\(AB=(1/2)BC=CD=1\)，则该双曲线的焦距为（　　）
\bindopt {\fontsize{10.5pt}{\qpTJgLeadLiti}\selectfont \makebox[\dimexpr0.5\linewidth-0.5\lxhang-0.25em\relax][l]{\(A\)．\(\sqrt{2}\)}\makebox[\dimexpr0.5\linewidth-0.5\lxhang-0.25em\relax][l]{\(B\)．\(\sqrt{6}/2\)}\par}
\bindopt {\fontsize{10.5pt}{\qpTJgLeadLiti}\selectfont \makebox[\dimexpr0.5\linewidth-0.5\lxhang-0.25em\relax][l]{\(C\)．\(2\sqrt{3}\)}\makebox[\dimexpr0.5\linewidth-0.5\lxhang-0.25em\relax][l]{\(D\)．\(4\sqrt{7}/7\)}\par}
\begin{ansblock}[2章-拓-课时14-34-2]
% ans:2章-拓-课时14-34-2
\ansitem{34-2}{\(C\)}
\ansnote{详解}{以 \(O\) 为原点、\(AD\) 所在直线为 \(x\) 轴建系。八等分结合 \(AB=1\)、\(BC=2\)、\(CD=1\)：圆半径 \(R=2\)，双曲线过点 \((\sqrt{2},\sqrt{2})\)，\(a=1\)。 \(2/1-2/b^{2}=1 \rightarrow  b^{2}=2\)，\(c^{2}=1+2=3\)，\(c=\sqrt{3}\)，焦距 \(2\sqrt{3}\)，故选\(C\)。}
\end{ansblock}

\zux{拓展·对称点垂直求渐近线}
\tuhao{34-3} \tieside{难}已知双曲线 \(C\)：\(x^{2}/a^{2}-y^{2}/b^{2}=1\)（\(a>0\)，\(b>0\)）的左焦点为 \(F_{1}\)，\(P\) 为左支上一点，\(|PF_{1}|=a\)，\(P_{0}\) 与 \(P\) 关于原点对称，且 \(P_{0}F_{1}\cdot PF_{1}=0\)，则双曲线的渐近线方程为（　　）
\bindopt {\fontsize{10.5pt}{\qpTJgLeadLiti}\selectfont \makebox[\dimexpr0.5\linewidth-0.5\lxhang-0.25em\relax][l]{\(A\)．\(y=\pm x\)}\makebox[\dimexpr0.5\linewidth-0.5\lxhang-0.25em\relax][l]{\(B\)．\(y=\pm (\sqrt{6}/2)x\)}\par}
\bindopt {\fontsize{10.5pt}{\qpTJgLeadLiti}\selectfont \makebox[\dimexpr0.5\linewidth-0.5\lxhang-0.25em\relax][l]{\(C\)．\(y=\pm (\sqrt{3}/2)x\)}\makebox[\dimexpr0.5\linewidth-0.5\lxhang-0.25em\relax][l]{\(D\)．\(y=\pm 2x\)}\par}
\begin{ansblock}[2章-拓-课时14-34-3]
% ans:2章-拓-课时14-34-3
\ansitem{34-3}{\(B\)}
\ansnote{详解}{\(P_{0}F_{1}\cdot PF_{1}=0 \rightarrow  P_{0}F_{1}\perp PF_{1}\)；\(P_{0}\) 与 \(P\) 关于原点对称 \(\rightarrow  |P_{0}F_{2}|=|PF_{1}|=a\) 且 \(P_{0}F_{1}\perp P_{0}F_{2}\)（\(P_{0}F_{2}\parallel PF_{1}\)）。 \(P_{0}\) 在右支：\(|P_{0}F_{1}|-|P_{0}F_{2}|=2a \rightarrow  |P_{0}F_{1}|=3a\)。 \(Rt\triangle P_{0}F_{1}F_{2}\)（直角在 \(P_{0}\)）：\((3a)^{2}+a^{2}=4c^{2} \rightarrow  10a^{2}=4c^{2}=4(a^{2}+b^{2}) \rightarrow  b^{2}/a^{2}=3/2\)。 渐近线 \(y=\pm (\sqrt{6}/2)x\)，故选\(B\)。}
\end{ansblock}

\zux{拓展·对称点数量积范围}
\tuhao{34-4} \tieside{难}已知点 \(M(0,1)\)，点 \(P\) 是双曲线 \(x^{2}/3-y^{2}=1\) 上的点，点 \(Q\) 是点 \(P\) 关于原点的对称点，则 \(MP\cdot MQ\) 的取值范围是\kongda{\((-\infty ,-2]\)}。
\begin{ansblock}[2章-拓-课时14-34-4]
% ans:2章-拓-课时14-34-4
\ansitem{34-4}{\((-\infty ,-2]\)}
\ansnote{详解}{设 \(P(x_{0},y_{0})\)（\(|x_{0}|\geqslant \sqrt{3}\)），\(Q(-x_{0},-y_{0})\)。 \(MP=(x_{0},y_{0}-1)\)，\(MQ=(-x_{0},-y_{0}-1)\)，\(MP\cdot MQ=-x_{0}^{2}-y_{0}^{2}+1\)。 由 \(y_{0}^{2}=x_{0}^{2}/3-1\)：\(MP\cdot MQ=-x_{0}^{2}-x_{0}^{2}/3+1+1=2-(4/3)x_{0}^{2}\leqslant 2-(4/3)\times 3=-2\)。 取值范围 \((-\infty ,-2]\)。}
\end{ansblock}

\zux{拓展·焦半径比例求e最大}
\tuhao{34-5} \tieside{难}已知双曲线 \(x^{2}/a^{2}-y^{2}/b^{2}=1\)（\(a>0\)，\(b>0\)）的左、右焦点分别为 \(F_{1}\)，\(F_{2}\)，点 \(P\) 在双曲线的右支上，且 \(|PF_{1}|=4|PF_{2}|\)，求双曲线的离心率 \(e\) 的最大值．
\liubai[16mm]{解答书写区}
\begin{ansblock}[2章-拓-课时14-34-5]
% ans:2章-拓-课时14-34-5
\ansitem{34-5}{最大值为 \(5/3\)。}
\ansnote{详解}{\(|PF_{1}|-|PF_{2}|=3|PF_{2}|=2a \rightarrow  |PF_{2}|=(2/3)a\)。 右支焦半径 \(|PF_{2}|\geqslant c-a\)：\((2/3)a\geqslant c-a \rightarrow  c\leqslant (5/3)a \rightarrow  e\leqslant 5/3\)；又 \(e>1\)。 故 \(e\) 的最大值 \(5/3\)。}
\end{ansblock}

\zux{拓展·大教堂外形求方程}
% 图债注：「如图」指源件配图，印面无图（冻片【注】裁定：去装饰图/条件全在文字/尺寸随注——图债登记汇编报告）
\tuhao{34-6} \tieside{难}南非双曲线大教堂由伦敦著名的建筑事务所完成．若将如图所示的双曲线大教堂外形弧线的一段近似看成双曲线 \(y^{2}/a^{2}-x^{2}/b^{2}=1\)（\(a>0\)，\(b>0\)）下支的一部分，且此双曲线过点 \((1,-2)\)，离心率为 \(\sqrt{6}/2\)，则此双曲线的方程为（　　）
\bindopt {\fontsize{10.5pt}{\qpTJgLeadLiti}\selectfont \makebox[\dimexpr0.5\linewidth-0.5\lxhang-0.25em\relax][l]{\(A\)．\(y^{2}/2-x^{2}=1\)}\makebox[\dimexpr0.5\linewidth-0.5\lxhang-0.25em\relax][l]{\(B\)．\(y^{2}/2-x^{2}/3=1\)}\par}
\bindopt {\fontsize{10.5pt}{\qpTJgLeadLiti}\selectfont \makebox[\dimexpr0.5\linewidth-0.5\lxhang-0.25em\relax][l]{\(C\)．\(y^{2}/2-x^{2}/4=1\)}\makebox[\dimexpr0.5\linewidth-0.5\lxhang-0.25em\relax][l]{\(D\)．\(y^{2}/3-x^{2}/3=1\)}\par}
\begin{ansblock}[2章-拓-课时14-34-6]
% ans:2章-拓-课时14-34-6
\ansitem{34-6}{\(A\)}
\ansnote{详解}{\(e^{2}=3/2=c^{2}/a^{2}=(a^{2}+b^{2})/a^{2} \rightarrow  b^{2}=a^{2}/2\)。 过 \((1,-2)\)：\(4/a^{2}-1/b^{2}=1 \rightarrow  4/a^{2}-2/a^{2}=1 \rightarrow  a^{2}=2\)，\(b^{2}=1\)。 方程 \(y^{2}/2-x^{2}=1\)，故选\(A\)。}
\end{ansblock}

\jietitle{2.7.1　抛物线方程（课时15·拓区）}
\zux{拓展·桥拱4溢流孔建模}
% 图债注：「如图」指源件配图，印面无图（冻片【注】裁定：去装饰图/条件全在文字/尺寸随注——图债登记汇编报告）
\tuhao{01} \tieside{难}早在一千多年之前，我国已经把溢流孔用于造桥技术，以减轻桥身重量和水流对桥身的冲击．现设桥拱上有如图所示的 \(4\) 个溢流孔，桥拱和溢流孔的轮廓线均为抛物线的一部分，且 \(4\) 个溢流孔的轮廓线相同．根据图上尺寸，试分别求出这些抛物线的方程，及溢流孔与桥拱交点 \(A\)，\(B\)，\(C\) 的位置．
\liubai[16mm]{解答书写区}
\begin{ansblock}[2章-拓-课时15-01]
% ans:2章-拓-课时15-01
\ansitem{01}{桥拱所在抛物线 \(y=-(1/80)x^{2}\)；溢流孔所在抛物线分别为 \(y=-(5/36)(x+14)^{2}\)、\(y=-(5/36)(x+7)^{2}\)、\(y=-(5/36)(x-7)^{2}\)、\(y=-(5/36)(x-14)^{2}\)；交点 \(A((140/13),-(245/169))\)、\(B(10,-5/4)\)、\(C((70/13),-(245/676))\)。}
\ansnote{详解}{设桥拱 \(y=ax^{2}\)。图注（转置说明）：题面与源文按图代入的是点 \((20,-5)\)（源详解 \(-5=a\times 20^{2}\)，\(dump\) 文字「过点\((-5,20)\)」系转置，标注尺寸以源图为准）： \(-5=a\times 20^{2} \rightarrow  a=-1/80\)，桥拱 \(y=-(1/80)x^{2}\)。① 设 \(A\) 所在溢流孔 \(y=a_{1}(x-14)^{2}\)（顶点 \((14,0)\)），亦过 \((20,-5)\)：\(-5=a_{1}\times 36 \rightarrow  a_{1}=-5/36\)。 \(4\) 孔轮廓相同、顶点 \((\pm 7,0)\)、\((\pm 14,0)\)：\(y=-(5/36)(x\mp 7)^{2}\)、\(y=-(5/36)(x\mp 14)^{2}\)。②③ 联立①②：\(400(x-14)^{2}=36x^{2}\)（两边乘 \(2880\)）\(\rightarrow  91x^{2}-2800x+19600=0\)，\(\Delta =705600\)，\(\sqrt{}\Delta =840 \rightarrow  x=20\)（\(y=-5\)）或 \(x=140/13\)（\(y=-245/169\)），取交点 \(A(140/13,-245/169)\)； 联立①③：\(400(x-7)^{2}=36x^{2} \rightarrow  91x^{2}-1400x+4900=0\)，\(\Delta =176400\)，\(\sqrt{}\Delta =420 \rightarrow  x=10\)（\(y=-5/4\)）或 \(x=70/13\)（\(y=-245/676\)），得 \(C(70/13,-245/676)\)、\(B(10,-5/4)\)。 验算：\(x=140/13\)：\(y=-(1/80)(140/13)^{2}=-19600/(80\times 169)=-245/169\) ；\(x=10\)：\(-100/80=-5/4\) ；\(x=70/13\)：\(-4900/(80\times 169)=-245/676\) 。}
\end{ansblock}

\zux{拓展·阿基米德三角形新定义}
\tuhao{02} \tieside{难}阿基米德（公元前\(287\)年\(\sim \)公元前\(212\)年）是古希腊伟大的物理学家，数学家和天文学家，并享有\("\)数学之神\("\)的称号．他研究抛物线的求积法，得出了著名的阿基米德定理．在该定理中，抛物线的弦与过弦的端点的两切线所围成的三角形被称为\("\)阿基米德三角形\("\)．若抛物线上任意两点 \(A\)，\(B\) 处的切线交于点 \(P\)，则 \(\triangle PAB\) 为\("\)阿基米德三角形\("\)，且当线段 \(AB\) 经过抛物线的焦点 \(F\) 时，\(\triangle PAB\) 具有以下特征：\((1) P\) 点必在抛物线的准线上；\((2) PA\perp PB\)；\((3) PF\perp AB\)．若经过抛物线 \(y^{2}=8x\) 的焦点的一条弦为 \(AB\)，\("\)阿基米德三角形\("\)为 \(\triangle PAB\)，且点 \(P\) 在直线 \(x-y+6=0\) 上，则直线 \(AB\) 的方程为（　　）
\bindopt {\fontsize{10.5pt}{\qpTJgLeadLiti}\selectfont \makebox[\dimexpr0.5\linewidth-0.5\lxhang-0.25em\relax][l]{\(A\)．\(x-y-2=0\)}\makebox[\dimexpr0.5\linewidth-0.5\lxhang-0.25em\relax][l]{\(B\)．\(x-2y-2=0\)}\par}
\bindopt {\fontsize{10.5pt}{\qpTJgLeadLiti}\selectfont \makebox[\dimexpr0.5\linewidth-0.5\lxhang-0.25em\relax][l]{\(C\)．\(x+y-2=0\)}\makebox[\dimexpr0.5\linewidth-0.5\lxhang-0.25em\relax][l]{\(D\)．\(x+2y-2=0\)}\par}
\begin{ansblock}[2章-拓-课时15-02]
% ans:2章-拓-课时15-02
\ansitem{02}{\(A\)}
\ansnote{详解}{\(y^{2}=8x\)：\(p=4\)，准线 \(x=-2\)，焦点 \(F(2,0)\)。 \(P\) 在准线 \(x=-2\) 上，又在 \(x-y+6=0\) 上：\(P(-2,4)\)。 \(k\_PF=(0-4)/(2+2)=-1\)；由特征\((3) PF\perp AB \rightarrow  k\_AB=1\)。 \(AB\) 过 \(F(2,0)\)：\(y=x-2\)，即 \(x-y-2=0\)，故选\(A\)。 验算：\(P(-2,4)\) 代入 \(x-y+6=0\)：\(-2-4+6=0\) 。 （卷③自带详解照录，验算一致。）}
\end{ansblock}

\zux{拓展·正方体轨迹判定多选}
\tuhao{03}\duoxuan\hspace{0.6em} \tieside{难}已知正方体 \(ABCD-A_{1}B_{1}C_{1}D_{1}\) 的棱长为 \(2\)，\(M\) 为 \(DD_{1}\) 的中点，\(N\) 为正方形 \(ABCD\) 所在平面内一动点，则下列命题正确的有（　　）
\lxopt{\(A\)．若 \(MN=2\)，则 \(MN\) 的中点的轨迹所围成图形的面积为 \(\pi \)}
\lxopt{\(B\)．若 \(N\) 到直线 \(BB_{1}\) 与直线 \(DC\) 的距离相等，则 \(N\) 的轨迹为抛物线}
\lxopt{\(C\)．若 \(D_{1}N\) 与 \(AB\) 所成的角为 \(\pi /3\)，则 \(N\) 的轨迹为双曲线}
\lxopt{\(D\)．若 \(\angle ND_{1}B=\pi /6\)，则 \(N\) 的轨迹为椭圆}
\begin{ansblock}[2章-拓-课时15-03]
% ans:2章-拓-课时15-03
\ansitem{03}{\(BCD\)}
\ansnote{详解}{\(A\)：\(DD_{1}\perp \)底面 \(\rightarrow  DN=\sqrt{MN^{2}-MD^{2}}=\sqrt{4-1}=\sqrt{3}\)，\(N\) 的轨迹是以 \(D\) 为圆心、\(\sqrt{3}\) 为半径的圆。 设 \(MN\) 中点 \(P\)、\(MD\) 中点 \(Q\)，则 \(PQ\parallel ND\)、\(PQ=\sqrt{3}/2\)，\(P\) 的轨迹是以 \(Q\) 为圆心 \(\sqrt{3}/2\) 为半径的圆，面积 \(\pi \times 3/4=3\pi /4\neq \pi \)，\(A\) 错； \(B\)：\(N\) 到 \(BB_{1}\) 的距离\(=NB\)（\(BB_{1}\perp \)底面），即 \(N\) 到定点 \(B\) 与定直线 \(DC\) 距离相等（\(B\) 不在 \(DC\) 上），轨迹为抛物线，\(B\) 对； \(C\)：\(AB\parallel D_{1}C_{1}\)，\(D_{1}N\) 与 \(D_{1}C_{1}\) 所成角 \(\pi /3 \rightarrow  D_{1}N\) 为以 \(D_{1}C_{1}\) 为轴、轴截面顶角 \(2\pi /3\) 的圆锥母线，底面 \(ABCD\) 与轴 \(D_{1}C_{1}\) 平行 \(\rightarrow \) 交线为双曲线，\(C\) 对； \(D\)：\(D_{1}N\) 为以 \(D_{1}B\) 为轴、顶角 \(\pi /3\) 的圆锥母线，平面 \(ABCD\) 与轴 \(D_{1}B\) 斜交 \(\rightarrow \) 交线为椭圆，\(D\) 对。 故选 \(BCD\)。 （卷③自带详解照录校订，逐项判定一致。）}
\end{ansblock}

\zux{拓展·正方体轨迹与截面多选}
% 图债注：「如图」指源件配图，印面无图（冻片【注】裁定：去装饰图/条件全在文字/尺寸随注——图债登记汇编报告）
\tuhao{04}\duoxuan\hspace{0.6em} \tieside{难}如图，棱长为 \(2\) 的正方体 \(ABCD-A_{1}B_{1}C_{1}D_{1}\) 中，\(M\)，\(N\) 分别为棱 \(A_{1}D_{1}\)，\(AB\) 中点，\(P\) 为侧面 \(BCC_{1}B_{1}\) 一动点，下列说法正确的是（　　）
\lxopt{\(A\)．异面直线 \(AC\) 与 \(BM\) 所成角的余弦值为 \(\sqrt{2}/3\)}
\lxopt{\(B\)．若 \(\triangle AC_{1}P\) 的面积为 \(\sqrt{3}\)，则动点 \(P\) 的轨迹为椭圆的一部分}
\lxopt{\(C\)．若点 \(P\) 到直线 \(BC\) 与直线 \(C_{1}D_{1}\) 的距离相等，则动点 \(P\) 的轨迹为抛物线的一部分}
\lxopt{\(D\)．过直线 \(MN\) 的平面 \(\alpha \) 与面 \(ABCD\) 所成角最小时，平面 \(\alpha \) 截正方体所得的截面面积为 \(3\sqrt{3}\)}
\begin{ansblock}[2章-拓-课时15-04]
% ans:2章-拓-课时15-04
\ansitem{04}{\(BCD\)}
\ansnote{详解}{以 \(D\) 为原点、\(DA\)、\(DC\)、\(DD_{1}\) 为 \(x\)、\(y\)、\(z\) 轴建系：\(A(2,0,0)\)，\(C(0,2,0)\)，\(B(2,2,0)\)，\(M(1,0,2)\)。 \(A\)：\(AC=(-2,2,0)\)，\(BM=(-1,-2,2)\)，\(|cos\langle AC,BM\rangle |=|AC\cdot BM|/(|AC||BM|)=|2-4+0|/(2\sqrt{2}\times 3)=2/(6\sqrt{2})=\sqrt{2}/6\neq \sqrt{2}/3\)，\(A\) 错； \(B\)：\(S=\)½\(|AC_{1}|\cdot h=\sqrt{3}\)，\(|AC_{1}|=2\sqrt{3} \rightarrow  h=1\)：\(P\) 在以 \(AC_{1}\) 为轴、半径 \(1\) 的圆柱面上，该圆柱面与侧面 \(BCC_{1}B_{1}\)（与轴斜交）的交线为椭圆的一部分，\(B\) 对； \(C\)：\(P\) 在侧面 \(BCC_{1}B_{1}\) 上时，到 \(C_{1}D_{1}\) 的距离\(=|PC_{1}|\)，条件即 \(|PC_{1}|=P\) 到直线 \(BC\) 的距离——到定点与定直线等距且 \(C_{1}\) 不在 \(BC\) 上 \(\rightarrow \) 抛物线的一部分，\(C\) 对； \(D\)：取 \(AD\) 中点 \(H\)，\(MH=2\) 且 \(MH\perp \)底面；当交线与 \(AC\) 平行（\(HN\perp AC\)）时二面角最小，截面为过 \(M\)、\(N\) 及四棱中点的正六边形，边长 \(\sqrt{2}\)，面积 \(6\times (\sqrt{3}/4)\times 2=3\sqrt{3}\)，\(D\) 对。 故选 \(BCD\)。 （卷③自带详解照录校订，逐项判定一致；\(A\) 项 \(|AC|=2\sqrt{2}\)、\(|BM|=3\)，余弦 \(2/(6\sqrt{2})=\sqrt{2}/6\) 。）}
\end{ansblock}

\zux{拓展·折叠最值（椭圆轨迹）}
% 图债注：「如图」指源件配图，印面无图（冻片【注】裁定：去装饰图/条件全在文字/尺寸随注——图债登记汇编报告）
\tuhao{05} \tieside{难}如图，在平面四边形 \(ABCD\) 中，满足 \(AB=BC\)，\(CD=AD\)，且 \(AB+AD=10\)，\(BD=8\)，沿着 \(BD\) 把 \(\triangle ABD\) 折起，使点 \(A\) 到达点 \(P\) 的位置，且使 \(PC=2\)，则三棱锥 \(P-BCD\) 体积的 \(max\) 为（　　）
\bindopt {\fontsize{10.5pt}{\qpTJgLeadLiti}\selectfont \makebox[\dimexpr0.5\linewidth-0.5\lxhang-0.25em\relax][l]{\(A\)．\(12\)}\makebox[\dimexpr0.5\linewidth-0.5\lxhang-0.25em\relax][l]{\(B\)．\(12\sqrt{2}\)}\par}
\bindopt {\fontsize{10.5pt}{\qpTJgLeadLiti}\selectfont \makebox[\dimexpr0.5\linewidth-0.5\lxhang-0.25em\relax][l]{\(C\)．\(16\sqrt{2}/3\)}\makebox[\dimexpr0.5\linewidth-0.5\lxhang-0.25em\relax][l]{\(D\)．\(16/3\)}\par}
\begin{ansblock}[2章-拓-课时15-05]
% ans:2章-拓-课时15-05
\ansitem{05}{\(C\)}
\ansnote{详解}{折叠不变量：\(PB=AB\)，\(PD=AD \rightarrow  PB+PD=10>BD=8\)，\(P\) 在以 \(B\)、\(D\) 为焦点、长轴 \(10\)、焦距 \(8\) 的椭圆上（不在长轴端点）。 \(\triangle BPD\cong \triangle BCD\)（三边相等）\(\rightarrow  \angle PBD=\angle CBD\)。过 \(P\) 作 \(PE\perp BD\) 于 \(E\)，连 \(CE\)：\(\triangle BPE\cong \triangle BCE \rightarrow  CE\perp BD\) 且 \(PE=CE \rightarrow  BD\perp \)平面 \(PCE\)。 \(V(P-BCD)=V(B-PCE)+V(D-PCE)=(1/3)S\triangle PCE\cdot BD=(8/3)S\triangle PCE\)。 取 \(PC\) 中点 \(F\)，\(EF\perp PC\)（\(PE=CE\)）：\(S\triangle PCE=(1/2)\cdot PC\cdot EF=1\cdot \sqrt{PE^{2}-1}\)（\(PC=2 \rightarrow  CF=1\)）。 \(PE\) 的 \(max=\)椭圆短半轴\(=\sqrt{5^{2}-4^{2}}=3 \rightarrow  S\triangle PCE max=\sqrt{9-1}=2\sqrt{2}\)。 \(V max=(8/3)\times 2\sqrt{2}=16\sqrt{2}/3\)，故选\(C\)。 验算：\(S\triangle PCE=\sqrt{PE^{2}-1}\)：\(EF=\sqrt{PE^{2}-CF^{2}}=\sqrt{PE^{2}-1}\) ；\(PE=3 \rightarrow  EF=2\sqrt{2}\)，\(S=\)½\(\cdot 2\cdot 2\sqrt{2}=2\sqrt{2}\) 。 （卷③自带详解照录，验算一致。）}
\end{ansblock}

\zux{拓展·二面角=线面角→抛物线轨迹}
\tuhao{06}\duoxuan\hspace{0.6em} \tieside{难}已知长方体 \(ABCD-A_{1}B_{1}C_{1}D_{1}\) 中，\(AB=BC=2\)，\(AA_{1}=2\sqrt{2}\)，点 \(P\) 是四边形 \(A_{1}B_{1}C_{1}D_{1}\) 内（包含边界）的一动点，设二面角 \(P-AD-B\) 的大小为 \(\alpha \)，直线 \(PB\) 与平面 \(ABCD\) 所成的角为 \(\beta \)，若 \(\alpha =\beta \)，则（　　）
\lxopt{\(A\)．点 \(P\) 的轨迹为一条抛物线}
\lxopt{\(B\)．线段 \(PB\) 长的最小值为 \(3\)}
\lxopt{\(C\)．直线 \(PA_{1}\) 与直线 \(CD\) 所成角的最大值为 \(\pi /4\)}
\lxopt{\(D\)．三棱锥 \(P-A_{1}BC_{1}\) 体积的最大值为 \(2\sqrt{2}/3\)}
\begin{ansblock}[2章-拓-课时15-06]
% ans:2章-拓-课时15-06
\ansitem{06}{\(BCD\)}
\ansnote{详解}{过 \(P\) 作 \(PO\perp \)平面 \(ABCD\) 于 \(O\)，\(OH\perp AD\) 于 \(H\)。 \(AD\perp \)平面 \(POH \rightarrow  AD\perp PH \rightarrow  \angle PHO=\alpha \)；又 \(\angle PBO=\beta \)，\(\alpha =\beta  \rightarrow  OH=OB\)。 \(O\) 的轨迹：底面内到定点 \(B\) 与定直线 \(AD\) 等距的点集——以 \(B\) 为焦点、\(AD\) 为准线的抛物线在四边形 \(ABCD\) 内（含边界）的部分；\(P\) 的轨迹为其上平移（以 \(B_{1}\) 为焦点、\(A_{1}D_{1}\) 为准线），是抛物线的一部分而非整条，\(A\) 错； \(B\)：\(|PB|^{2}=(2\sqrt{2})^{2}+|OB|^{2}\)，\(O\) 为 \(AB\) 中点时 \(|OB|min=1 \rightarrow  |PB|min=\sqrt{8+1}=3\)，\(B\) 对； \(C\)：\(PA_{1}\) 与 \(CD\) 所成角即 \(OA\) 与 \(AB\) 所成角 \(\angle OAB\)。以 \(AB\) 中点 \(M\) 为原点建系，抛物线 \(y^{2}=4x\)，\(A(-1,0)\)。设切线 \(x=ty-1\)：代入得 \(y^{2}-4ty+4=0\)，\(\Delta =16t^{2}-16=0 \rightarrow  t=\pm 1\)，取 \(t=1\)（\(O\) 在四边形内）\(\rightarrow  \angle OAB=\pi /4\)，最大值 \(\pi /4\)，\(C\) 对； \(D\)：\(V(P-A_{1}BC_{1})=V(B-PA_{1}C_{1})=(1/3)S\triangle PA_{1}C_{1}\cdot BB_{1}\)，\(A_{1}C_{1}=2\sqrt{2}\)。\(P\) 到 \(A_{1}C_{1}\) 距离最大 \(\Leftrightarrow  O\) 到 \(AC\) 距离最大：\(O\) 为 \(AB\) 中点时最大，值 \((1/4)|BD|=\sqrt{2}/2\)。 \(Vmax=(2\sqrt{2}/3)\times (1/2)\times 2\sqrt{2}\times (\sqrt{2}/2)=2\sqrt{2}/3\)，\(D\) 对。 故选 \(BCD\)。 验算：\(y^{2}=4x\)（焦点 \(B(1,0)\)，准线 \(x=-1=AD\) 所在直线）：\(|OB|\) 与 \(|OH|\) 分别为到焦点、准线距离，等距即轨迹 ；切线 \(y=x+1\) 过 \(A(-1,0)\)，斜率\(1 \rightarrow  \angle OAB=45^{\circ}\) 。 （卷③自带详解照录校订，逐项判定一致。）}
\end{ansblock}

\zux{拓展·含参开壳焦点（y²=4ax）}
\tuhao{07} \tieside{难}写出抛物线 \(y^{2}=4ax\) 的焦点坐标．
\liubai[16mm]{解答书写区}
\begin{ansblock}[2章-拓-课时15-07]
% ans:2章-拓-课时15-07
\ansitem{07}{\((a,0)\)（\(a\neq 0\)）。}
\ansnote{详解}{按标准形 \(y^{2}=2px\) 对照：\(2p=4a \rightarrow  p=2a\)，焦点 \((p/2,0)=(a,0)\)。 \(a>0\)：开口向右，焦点 \((a,0)\)；\(a<0\)：开口向左，焦点仍为 \((a,0)\)（\(a<0\) 时即 \(x\) 轴负半轴）。 故焦点 \((a,0)\)，\(a\neq 0\)。 验算：\(a=-1\)：\(y^{2}=-4x\)，焦点 \((-1,0)\) 。 （本块自撰亲算；含参开壳，教材未限 \(a\) 正负，按两种开口统一写 \((a,0)\)。）}
\end{ansblock}

\zux{拓展·y=2x² 与 y=ax² 的焦点与准线}
\tuhao{08} \tieside{难}写出下列抛物线的焦点坐标和准线方程：
\((1) y=2x^{2}\)；\((2) y=ax^{2}\)（\(a\neq 0\)）．
\liubai[16mm]{解答书写区}
\begin{ansblock}[2章-拓-课时15-08]
% ans:2章-拓-课时15-08
\ansitem{08}{\((1)\) 焦点 \((0,1/8)\)，准线 \(y=-1/8\)；\((2)\) 焦点 \((0,1/(4a))\)，准线 \(y=-1/(4a)\)（\(a\neq 0\)）。}
\ansnote{详解}{\((1) y=2x^{2} \rightarrow  x^{2}=(1/2)y=2\cdot (1/4)y\)：\(2p=1/2\)，\(p=1/4\)。 开口向上：焦点 \((0,1/8)\)，准线 \(y=-1/8\)。 \((2) y=ax^{2} \rightarrow  x^{2}=(1/a)y=2\cdot (1/(2a))y\)：\(2p=1/a\)，\(p=1/(2a)\)。 \(a>0\)：开口向上，焦点 \((0,1/(4a))\)，准线 \(y=-1/(4a)\)；\(a<0\)：开口向下，焦点、准线表达式不变（\(1/(4a)<0\)，准线在 \(x\) 轴上方）。 验算：\(a=2\) 时与 \((1)\) 一致：\(1/(4\times 2)=1/8\) 。 （本块自撰亲算。）}
\end{ansblock}

\zux{拓展·路灯支架建模}
% 图债注：「如图」指源件配图，印面无图（冻片【注】裁定：去装饰图/条件全在文字/尺寸随注——图债登记汇编报告）
\tuhao{09} \tieside{难}某城市在主干道统一安装了一种新型节能路灯，该路灯由灯柱和支架组成．在如图所示的平面直角坐标系中，支架 \(ACB\) 是抛物线 \(y^{2}=2x\) 的一部分，灯柱 \(CD\) 经过该抛物线的焦点 \(F\) 且与路面垂直，其中 \(B\) 为抛物线的顶点，\(DH\) 表示道路路面，\(BF\parallel DH\)，\(A\) 为锥形灯罩的顶，灯罩轴线与抛物线在 \(A\) 处的切线垂直．安装时，要求锥形灯罩的顶到灯柱所在直线的距离是 \(1.5m\)，灯罩的轴线正好通过道路路面中的中线．
\((1)\) 求灯罩轴线所在的直线方程；
\((2)\) 若路宽为 \(10m\)，求灯柱的高．
\liubai[24mm]{解答书写区}
\begin{ansblock}[2章-拓-课时15-09]
% ans:2章-拓-课时15-09
\ansitem{09}{\((1) y=-2x+6\)；\((2) 6m\)。}
\ansnote{详解}{\((1) |BF|=p/2=1/2\)，灯柱 \(CD\) 过焦点 \(F(1/2,0)\) 且 \(A\) 到灯柱距离 \(1.5 \rightarrow  x\_A=1/2+3/2=2\)。 代入 \(y^{2}=2x\)：\(y\_A=2\)（取上半支），\(A(2,2)\)。 设 \(A\) 处切线 \(y-2=k(x-2)\)：与 \(y^{2}=2x\) 联立消 \(x\) 得 \(ky^{2}-2y+4-4k=0\)，相切 \(\Delta =4-4k(4-4k)=0 \rightarrow  16k^{2}-16k+4=0 \rightarrow  k=1/2\)。 灯罩轴线\(\perp \)切线：斜率 \(-2\)，方程 \(y-2=-2(x-2)\)，即 \(y=-2x+6\)。 \((2)\) 路宽 \(|DH|=10\)，灯罩轴线过路面中线 \(\rightarrow \) 轴线与路面交点到 \(y\) 轴距离 \(1/2+5=11/2\)。 对 \(y=-2x+6\)，\(x=11/2\) 时 \(y=-5 \rightarrow \) 交点 \((11/2,-5)\)，\(|FD|=5\)。 又 \(x=1/2\) 代入 \(y^{2}=2x \rightarrow  y=\pm 1\)，\(|CF|=1\)。 灯柱高 \(|CD|=|DF|+|FC|=5+1=6(m)\)。 验算：切线斜率方程 \(16k^{2}-16k+4=4(2k-1)^{2}=0\) ；\(|CD|=6\) 。}
\end{ansblock}

\zux{拓展·3x²+4y=0 焦点准线}
\tuhao{10} \tieside{难}抛物线 \(3x^{2}+4y=0\) 的焦点坐标为\kongda{\((0,-1/3)\)；\(y=1/3\)。}，准线方程是\(\kongbai{}\)．
\begin{ansblock}[2章-拓-课时15-10]
% ans:2章-拓-课时15-10
\ansitem{10}{\((0,-1/3)\)；\(y=1/3\)。}
\ansnote{详解}{\(3x^{2}+4y=0 \rightarrow  x^{2}=-(4/3)y=-2\cdot (2/3)y\)：\(2p=4/3\)，\(p=2/3\)。 开口向下：焦点 \((0,-p/2)=(0,-1/3)\)，准线 \(y=p/2=1/3\)。 验算：顶点到焦点、到准线等距 \(1/3\) 。}
\end{ansblock}

\zux{拓展·准线 y=−2 求方程}
\tuhao{11} \tieside{难}已知抛物线的准线方程为 \(y=-2\)，求抛物线的标准方程．
\liubai[16mm]{解答书写区}
\begin{ansblock}[2章-拓-课时15-11]
% ans:2章-拓-课时15-11
\ansitem{11}{\(x^{2}=8y\)}
\ansnote{详解}{准线 \(y=-2\) 在 \(x\) 轴下方 \(\rightarrow \) 开口向上：\(x^{2}=2py\)（\(p>0\)），\(-p/2=-2 \rightarrow  p=4\)，\(2p=8\)。 方程 \(x^{2}=8y\)。 验算：焦点 \((0,4)\)、准线 \(y=-2\) 。}
\end{ansblock}

\zux{拓展·椭圆中心顶点焦点抛物线}
\tuhao{12} \tieside{难}以椭圆 \(x^{2}/2+y^{2}=1\) 的对称中心为顶点，椭圆的焦点为焦点的抛物线的方程是（　　）．
\lxopt{\(A\)．\(y^{2}=4x\)}
\lxopt{\(B\)．\(y^{2}=-4x\) 或 \(x^{2}=4y\)}
\lxopt{\(C\)．\(x^{2}=4y\)}
\lxopt{\(D\)．\(y^{2}=4x\) 或 \(y^{2}=-4x\)}
\begin{ansblock}[2章-拓-课时15-12]
% ans:2章-拓-课时15-12
\ansitem{12}{\(D\)}
\ansnote{详解}{椭圆 \(x^{2}/2+y^{2}=1\)：\(a^{2}=2\)、\(b^{2}=1 \rightarrow  c^{2}=1\)，焦点 \((\pm 1,0)\)，对称中心 \(O(0,0)\)。 以 \(O\) 为顶点、焦点为 \((1,0)\) 或 \((-1,0)\)：\(p/2=1 \rightarrow  2p=4 \rightarrow  y^{2}=4x\) 或 \(y^{2}=-4x\)，故选\(D\)。 验算：\(y^{2}=\pm 4x\) 的焦点 \((\pm 1,0)\) 。}
\end{ansblock}

\zux{拓展·过点距准线定方程双解}
\tuhao{13} \tieside{难}已知抛物线 \(y^{2}=2px\)（\(p>0\)）经过点 \(M(x_{0},2\sqrt{2})\)，若点 \(M\) 到准线 \(l\) 的距离为 \(3\)，则该抛物线的方程为（　　）
\lxopt{\(A\)．\(y^{2}=4x\)}
\lxopt{\(B\)．\(y^{2}=2x\) 或 \(y^{2}=4x\)}
\lxopt{\(C\)．\(y^{2}=8x\)}
\lxopt{\(D\)．\(y^{2}=4x\) 或 \(y^{2}=8x\)}
\begin{ansblock}[2章-拓-课时15-13]
% ans:2章-拓-课时15-13
\ansitem{13}{\(D\)}
\ansnote{详解}{\((2\sqrt{2})^{2}=2px_{0} \rightarrow  x_{0}=4/p\)。 \(M\) 到准线距离\(=x_{0}+p/2=4/p+p/2=3 \rightarrow  p^{2}-6p+8=0 \rightarrow  p=2\) 或 \(p=4\)。 抛物线 \(y^{2}=4x\) 或 \(y^{2}=8x\)，故选\(D\)。 验算：\(p=2\)：\(x_{0}=2\)，\(M(2,2\sqrt{2})\)，到准线 \(x=-1\) 距离 \(3\) ；\(p=4\)：\(x_{0}=1\)，\(M(1,2\sqrt{2})\)，到准线 \(x=-2\) 距离 \(3\) 。}
\end{ansblock}

\zux{拓展·开口大小与系数（作图观察）}
\tuhao{14} \tieside{难}在同一平面直角坐标系中画出下列抛物线．
\((1) y^{2}=x\)；\((2) y^{2}=2x\)；\((3) y^{2}=4x\)．
通过观察这些图形，说明抛物线开口的大小与方程中 \(x\) 的系数有怎样的关系．
\liubai[16mm]{解答书写区}
\begin{ansblock}[2章-拓-课时15-14]
% ans:2章-拓-课时15-14
\ansitem{14}{答案见解析：\(x\) 的系数为正数且越大时，抛物线的开口向右且开口越大。}
\ansnote{详解}{三条抛物线同以 \(x\) 轴为对称轴、顶点在原点、开口向右。作图观察（或取同一 \(y\) 值比较横坐标 \(x=y^{2}/(2p)\)）：系数 \(1\)、\(2\)、\(4\) 递增时，同一高度处横坐标 \(1/(2p)\cdot y^{2}\) 递减，曲线横向收拢更快、开口更大。 故：\(x\) 的系数为正且越大，开口向右且开口越大。 （本块自撰亲算补足说理；作图观察题，答案见解析。）}
\end{ansblock}

\zux{拓展·到焦点距离恒大于1求p范围}
\tuhao{15} \tieside{难}若抛物线 \(y^{2}=2px\)（\(p>0\)）上任意一点到焦点的距离恒大于 \(1\)，则 \(p\) 的取值范围是（　　）
\bindopt {\fontsize{10.5pt}{\qpTJgLeadLiti}\selectfont \makebox[\dimexpr0.25\linewidth-0.25\lxhang-0.25em\relax][l]{\(A\)．\(p<1\)}\makebox[\dimexpr0.25\linewidth-0.25\lxhang-0.25em\relax][l]{\(B\)．\(p>1\)}\makebox[\dimexpr0.25\linewidth-0.25\lxhang-0.25em\relax][l]{\(C\)．\(p<2\)}\makebox[\dimexpr0.25\linewidth-0.25\lxhang-0.25em\relax][l]{\(D\)．\(p>2\)}\par}
\begin{ansblock}[2章-拓-课时15-15]
% ans:2章-拓-课时15-15
\ansitem{15}{\(D\)}
\ansnote{详解}{由定义，点到焦点距离\(=\)点到准线 \(x=-p/2\) 的距离，其最小值在顶点处取得，为 \(p/2\)。 恒大于 \(1 \Leftrightarrow  p/2>1 \Leftrightarrow  p>2\)，故选\(D\)。 验算：\(p=2\) 时顶点到焦点距离\(=1\)，不满足\("\)恒大于\("\)  开区间正确。}
\end{ansblock}

\jietitle{2.7.2　抛物线的性质（课时16·拓区）}
% 【16-18｜让位席（教材版让位）无键留块——印面跳号（17 后直接 19），注记随席（台账抄件）。】
\zux{拓展·曲线上点二次和最小}
\tuhao{01} \tieside{难}已知点 \((x\)，\(y)\) 在抛物线 \(y^{2}=4x\) 上，则 \(z=x^{2}+(1/2)y^{2}+3\) 的最小值是（　）
\bindopt {\fontsize{10.5pt}{\qpTJgLeadLiti}\selectfont \makebox[\dimexpr0.25\linewidth-0.25\lxhang-0.25em\relax][l]{\(A\)．\(2\)}\makebox[\dimexpr0.25\linewidth-0.25\lxhang-0.25em\relax][l]{\(B\)．\(3\)}\makebox[\dimexpr0.25\linewidth-0.25\lxhang-0.25em\relax][l]{\(C\)．\(4\)}\makebox[\dimexpr0.25\linewidth-0.25\lxhang-0.25em\relax][l]{\(D\)．\(0\)}\par}
\begin{ansblock}[2章-拓-课时16-01]
% ans:2章-拓-课时16-01
\ansitem{01}{\(B\)}
\ansnote{详解}{点在 \(y^{2}=4x\) 上，故 \(x\geqslant 0\) 且 \(y^{2}=4x\)。代入得 \(z=x^{2}+(1/2)\cdot 4x+3=x^{2}+2x+3=(x+1)^{2}+2\)。由 \(x\geqslant 0\)，\(x=0\) 时 \(z\) 最小，最小值为 \(3\)。故选：\(B\)。〔卷③自带详解照录校订〕}
\end{ansblock}

\zux{拓展·|PM|+|PF|最小反求p}
\tuhao{02} \tieside{难}点 \(M(20,40)\)，抛物线 \(y^{2}=2px\)（\(p>0\)）的焦点为 \(F\)，若对于抛物线上的任意点 \(P\)，\(|PM|+|PF|\) 的最小值为 \(41\)，则 \(p\) 的值等于\kongda{\(42\) 或 \(22\)}．
\begin{ansblock}[2章-拓-课时16-02]
% ans:2章-拓-课时16-02
\ansitem{02}{\(42\) 或 \(22\)}
\ansnote{详解}{按 \(M\) 与抛物线的位置分类。\((1) M\) 在抛物线内部或上：\(40^{2}\leqslant 2p\cdot 20\)，即 \(p\geqslant 40\)。设 \(P\) 在准线 \(x=-p/2\) 上的射影为 \(D\)，\(|PF|=|PD|\)，\(M\)、\(P\)、\(D\) 共线时取最小，最小值为 \(M\) 到准线的距离 \(20+p/2=41\)，得 \(p=42\)（\(\geqslant 40\) ）。\((2) M\) 在抛物线外部：\(p<40\)，\(P\)、\(M\)、\(F\) 共线时取最小，最小值为 \(|MF|=\sqrt{}(40^{2}+(20-p/2)^{2})=41\)，\((20-p/2)^{2}=81\)、\(p/2=11\) 或 \(29\)，\(p=22\) 或 \(p=58\)（\(58>40\) 与外部前提矛盾，舍）。综上 \(p=42\) 或 \(22\)。〔卷③自带详解照录校订；内外分类口径随源〕}
\end{ansblock}

\zux{拓展·|PA|+|PQ|最小}
\tuhao{03} \tieside{难}已知点 \(P\) 是抛物线 \(x^{2}=4y\) 上的动点，点 \(P\) 在 \(x\) 轴上的射影是点 \(Q\)，点 \(A\) 的坐标是 \((8,7)\)，则 \(|PA|+|PQ|\) 的最小值为（　）
\bindopt {\fontsize{10.5pt}{\qpTJgLeadLiti}\selectfont \makebox[\dimexpr0.25\linewidth-0.25\lxhang-0.25em\relax][l]{\(A\)．\(7\)}\makebox[\dimexpr0.25\linewidth-0.25\lxhang-0.25em\relax][l]{\(B\)．\(8\)}\makebox[\dimexpr0.25\linewidth-0.25\lxhang-0.25em\relax][l]{\(C\)．\(9\)}\makebox[\dimexpr0.25\linewidth-0.25\lxhang-0.25em\relax][l]{\(D\)．\(10\)}\par}
\begin{ansblock}[2章-拓-课时16-03]
% ans:2章-拓-课时16-03
\ansitem{03}{\(C\)}
\ansnote{详解}{\(x^{2}=4y\) 的焦点 \(F(0,1)\)，准线 \(y=-1\)。延长 \(QP\) 交准线于 \(M\)，则 \(|PM|=y\_P+1=|PQ|+1\)；由定义 \(|PM|=|PF|\)，故 \(|PQ|=|PF|-1\)。于是 \(|PA|+|PQ|=|PA|+|PF|-1\geqslant |AF|-1=\sqrt{}(8^{2}+(7-1)^{2})-1=10-1=9\)，当且仅当 \(A\)、\(P\)、\(F\) 三点共线时取等（\(A(8,7)\) 在抛物线外：\(64>4\cdot 7=28\)，线段 \(AF\) 与抛物线相交，取等可行）。最小值为 \(9\)，选 \(C\)。〔卷③自带详解照录校订〕}
\end{ansblock}

\zux{拓展·双根式和最小}
\tuhao{04} \tieside{难}已知点 \(P(m,n)\) 是抛物线 \(y=-(1/4)x^{2}\) 上一动点，则 \(\sqrt{}(m^{2}+(n+1)^{2})+\sqrt{}((m-4)^{2}+(n+5)^{2})\) 的最小值为（　）
\bindopt {\fontsize{10.5pt}{\qpTJgLeadLiti}\selectfont \makebox[\dimexpr0.25\linewidth-0.25\lxhang-0.25em\relax][l]{\(A\)．\(4\)}\makebox[\dimexpr0.25\linewidth-0.25\lxhang-0.25em\relax][l]{\(B\)．\(5\)}\makebox[\dimexpr0.25\linewidth-0.25\lxhang-0.25em\relax][l]{\(C\)．\(\sqrt{30}\)}\makebox[\dimexpr0.25\linewidth-0.25\lxhang-0.25em\relax][l]{\(D\)．\(6\)}\par}
\begin{ansblock}[2章-拓-课时16-04]
% ans:2章-拓-课时16-04
\ansitem{04}{\(D\)}
\ansnote{详解}{化标准形 \(x^{2}=-4y\)：焦点 \(F(0,-1)\)，准线 \(l\)：\(y=1\)。第一根式是 \(P\) 到 \(F(0,-1)\) 的距离 \(|PF|\)，第二根式是 \(P\) 到 \(A(4,-5)\) 的距离 \(|PA|\)。由定义 \(|PF|\) 等于 \(P\) 到准线 \(y=1\) 的距离（\(n\leqslant 0\)，即 \(1-n\)），故 \(|PF|+|PA|\) 不小于 \(A\) 到准线的垂线段长 \(1-(-5)=6\)（\(A\) 在抛物线内部：\(4^{2}=16<-4\cdot (-5)=20\) ，\(P\) 取垂线段与抛物线交点即可取等）。最小值 \(6\)，选 \(D\)。〔卷③自带详解照录校订〕}
\end{ansblock}

\zux{拓展·阿氏圆构造最小值}
\tuhao{05} \tieside{难}已知 \(A(2,0)\)，\(B(6,0)\)，\(|BM|=2\)，点 \(N\) 在抛物线 \(y^{2}=8x\) 上，则 \(|MN|+(1/2)|MA|\) 的最小值为（　）
\bindopt {\fontsize{10.5pt}{\qpTJgLeadLiti}\selectfont \makebox[\dimexpr0.25\linewidth-0.25\lxhang-0.25em\relax][l]{\(A\)．\(6\)}\makebox[\dimexpr0.25\linewidth-0.25\lxhang-0.25em\relax][l]{\(B\)．\(2\sqrt{5}\)}\makebox[\dimexpr0.25\linewidth-0.25\lxhang-0.25em\relax][l]{\(C\)．\(5\)}\makebox[\dimexpr0.25\linewidth-0.25\lxhang-0.25em\relax][l]{\(D\)．\(2\sqrt{6}\)}\par}
\begin{ansblock}[2章-拓-课时16-05]
% ans:2章-拓-课时16-05
\ansitem{05}{\(D\)}
\ansnote{详解}{\(M(x,y)\) 在以 \(B(6,0)\) 为圆心、\(2\) 为半径的圆上：\((x-6)^{2}+y^{2}=4\)；\(N(x_{0},y_{0})\) 满足 \(y_{0}^{2}=8x_{0}\)（\(x_{0}\geqslant 0\)）。在圆约束下加权配凑：\((1/4)|MA|^{2}=(1/4)[(x-2)^{2}+y^{2}]=(1/4)[(x-2)^{2}+y^{2}]+(3/4)[(x-6)^{2}+y^{2}-4]=(x-5)^{2}+y^{2}\)（所加项在圆上为 \(0\)），即 \((1/2)|MA|=|MC|\)，其中 \(C(5,0)\)。于是 \(|MN|+(1/2)|MA|=|MN|+|MC|\geqslant |CN|=\sqrt{}((x_{0}-5)^{2}+y_{0}^{2})=\sqrt{}((x_{0}-5)^{2}+8x_{0})=\sqrt{}((x_{0}-1)^{2}+24)\geqslant 2\sqrt{6}\)，当且仅当 \(x_{0}=1\)（\(N(1,\pm 2\sqrt{2})\)）且 \(C\)、\(M\)、\(N\) 三点共线（\(M\) 为线段 \(CN\) 与圆的交点）时取等——取等可行：\(C(5,0)\) 在圆内（\(|CB|=1<2\)）而 \(N\) 在圆外（\(|NB|^{2}=(x_{0}-6)^{2}+8x_{0}=(x_{0}-2)^{2}+32>4\) 恒成立），线段 \(CN\) 必与圆相交。最小值 \(2\sqrt{6}\)，选 \(D\)。〔卷③自带详解照录校订；配凑恒等式本件复核：¼\([(x-2)^{2}+y^{2}]-[(x-5)^{2}+y^{2}]=-\)¾\([(x-6)^{2}+y^{2}-4]=0\) 于圆上成立 〕}
\end{ansblock}

\zux{拓展·动点张角最大}
\tuhao{06} \tieside{难}已知点 \(P\) 为抛物线 \(y^{2}=4x\) 上一动点，\(A(1,0)\)，\(B(3,0)\)，则 \(\angle APB\) 的最大值为（　）
\bindopt {\fontsize{10.5pt}{\qpTJgLeadLiti}\selectfont \makebox[\dimexpr0.25\linewidth-0.25\lxhang-0.25em\relax][l]{\(A\)．\(\pi /6\)}\makebox[\dimexpr0.25\linewidth-0.25\lxhang-0.25em\relax][l]{\(B\)．\(\pi /4\)}\makebox[\dimexpr0.25\linewidth-0.25\lxhang-0.25em\relax][l]{\(C\)．\(\pi /3\)}\makebox[\dimexpr0.25\linewidth-0.25\lxhang-0.25em\relax][l]{\(D\)．\(\pi /2\)}\par}
\begin{ansblock}[2章-拓-课时16-06]
% ans:2章-拓-课时16-06
\ansitem{06}{\(B\)}
\ansnote{详解}{由对称性不妨设 \(P(x,y)\)（\(y>0\)）。特殊位：\(x=1\) 时 \(P(1,2)\)，\(PA\perp x\) 轴，\(tan\angle APB=|AB|/|PA|=2/2=1\)，\(\angle APB=\pi /4\)；\(x=3\) 时 \(P(3,2\sqrt{3})\)，\(PB\perp x\) 轴，\(tan\angle APB=|AB|/|PB|=2/(2\sqrt{3})=\sqrt{3}/3\)，\(\angle APB=\pi /6\)。一般位 \(x\neq 1\) 且 \(x\neq 3\)：\(tan\angle APB=(k\_PB-k\_PA)/(1+k\_PA\cdot k\_PB)\)，其中 \(k\_PA=y/(x-1)\)、\(k\_PB=y/(x-3)\)，代入化简得 \(tan\angle APB=2y/(x^{2}-4x+3+y^{2})\)；以 \(y^{2}=4x\) 代 \(x=y^{2}/4\)：\(tan\angle APB=2y/(y^{4}/16+3)=2/((1/16)y^{3}+3/y)=2/((1/16)y^{3}+1/y+1/y+1/y)\)。分母四项用均值不等式：\((1/16)y^{3}+1/y+1/y+1/y \allowbreak \geqslant  4\cdot ((1/16)y^{3}\cdot (1/y)^{3})^{1/4}=4\cdot (1/16)^{1/4}=2\)，故 \(tan\angle APB\leqslant 2/2=1\)；等号要求 \((1/16)y^{3}=1/y\) 即 \(y=2\)（\(P(1,2)\)，落在 \(x=1\) 特殊位，不在一般位内），故一般位 \(0<tan\angle APB<1\)、\(0<\angle APB<\pi /4\)。综上 \(\angle APB\) 最大值 \(\pi /4\)，选 \(B\)。〔卷③自带详解照录校订；装配注：源解用两直线到角公式，教师版可改向量夹角 \(cos\angle APB=(PA\cdot PB)/(|PA||PB|)\) 路径，结论不变〕}
\end{ansblock}

\zux{拓展·公共弦长求方程}
\tuhao{07} \tieside{难}已知抛物线的顶点为坐标原点，对称轴为 \(x\) 轴，且与圆 \(x^{2}+y^{2}=4\) 相交的公共弦长为 \(2\sqrt{3}\)，则抛物线的方程为\kongda{\(y^{2}=3x\) 或 \(y^{2}=-3x\)}．
\begin{ansblock}[2章-拓-课时16-07]
% ans:2章-拓-课时16-07
\ansitem{07}{\(y^{2}=3x\) 或 \(y^{2}=-3x\)}
\ansnote{详解}{设方程为 \(y^{2}=2px\) 或 \(y^{2}=-2px\)（\(p>0\)），交点 \(A(x_{1},y_{1})\)、\(B(x_{2},y_{2})\)（\(y_{1}>0>y_{2}\)）。两曲线均关于 \(x\) 轴对称，公共弦被 \(x\) 轴垂直平分：\(x_{1}=x_{2}\)、\(y_{2}=-y_{1}\)，弦长 \(y_{1}-y_{2}=2y_{1}=2\sqrt{3} \Rightarrow  y_{1}=\sqrt{3}\)。代入圆：\(x_{1}^{2}=4-3=1\)、\(x_{1}=\pm 1\)。点 \((1,\sqrt{3})\) 在 \(y^{2}=2px\) 上：\(3=2p\)、\(p=3/2 \Rightarrow  y^{2}=3x\)；点 \((-1,\sqrt{3})\) 在 \(y^{2}=-2px\) 上：同理 \(p=3/2 \Rightarrow  y^{2}=-3x\)。故方程为 \(y^{2}=3x\) 或 \(y^{2}=-3x\)。〔卷③自带详解照录校订〕}
\end{ansblock}

\zux{拓展·圆与抛物线劣弧动点}
\tuhao{08}\duoxuan\hspace{0.6em} \tieside{难}已知抛物线 \(E\)：\(x^{2}=4y\) 的焦点为 \(F\)，圆 \(C\)：\(x^{2}+(y-1)^{2}=16\) 与抛物线 \(E\) 交于 \(A\)，\(B\) 两点，点 \(P\) 为劣弧 \(AB\) 上不同于 \(A\)，\(B\) 的一个动点，过点 \(P\) 作平行于 \(y\) 轴的直线 \(l\) 交抛物线 \(E\) 于点 \(N\)，则（　）
\lxopt{\(A\)．点 \(P\) 的纵坐标的取值范围是 \((2\sqrt{3},5)\)}
\lxopt{\(B\)．\(PN+NF\) 等于点 \(P\) 到抛物线 \(E\) 的准线的距离}
\lxopt{\(C\)．圆 \(C\) 的圆心到抛物线 \(E\) 的准线的距离为 \(2\)}
\lxopt{\(D\)．\(\triangle PFN\) 周长的取值范围是 \((8,10)\)}
\begin{ansblock}[2章-拓-课时16-08]
% ans:2章-拓-课时16-08
\ansitem{08}{\(BCD\)}
\ansnote{详解}{圆 \(C\) 圆心 \((0,1)\)、半径 \(4\)；抛物线 \(x^{2}=4y\) 的焦点恰为 \(F(0,1)\)、准线 \(y=-1\)，圆心与焦点重合。联立 \(x^{2}=4y\) 与 \(x^{2}+(y-1)^{2}=16\)：\(4y+(y-1)^{2}=16 \Rightarrow  y^{2}+2y-15=0 \Rightarrow  y=3\)（\(y=-5\) 舍，\(x^{2}=-20\) 无解），交点 \((\pm 2\sqrt{3},3)\)。劣弧 \(AB\) 为弦 \(y=3\) 上方含点 \((0,5)\) 的短弧（圆心到弦距离 \(2<\)半径 \(4\)），其上动点 \(P\) 纵坐标 \(y\_P\in (3,5)\)，\(A\) 项误以横坐标 \(2\sqrt{3}\) 为下界，错。\(B\)：\(P\)、\(N\) 同横坐标且 \(P\) 在 \(N\) 上方，\(|PN|=y\_P-y\_N\)；由定义 \(|NF|=y\_N+1\)（\(N\) 到准线距离），故 \(|PN|+|NF|=y\_P+1\)，恰为 \(P\) 到准线 \(y=-1\) 的距离，对（特位 \(x=0\)：\(P(0,4)\)、\(N(0,0)\)，\(4+1=5=y\_P+1\) ）。\(C\)：圆心 \((0,1)\) 到准线 \(y=-1\) 距离为 \(2\)，对。\(D\)：\(|PF|=\)半径 \(4\)，\(\triangle PFN\) 周长\(=|PF|+|PN|+|NF|=4+(y\_P-y\_N)+(y\_N+1)=y\_P+5\in (8,10)\)，对。故选 \(BCD\)。〔卷③自带详解照录校订；装配图位：源含简图〕}
\end{ansblock}

\zux{拓展·AP·BP最小与P坐标}
\tuhao{09} \tieside{难}已知点 \(A(2,0)\)，\(B(4,0)\)，动点 \(P\) 在抛物线 \(y^{2}=-4x\) 上运动，则使 \(AP\cdot BP\) 取得最小值的点 \(P\) 的坐标是\kongda{\((0,0)\)（最小值 \(8\)）}．
\begin{ansblock}[2章-拓-课时16-09]
% ans:2章-拓-课时16-09
\ansitem{09}{\((0,0)\)（最小值 \(8\)）}
\ansnote{详解}{设 \(P(x_{0},y_{0})\)，\(y_{0}^{2}=-4x_{0}\)（\(x_{0}\leqslant 0\)）。\(AP\cdot BP=(x_{0}-2)(x_{0}-4)+y_{0}^{2}=x_{0}^{2}-6x_{0}+8-4x_{0}=x_{0}^{2}-10x_{0}+8=(x_{0}-5)^{2}-17\)。由 \(x_{0}\leqslant 0\)，\((x_{0}-5)^{2}\geqslant 25\) 且随 \(x_{0}\) 增大而减小，故 \(x_{0}=0\) 时 \(AP\cdot BP\) 最小，最小值 \(8\)，此时 \(y_{0}=0\)，点 \(P(0,0)\)。〔卷③自带详解照录校订〕}
\end{ansblock}

\zux{拓展·外角平分线与线段等式}
% 图债注：「如图」指源件配图，印面无图（冻片【注】裁定：去装饰图/条件全在文字/尺寸随注——图债登记汇编报告）
\tuhao{10}\duoxuan\hspace{0.6em} \tieside{难}如图，在平面直角坐标系 \(xOy\) 中，抛物线 \(C\)：\(y^{2}=2px\)（\(p>0\)）的焦点为 \(F\)，准线为 \(l\)．设 \(l\) 与 \(x\) 轴的交点为 \(K\)，\(P\) 为抛物线 \(C\) 上异于 \(O\) 的任意一点，\(P\) 在 \(l\) 上的射影为 \(E\)，\(\angle EPF\) 的外角平分线交 \(x\) 轴于点 \(Q\)，过 \(Q\) 作 \(QM\perp PF\) 交 \(PF\) 于 \(M\)，过 \(Q\) 作 \(QN\perp EP\) 交线段 \(EP\) 的延长线于 \(N\)，则（　）
\bindopt {\fontsize{10.5pt}{\qpTJgLeadLiti}\selectfont \makebox[\dimexpr0.5\linewidth-0.5\lxhang-0.25em\relax][l]{\(A\)．\(|PE|=|PF|\)}\makebox[\dimexpr0.5\linewidth-0.5\lxhang-0.25em\relax][l]{\(B\)．\(|PF|=|QF|\)}\par}
\bindopt {\fontsize{10.5pt}{\qpTJgLeadLiti}\selectfont \makebox[\dimexpr0.5\linewidth-0.5\lxhang-0.25em\relax][l]{\(C\)．\(|PN|=|MF|\)}\makebox[\dimexpr0.5\linewidth-0.5\lxhang-0.25em\relax][l]{\(D\)．\(|PN|=|KF|\)}\par}
\begin{ansblock}[2章-拓-课时16-10]
% ans:2章-拓-课时16-10
\ansitem{10}{\(ABD\)}
\ansnote{详解}{\(A\)：由定义，\(|PF|\) 等于 \(P\) 到准线 \(l\) 的距离，而 \(PE\perp l\) 于 \(E\)，故 \(|PE|=|PF|\)，对。\(B\)：\(PQ\) 平分 \(\angle EPF\) 的外角 \(\Rightarrow  \angle FPQ=\angle NPQ\)；又 \(EN\parallel x\) 轴（同垂直于准线）\(\Rightarrow  \angle NPQ=\angle PQF\)（内错角），得 \(\angle FPQ=\angle PQF\)，\(\triangle FPQ\) 中 \(|PF|=|QF|\)，对。\(C\)：若 \(|PN|=|MF|\)，又 \(|QM|=|QN|\)（\(Q\) 在外角平分线上，到角两边等距）、\(\angle QMF=\angle QNP=90^{\circ}\)，则 \(\triangle FMQ\cong \triangle PNQ\)，得 \(|FQ|=|PQ|\)，结合 \(B\) 得 \(\triangle PFQ\) 为正三角形、\(\angle PFQ=\pi /3\)，则 \(P\) 被锁定为定点，与「\(P\) 为异于 \(O\) 的任意一点」矛盾，\(C\) 不恒成立。\(D\)：四边形 \(KQNE\) 为矩形（\(\angle KEQ=\angle EQN=\angle K=90^{\circ}\)）\(\Rightarrow  |EN|=|KQ|\)；\(|PN|=|EN|-|EP|=|KQ|-|PF|=|KF|+|FQ|-|PF|=|KF|\)（用 \(|EP|=|PF|=|FQ|\)），对。故选 \(ABD\)。〔卷③自带详解照录校订；装配图位：题面图嵌于源；\(D\) 的线段加减链已按源补全〕}
\end{ansblock}

\zux{拓展·焦点弦结论判定}
\tuhao{11}\duoxuan\hspace{0.6em} \tieside{难}已知抛物线 \(C\)：\(y^{2}=2px\)（\(p>0\)）的焦点为 \(F\)，直线 \(l\) 的斜率为 \(\sqrt{3}\) 且经过点 \(F\)，直线 \(l\) 与抛物线 \(C\) 交于点 \(A\)、\(B\) 两点（点 \(A\) 在第一象限），与抛物线的准线交于点 \(D\)，若 \(|AF|=8\)，则以下结论正确的是（　）
\bindopt {\fontsize{10.5pt}{\qpTJgLeadLiti}\selectfont \makebox[\dimexpr0.5\linewidth-0.5\lxhang-0.25em\relax][l]{\(A\)．\(p=4\)}\makebox[\dimexpr0.5\linewidth-0.5\lxhang-0.25em\relax][l]{\(B\)．\(DF=FA\)}\par}
\bindopt {\fontsize{10.5pt}{\qpTJgLeadLiti}\selectfont \makebox[\dimexpr0.5\linewidth-0.5\lxhang-0.25em\relax][l]{\(C\)．\(|BD|=2|BF|\)}\makebox[\dimexpr0.5\linewidth-0.5\lxhang-0.25em\relax][l]{\(D\)．\(|BF|=4\)}\par}
\begin{ansblock}[2章-拓-课时16-11]
% ans:2章-拓-课时16-11
\ansitem{11}{\(ABC\)}
\ansnote{详解}{准线 \(m\) 交 \(x\) 轴于 \(P\)，\(|PF|=p\)。分别过 \(A\)、\(B\) 作准线的垂线，垂足 \(E\)、\(M\)。倾斜角 \(60^{\circ}\)：\(AE\parallel x\) 轴 \(\Rightarrow  \angle EAF=60^{\circ}\)；由定义 \(|AE|=|AF| \Rightarrow  \triangle AEF\) 为等边三角形 \(\Rightarrow  \angle EFP=\angle AEF=60^{\circ}\)、\(\angle FEP=30^{\circ}\)（直角\(\triangle EPF\) 中），于是 \(|AF|=|EF|=2|PF|=2p=8\)，\(p=4\)，\(A\) 对。又 \(|AE|=|EF|=2|PF|\) 且 \(PF\parallel AE \Rightarrow  F\) 为 \(AD\) 中点，\(|DF|=|FA|\)，\(B\) 对。\(\angle DAE=60^{\circ}\) 的补算：\(\triangle ADE\) 中 \(\angle AED=90^{\circ}\)、\(\angle ADE=30^{\circ} \Rightarrow  |BD|=2|BM|=2|BF|\)（末步用定义 \(|BM|=|BF|\)），\(C\) 对。由 \(|BD|=2|BF|\) 得 \(|BF|=(1/3)|DF|=(1/3)|AF|=8/3\neq 4\)，\(D\) 错。故选 \(ABC\)。〔卷③自带详解照录校订；装配图位：源含简图〕}
\end{ansblock}

\zux{拓展·抛物线与圆四交点定值}
\tuhao{12} \tieside{难}已知抛物线 \(y^{2}=4x\) 与圆 \(F\)：\(x^{2}+y^{2}-2x=0\)，过点 \(F\) 作直线 \(l\)，自上而下顺次与上述两曲线交于点 \(A\)，\(B\)，\(C\)，\(D\)，则下列关于 \(|AB|\cdot |CD|\) 的值的说法中，正确的是（　）
\bindopt {\fontsize{10.5pt}{\qpTJgLeadLiti}\selectfont \makebox[\dimexpr0.5\linewidth-0.5\lxhang-0.25em\relax][l]{\(A\)．等于 \(1\)}\makebox[\dimexpr0.5\linewidth-0.5\lxhang-0.25em\relax][l]{\(B\)．等于 \(16\)}\par}
\bindopt {\fontsize{10.5pt}{\qpTJgLeadLiti}\selectfont \makebox[\dimexpr0.5\linewidth-0.5\lxhang-0.25em\relax][l]{\(C\)．最小值为 \(4\)}\makebox[\dimexpr0.5\linewidth-0.5\lxhang-0.25em\relax][l]{\(D\)．最大值为 \(4\)}\par}
\begin{ansblock}[2章-拓-课时16-12]
% ans:2章-拓-课时16-12
\ansitem{12}{\(A\)}
\ansnote{详解}{圆化标准形：\((x-1)^{2}+y^{2}=1\)，圆心 \(F(1,0)\)、半径 \(1\)，恰为抛物线 \(y^{2}=4x\) 的焦点。设 \(A(x_{1},y_{1})\)、\(D(x_{2},y_{2})\)（\(A\)、\(D\) 为直线与抛物线的两交点，\(B\)、\(C\) 为与圆的两交点，自上而下顺次），由焦半径 \(|AF|=x_{1}+1\)、\(|DF|=x_{2}+1\)，得 \(|AB|=|AF|-|BF|=x_{1}+1-1=x_{1}\)，\(|CD|=|DF|-|CF|=x_{2}\)。设 \(l\)：\(x=my+1\) 代入 \(y^{2}=4x\) 得 \(y^{2}-4my-4=0\)，\(y_{1}y_{2}=-4\)，故 \(|AB|\cdot |CD|=x_{1}x_{2}=(y_{1}^{2}/4)(y_{2}^{2}/4)=(y_{1}y_{2})^{2}/16=16/16=1\)，为定值 \(1\)（与 \(m\) 无关，最值之说随之落空）。选 \(A\)。〔卷③自带详解照录校订；源用「平均性质」名目，按 §五 口径已改写为常规联立韦达，答案不变〕}
\end{ansblock}

\zux{拓展·OA·OB定值}
\tuhao{13} \tieside{难}直线 \(y=kx+2\) 与抛物线 \(x^{2}=4y\) 交于 \(A\)，\(B\) 两点，则 \(OA\cdot OB\)（\(O\) 为抛物线顶点）的值为（　）
\bindopt {\fontsize{10.5pt}{\qpTJgLeadLiti}\selectfont \makebox[\dimexpr0.5\linewidth-0.5\lxhang-0.25em\relax][l]{\(A\)．【图嵌未回】}\makebox[\dimexpr0.5\linewidth-0.5\lxhang-0.25em\relax][l]{\(B\)．【图嵌未回】}\par}
\bindopt {\fontsize{10.5pt}{\qpTJgLeadLiti}\selectfont \makebox[\dimexpr0.5\linewidth-0.5\lxhang-0.25em\relax][l]{\(C\)．\(4\)}\makebox[\dimexpr0.5\linewidth-0.5\lxhang-0.25em\relax][l]{\(D\)．\(12\)}\par}
\begin{ansblock}[2章-拓-课时16-13]
% ans:2章-拓-课时16-13
\ansitem{13}{\(B\)（\(-4\)）}
\ansnote{详解}{联立 \(x^{2}=4(kx+2)\)：\(x^{2}-4kx-8=0\)（\(\Delta =16k^{2}+32>0\) 恒成立），\(x_{1}+x_{2}=4k\)、\(x_{1}x_{2}=-8\)。\(y_{1}y_{2}=(kx_{1}+2)(kx_{2}+2)=k^{2}x_{1}x_{2}+2k(x_{1}+x_{2})+4=-8k^{2}+8k^{2}+4=4\)（即过 \((0,2)\) 的弦两端纵坐标之积为定值 \(4\)，\(A\)、\(B\) 分居 \(y\) 轴两侧故 \(x_{1}x_{2}<0\) 与 \(-8\) 相符）。\(OA\cdot OB=x_{1}x_{2}+y_{1}y_{2}=-8+4=-4\)，与 \(k\) 无关。值 \(-4\) 对应选项 \(B\)，选 \(B\)。〔校订注：卷③此题面之直线、抛物线方程及选项 \(A\)、\(B\) 数值在原 \(dump\) 中均为图嵌（【图】占位，图内文本未入对象层），题方按源【编注】【分析】【详解】读数 \(y=kx+2\)、\(x^{2}=4y\) 与 \(OA\cdot OB=-4\) 补全（联立恒等于源算路复核 ，见本块推导）；选项 \(A\) 数值未能回，装配回源图或改选项；「平均性质」名目已按 §五 改常规联立〕}
\end{ansblock}

\zux{拓展·焦点弦切线共线判定}
\tuhao{14}\duoxuan\hspace{0.6em} \tieside{难}在平面直角坐标系 \(xOy\) 中，过抛物线 \(x^{2}=2y\) 的焦点的直线 \(l\) 与该抛物线的两个交点为 \(A(x_{1},y_{1})\)，\(B(x_{2},y_{2})\)，则（　）
\lxopt{\(A\)．抛物线在点 \(x=1\) 处切线方程为 \(2x-2y-1=0\)}
\lxopt{\(B\)．若点 \(M\) 坐标为 \((0,-1/2)\)，则 \(AM\cdot BM=0\)}
\lxopt{\(C\)．\(|OA|+|OB|>\sqrt{5}\)}
\lxopt{\(D\)．若 \(BN\) 垂直抛物线准线于点 \(N\)，则 \(A\)，\(O\)，\(N\) 三点在一条直线上}
\begin{ansblock}[2章-拓-课时16-14]
% ans:2章-拓-课时16-14
\ansitem{14}{\(AD\)}
\ansnote{详解}{抛物线 \(x^{2}=2y\) 的焦点 \((0,1/2)\)、准线 \(y=-1/2\)。设 \(l\)：\(y=kx+1/2\)，代入 \(x^{2}=2y\) 得 \(x^{2}-2kx-1=0\)，\(x_{1}+x_{2}=2k\)、\(x_{1}x_{2}=-1\)；\(y_{1}+y_{2}=k(x_{1}+x_{2})+1=2k^{2}+1\)、\(y_{1}y_{2}=(kx_{1}+1/2)(kx_{2}+1/2)=k^{2}x_{1}x_{2}+(k/2)(x_{1}+x_{2})+1/4=-k^{2}+k^{2}+1/4=1/4\)。\(A\)：过点 \((1,1/2)\)、斜率 \(t\) 的直线 \(y=t(x-1)+1/2\) 代入 \(x^{2}=2y\) 得 \(x^{2}-2tx+2t-1=0\)，即 \((x-1)(x-(2t-1))=0\)，另一交点横标 \(2t-1\)；相切 \(\Leftrightarrow  2t-1=1 \Leftrightarrow  t=1\)，切线 \(y=x-1/2\) 即 \(2x-2y-1=0\)，\(A\) 对。\(B\)：\(AM\cdot BM=(-x_{1},-1/2-y_{1})\cdot (-x_{2},-1/2-y_{2})=x_{1}x_{2}+(1/2+y_{1})(1/2+y_{2})=x_{1}x_{2}+(1/4)+(1/2)(y_{1}+y_{2})+y_{1}y_{2}=-1+1/4+(1/2)(2k^{2}+1)+1/4=k^{2}\)，仅 \(k=0\)（通径位）时为 \(0\)，不恒成立，\(B\) 错。\(C\)：取 \(l\parallel x\) 轴（\(k=0\)，即通径），\(A(1,1/2)\)、\(B(-1,1/2)\)，\(|OA|=|OB|=\sqrt{1+1/4}=\sqrt{5}/2\)，\(|OA|+|OB|=\sqrt{5}\)，严格不等式不成立，\(C\) 错。\(D\)：\(BN\perp \)准线 \(\Rightarrow  N(x_{2},-1/2)\)；直线 \(OA\)：\(y=(y_{1}/x_{1})x=(x_{1}/2)x\)（用 \(y_{1}=x_{1}^{2}/2\)），代入 \(x=x_{2}\) 得纵标 \((x_{1}x_{2})/2=-1/2\)，即直线 \(OA\) 恰过 \(N(x_{2},-1/2)\)，\(A\)、\(O\)、\(N\) 共线，\(D\) 对。故选 \(AD\)。〔卷③自带详解照录校订；\(A\) 项源解用导数 \(y'=x\)（系选修\(2\) 工具），已消前引改「联立重根」路径，结论一致；\(C\) 项源 \(dump\) 排印「\(B(1,-1/2)\)」系讹（该点不在抛物线上），依 \(k=0\) 弦 \(x^{2}=1\) 校正为 \(B(-1,1/2)\)〕}
\end{ansblock}

\zux{拓展·垂线段中点轨迹}
\tuhao{15} \tieside{难}从抛物线 \(y^{2}=2px\)（\(p>0\)）上各点向 \(x\) 轴作垂线段，求垂线段中点的轨迹方程，并说明轨迹是什么曲线．
\liubai[16mm]{解答书写区}
\begin{ansblock}[2章-拓-课时16-15]
% ans:2章-拓-课时16-15
\ansitem{15}{\(y^{2}=(p/2)x\)（\(x\geqslant 0\)），轨迹是顶点在原点、对称轴为 \(x\) 轴、开口向右的抛物线（\(2\tilde p \allowbreak =p/2\)，焦点 \((p/8,0)\)）}
\ansnote{详解}{设抛物线上点 \((x_{0},y_{0})\)，\(y_{0}^{2}=2px_{0}\)（\(x_{0}\geqslant 0\)），垂线段中点 \((x,y)=(x_{0}, y_{0}/2)\)。消参：\(x_{0}=x\)、\(y_{0}=2y\) 代入得 \(4y^{2}=2px\)，即 \(y^{2}=(p/2)x\)（\(x\geqslant 0\)）。这是顶点在原点、对称轴为 \(x\) 轴、开口向右的抛物线，其 \(2\tilde p \allowbreak =p/2\)（\(\tilde p \allowbreak =p/4\)），焦点 \((p/8,0)\)。（原点处垂线段退化为点 \(O\)，中点仍为 \(O\)，轨迹含顶点 。）〔亲算印证：消参直推，读数自验〕}
\end{ansblock}

\zux{拓展·数量积定方程}
\tuhao{16} \tieside{难}已知抛物线的顶点是坐标原点 \(O\)，对称轴为 \(x\) 轴，焦点为 \(F\)，抛物线上的点 \(A\) 的横坐标为 \(2\)，且 \(FA\cdot OA=16\)，求此抛物线的方程．
\liubai[16mm]{解答书写区}
\begin{ansblock}[2章-拓-课时16-16]
% ans:2章-拓-课时16-16
\ansitem{16}{\(y^{2}=8x\)}
\ansnote{详解}{点 \(A\) 横坐标 \(2>0\) 在 \(y\) 轴右侧，故抛物线开口向右，设 \(y^{2}=2px\)（\(p>0\)），\(F(p/2,0)\)。设 \(A(2,y\_A)\)，\(y\_A^{2}=4p\)。\(FA\rightarrow =(2-p/2, y\_A)\)、\(OA\rightarrow =(2, y\_A)\)，\(FA\cdot OA=2(2-p/2)+y\_A^{2}=4-p+4p=4+3p=16\)，解得 \(p=4\)，抛物线为 \(y^{2}=8x\)。（回验：\(A(2,\pm 4)\)，\(FA\cdot OA=2(2-2)+16=16\) 。）〔亲算印证：自验读数 \(y^{2}=8x\)；教材答案条乱码段（\(raw\) 行\(4554\)起）未取，与 §三\(.1\) 口径一致〕}
\end{ansblock}

\zux{拓展·焦半径公式证明与总结}
\tuhao{17} \tieside{难}设抛物线 \(y^{2}=2px\)（\(p>0\)）上一点 \(M\) 的横坐标为 \(x_{0}\)，证明 \(M\) 到抛物线焦点的距离为 \(x_{0}+p/2\)，并总结出关于抛物线其他形式的标准方程的类似结论．
\liubai[16mm]{解答书写区}
\begin{ansblock}[2章-拓-课时16-17]
% ans:2章-拓-课时16-17
\ansitem{17}{证明见详解；类似结论：\(y^{2}=-2px\) 上 \(|MF|=p/2-x_{0}\)（\(x_{0}\leqslant 0\)）；\(x^{2}=2py\) 上 \(|MF|=y_{0}+p/2\)；\(x^{2}=-2py\) 上 \(|MF|=p/2-y_{0}\)}
\ansnote{详解}{证明：准线 \(l\)：\(x=-p/2\)。由抛物线定义，\(|MF|\) 等于 \(M\) 到 \(l\) 的距离；\(M\) 横坐标 \(x_{0}\geqslant 0\)，故距离为 \(x_{0}-(-p/2)=x_{0}+p/2\)。\(\rule{1.05ex}{1.05ex}\) 类似结论（同法：焦半径\(=\)到准线距离，注意坐标符号范围）：①\(y^{2}=-2px\)（\(x_{0}\leqslant 0\)，准线 \(x=p/2\)）：\(|MF|=p/2-x_{0}\)；②\(x^{2}=2py\)（\(y_{0}\geqslant 0\)，准线 \(y=-p/2\)）：\(|MF|=y_{0}+p/2\)；③\(x^{2}=-2py\)（\(y_{0}\leqslant 0\)，准线 \(y=p/2\)）：\(|MF|=p/2-y_{0}\)。统一记法：焦半径\(=p/2+\)（沿开口方向的点到顶点有向距离）。〔亲算印证：定义直推，与 §三\(.1\)「\(|MF|=|y_{0}|+p/2\)」型读数同族〕}
\end{ansblock}

\zux{拓展·FQ垂直平分线过哪点}
\tuhao{19} \tieside{难}设抛物线的顶点为 \(O\)，焦点为 \(F\)，准线为 \(l\)，\(P\) 是抛物线上异于 \(O\) 的一点，过 \(P\) 作 \(PQ\perp l\) 于 \(Q\)．则线段 \(FQ\) 的垂直平分线（　）
\bindopt {\fontsize{10.5pt}{\qpTJgLeadLiti}\selectfont \makebox[\dimexpr0.5\linewidth-0.5\lxhang-0.25em\relax][l]{\(A\)．经过点 \(O\)}\makebox[\dimexpr0.5\linewidth-0.5\lxhang-0.25em\relax][l]{\(B\)．经过点 \(P\)}\par}
\bindopt {\fontsize{10.5pt}{\qpTJgLeadLiti}\selectfont \makebox[\dimexpr0.5\linewidth-0.5\lxhang-0.25em\relax][l]{\(C\)．平行于直线 \(OP\)}\makebox[\dimexpr0.5\linewidth-0.5\lxhang-0.25em\relax][l]{\(D\)．垂直于直线 \(OP\)}\par}
\begin{ansblock}[2章-拓-课时16-19]
% ans:2章-拓-课时16-19
\ansitem{19}{\(B\)}
\ansnote{详解}{连接 \(PF\)。由抛物线定义，\(P\) 到焦点的距离等于 \(P\) 到准线的距离，即 \(|PF|=|PQ|\)，\(\triangle QPF\) 为等腰三角形，\(P\) 到线段 \(FQ\) 两端点等距，故 \(P\) 在线段 \(FQ\) 的垂直平分线上，选 \(B\)。（反例排 \(A\)：取 \(y^{2}=4x\)、\(P(1,2)\)，则 \(F(1,0)\)、\(Q(-1,2)\)，\(|OF|=1\) 而 \(|OQ|=\sqrt{5}\)，\(O\) 不在垂直平分线上；\(C\)、\(D\) 随之不恒成立。）〔亲算印证：定义直推，与轮\(2\)亲算④抽样档「过 \(P\)，选\(B\)」一致〕}
\end{ansblock}

\zux{拓展·到椭圆左顶点距离最小}
\tuhao{20} \tieside{难}已知点 \(P\) 在抛物线 \(y^{2}=-4x\) 上，求点 \(P\) 到椭圆 \(x^{2}/16+y^{2}/15=1\) 左顶点的距离最小值．
\liubai[16mm]{解答书写区}
\begin{ansblock}[2章-拓-课时16-20]
% ans:2章-拓-课时16-20
\ansitem{20}{\(2\sqrt{3}\)}
\ansnote{详解}{椭圆左顶点 \((-4,0)\)。设 \(P(-y^{2}/4,y)\)，则 \(d^{2}=(-y^{2}/4+4)^{2}+y^{2}=(1/16)y^{4}-y^{2}+16=(1/16)(y^{2}-8)^{2}+12\geqslant 12\)，当 \(y^{2}=8\)（\(y=\pm 2\sqrt{2}\)，\(x=-2\)，点 \(P(-2,\pm 2\sqrt{2})\) 在抛物线上 ）时取等，\(d\) 的最小值为 \(\sqrt{12}=2\sqrt{3}\)。〔亲算印证：轮\(2\)亲算④抽样档 \(2\sqrt{3}\) 与 \(3\)章件\(14-\#10\) 源答案一致〕}
\end{ansblock}

\zux{拓展·|PA|+|PF|最小与P坐标}
\tuhao{21} \tieside{难}已知抛物线 \(y^{2}=2x\) 的焦点是 \(F\)，点 \(P\) 是抛物线上的动点，若 \(A(3,2)\)，则 \(|PA|+|PF|\) 的最小值为\kongda{\(7/2\)；\(P(2,2)\)}，此时点 \(P\) 的坐标为\(\kongbai{}\)．
\begin{ansblock}[2章-拓-课时16-21]
% ans:2章-拓-课时16-21
\ansitem{21}{\(7/2\)；\(P(2,2)\)}
\ansnote{详解}{\(A(3,2)\)：\(2^{2}=4<2\cdot 3=6\)，\(A\) 在抛物线内部。准线 \(l\)：\(x=-1/2\)。设 \(P\) 在 \(l\) 上的射影为 \(Q\)，由定义 \(|PF|=|PQ|\)，\(|PA|+|PF|=|PA|+|PQ|\geqslant A\) 到准线的水平距离\(=3-(-1/2)=7/2\)，当 \(A\)、\(P\)、\(Q\) 共线（\(y\_P=y\_A=2\)）时取等。代入 \(y^{2}=2x\) 得 \(x=2\)，即 \(P(2,2)\)，最小值 \(7/2\)。〔亲算印证：轮\(2\)亲算④抽样档「\(7/2\)、\(P(2,2)\)」与源答案一致〕}
\end{ansblock}

\zux{拓展·焦半径求PF长}
\tuhao{22} \tieside{难}已知抛物线 \(y^{2}=4x\) 的焦点为 \(F\)，点 \(P(m,-4)\) 在抛物线上，则 \(PF\) 的长为（　）
\bindopt {\fontsize{10.5pt}{\qpTJgLeadLiti}\selectfont \makebox[\dimexpr0.25\linewidth-0.25\lxhang-0.25em\relax][l]{\(A\)．\(2\)}\makebox[\dimexpr0.25\linewidth-0.25\lxhang-0.25em\relax][l]{\(B\)．\(3\)}\makebox[\dimexpr0.25\linewidth-0.25\lxhang-0.25em\relax][l]{\(C\)．\(4\)}\makebox[\dimexpr0.25\linewidth-0.25\lxhang-0.25em\relax][l]{\(D\)．\(5\)}\par}
\begin{ansblock}[2章-拓-课时16-22]
% ans:2章-拓-课时16-22
\ansitem{22}{\(D\)}
\ansnote{详解}{\(P\) 在抛物线上：\(16=4m\)、\(m=4\)，\(P(4,-4)\)。焦半径（定义）：准线 \(x=-1\)，\(|PF|=x\_P+1=4+1=5\)。（两点距复核：\(F(1,0)\)，\(|PF|=\sqrt{9+16}=5\) 。）选 \(D\)。〔亲算印证：轮\(2\)亲算④抽样档「\(5\)，选\(D\)」与源答案一致〕}
\end{ansblock}

\zux{拓展·动圆相切圆心轨迹}
\tuhao{23} \tieside{难}已知直线 \(l\)：\(x=-2\) 和圆 \(C\)：\((x-3)^{2}+y^{2}=1\)．若圆 \(M\) 与直线 \(l\) 相切，与圆 \(C\) 外切，求圆 \(M\) 的圆心 \(M\) 的轨迹方程．
\liubai[16mm]{解答书写区}
\begin{ansblock}[2章-拓-课时16-23]
% ans:2章-拓-课时16-23
\ansitem{23}{\(y^{2}=12x\)}
\ansnote{详解}{设圆 \(M\) 半径为 \(R\)。\(\odot C\) 圆心 \(C(3,0)\)、半径 \(1\)，外切 \(\Rightarrow  |MC|=R+1\)；\(\odot M\) 与直线 \(x=-2\) 相切 \(\Rightarrow \) 圆心 \(M\) 到 \(x=-2\) 的距离为 \(R\)，从而 \(M\) 到直线 \(x=-3\) 的距离为 \(R+1\)。于是动点 \(M\) 到定点 \(C(3,0)\) 的距离等于它到定直线 \(x=-3\) 的距离，轨迹是以 \(C\) 为焦点、\(x=-3\) 为准线的抛物线（定义）：\(p/2=3\)、\(p=6\)，标准方程 \(y^{2}=12x\)。（核验 \(x\geqslant 0\)：\(M\) 到 \(x=-3\) 距离\(=\)横坐标\(+3=R+1>0\) 且 \(|MC|=R+1\geqslant 1>0\) 自洽。）〔亲算印证：轮\(2\)亲算④抽样档 \(y^{2}=12x\) 与源答案一致〕}
\end{ansblock}

\zux{拓展·等腰三角形顶角}
\tuhao{24} \tieside{难}\(A\)，\(B\) 是抛物线 \(x^{2}=2y\) 上的两点，\(O\) 为坐标原点．若 \(|OA|=|OB|\)，且 \(\triangle AOB\) 的面积为 \(12\sqrt{3}\)，则 \(\angle AOB=\)（　）
\bindopt {\fontsize{10.5pt}{\qpTJgLeadLiti}\selectfont \makebox[\dimexpr0.25\linewidth-0.25\lxhang-0.25em\relax][l]{\(A\)．\(30^{\circ}\)}\makebox[\dimexpr0.25\linewidth-0.25\lxhang-0.25em\relax][l]{\(B\)．\(45^{\circ}\)}\makebox[\dimexpr0.25\linewidth-0.25\lxhang-0.25em\relax][l]{\(C\)．\(60^{\circ}\)}\makebox[\dimexpr0.25\linewidth-0.25\lxhang-0.25em\relax][l]{\(D\)．\(120^{\circ}\)}\par}
\begin{ansblock}[2章-拓-课时16-24]
% ans:2章-拓-课时16-24
\ansitem{24}{\(C\)}
\ansnote{详解}{\(|OA|=|OB|\)：\(x_{1}^{2}+y_{1}^{2}=x_{2}^{2}+y_{2}^{2}\)，而 \(y\_i=x\_i^{2}/2\)，\((x_{1}^{2}-x_{2}^{2})(1+(x_{1}^{2}+x_{2}^{2})/4)=0 \Rightarrow  x_{1}^{2}=x_{2}^{2}\)，\(A\)、\(B\) 异点故 \(x_{2}=-x_{1}\)，两点关于 \(y\) 轴对称。设 \(B(a,a^{2}/2)\)、\(A(-a,a^{2}/2)\)（\(a>0\)），\(S=\)½\(\cdot 2a\cdot (a^{2}/2)=a^{3}/2=12\sqrt{3} \Rightarrow  a^{3}=24\sqrt{3}=(2\sqrt{3})^{3} \Rightarrow  a=2\sqrt{3}\)，\(B(2\sqrt{3},6)\)。记 \(OB\) 与 \(y\) 轴夹角 \(\theta \)：\(tan\theta =a/(a^{2}/2)=2\sqrt{3}/6=\sqrt{3}/3 \Rightarrow  \theta =30^{\circ}\)，\(\angle AOB=2\theta =60^{\circ}\)。选 \(C\)。〔亲算印证：轮\(2\)亲算④抽样档「\(60^{\circ}\)，选\(C\)」与源答案一致〕}
\end{ansblock}

\zux{拓展·等腰梯形内接求焦半径}
\tuhao{25} \tieside{难}已知等腰梯形 \(ABCD\) 的顶点都在抛物线 \(y^{2}=2px\)（\(p>0\)）上，且 \(AB\parallel CD\)，\(|AB|=2\)，\(|CD|=4\)，\(\angle ADC=60^{\circ}\)，则点 \(A\) 到抛物线的焦点的距离是\kongda{\(7\sqrt{3}/12\)}．
\begin{ansblock}[2章-拓-课时16-25]
% ans:2章-拓-课时16-25
\ansitem{25}{\(7\sqrt{3}/12\)}
\ansnote{详解}{等腰梯形两底 \(AB\parallel CD\)，而抛物线上平行于 \(x\) 轴方向的弦（两交点同纵坐标）不存在（\(y^{2}=2px\) 单值），故两底必为垂直于 \(x\) 轴的弦：\(A\)、\(B\) 关于 \(x\) 轴对称，\(C\)、\(D\) 关于 \(x\) 轴对称。\(|AB|=2 \Rightarrow  y\_A=1\)（取上半），\(|CD|=4 \Rightarrow  y\_D=2\)。设 \(A(m,1)\)；腰与底夹角 \(\angle ADC=60^{\circ}\)，底铅直，故 \(DA\) 与铅直线夹 \(60^{\circ}\)，纵差 \(2-1=1 \Rightarrow \) 横差\(=tan60^{\circ}=\sqrt{3}\)，\(D(m+\sqrt{3},2)\)。两点在抛物线上：\(1=2pm\)、\(4=2p(m+\sqrt{3})\)，两式相减 \(3=2\sqrt{3} p \Rightarrow  p=\sqrt{3}/2\)、\(m=1/(2p)=\sqrt{3}/3\)。由焦半径（拓\(16-17\) 结论）\(|AF|=m+p/2=\sqrt{3}/3+\sqrt{3}/4=7\sqrt{3}/12\)。〔亲算印证：轮\(2\)亲算②\(-1\) 关键步「\(p=\sqrt{3}/2\)、\(m=\sqrt{3}/3\)、\(|AF|=7\sqrt{3}/12\)」与 \(3\)章件\(15-\#14\) 源答案一致〕}
\end{ansblock}

\zux{拓展·锐角存在性求m范围}
\tuhao{26} \tieside{难}已知 \(F\) 为抛物线 \(y^{2}=4x\) 的焦点，过 \(F\) 的直线 \(l\) 与抛物线交于 \(A\)，\(B\) 两点，若在 \(x\) 轴负半轴上存在一点 \(M(m,0)\)，使得 \(\angle AMB\) 为锐角，则 \(m\) 的取值范围为（　）
\bindopt {\fontsize{10.5pt}{\qpTJgLeadLiti}\selectfont \makebox[\dimexpr0.5\linewidth-0.5\lxhang-0.25em\relax][l]{\(A\)．\((-\infty ,-1)\)}\makebox[\dimexpr0.5\linewidth-0.5\lxhang-0.25em\relax][l]{\(B\)．\((-\infty ,-1]\)}\par}
\bindopt {\fontsize{10.5pt}{\qpTJgLeadLiti}\selectfont \makebox[\dimexpr0.5\linewidth-0.5\lxhang-0.25em\relax][l]{\(C\)．\((-\infty ,-3)\)}\makebox[\dimexpr0.5\linewidth-0.5\lxhang-0.25em\relax][l]{\(D\)．\((-\infty ,-3]\)}\par}
\begin{ansblock}[2章-拓-课时16-26]
% ans:2章-拓-课时16-26
\ansitem{26}{\(A\)}
\ansnote{详解}{\(F(1,0)\)。设 \(l\)：\(x=ty+1\)（含斜率不存在），代入 \(y^{2}=4x\)：\(y^{2}-4ty-4=0\)，\(y_{1}+y_{2}=4t\)、\(y_{1}y_{2}=-4\)，\(x_{1}+x_{2}=t(y_{1}+y_{2})+2=4t^{2}+2\)、\(x_{1}x_{2}=(y_{1}y_{2})^{2}/16=1\)。\(\angle AMB\) 为锐角 \(\Leftrightarrow  MA\cdot MB>0\)（且 \(A\)、\(M\)、\(B\) 不共线，\(M\) 在 \(x\) 轴负半轴而弦过 \(F\) 恒成立）：\(MA\cdot MB=(x_{1}-m)(x_{2}-m)+y_{1}y_{2}=x_{1}x_{2}-m(x_{1}+x_{2})+m^{2}-4=m^{2}-2m-3-4mt^{2}\)。题源解答按「对一切过 \(F\) 的弦（\(\forall t\in R\)）成立」求解：\(m<0\) 时 \(-4mt^{2}\geqslant 0\) 且随 \(t^{2}\) 递增，最严处 \(t=0\)，须 \(m^{2}-2m-3>0 \Rightarrow  m<-1\) 或 \(m>3\)，配 \(m<0\) 得 \(m\in (-\infty ,-1)\)，选 \(A\)。〔措辞存疑：题面「存在 \(M\) 使 \(\angle AMB\) 为锐角」若按「对某条弦成立」解读，则一切 \(m<0\) 皆可（\(t\) 足够大即成立），与选项无合——故题面取「对任意焦点弦恒可为锐角」义，与源答案 \(A\) 一致，已登本件存疑清单〕〔亲算印证：轮\(2\)亲算②\(-1\)「\((-\infty ,-1)\)，选\(A\)」与 \(3\)章件\(18-\#1\) 源答案一致〕}
\end{ansblock}

\zux{拓展·弦中点到准线距离最小}
\tuhao{27} \tieside{难}已知抛物线 \(C\)：\(y^{2}=8x\) 的焦点为 \(F\)，过点 \(F\) 的直线交 \(C\) 于 \(A\)，\(B\) 两点，则 \(AB\) 的中点 \(M\) 到 \(C\) 的准线 \(l\) 的距离的最小值为（　）
\bindopt {\fontsize{10.5pt}{\qpTJgLeadLiti}\selectfont \makebox[\dimexpr0.25\linewidth-0.25\lxhang-0.25em\relax][l]{\(A\)．\(2\)}\makebox[\dimexpr0.25\linewidth-0.25\lxhang-0.25em\relax][l]{\(B\)．\(4\)}\makebox[\dimexpr0.25\linewidth-0.25\lxhang-0.25em\relax][l]{\(C\)．\(5\)}\makebox[\dimexpr0.25\linewidth-0.25\lxhang-0.25em\relax][l]{\(D\)．\(6\)}\par}
\begin{ansblock}[2章-拓-课时16-27]
% ans:2章-拓-课时16-27
\ansitem{27}{\(B\)}
\ansnote{详解}{\(p=4\)，准线 \(x=-2\)。分别过 \(A\)、\(M\)、\(B\) 作准线的垂线，垂足 \(C\)、\(D\)、\(E\)，则 \(MD\) 为梯形（退化情形为中点公式仍合）\(ACEB\) 的中位线，\(|MD|=(|AC|+|BE|)/2=(|AF|+|BF|)/2=|AB|/2\)（定义）。设 \(l\)：\(x=my+2\) 代入 \(y^{2}=8x\)：\(y^{2}-8my-16=0\)，\(|AB|=\sqrt{1+m^{2}}|y_{1}-y_{2}|=\sqrt{1+m^{2}}\cdot \sqrt{64m^{2}+64}=8(1+m^{2})\geqslant 8\)（\(m=0\) 即通径取等）。故 \(|MD|=|AB|/2\geqslant 4\)，最小值 \(4\)，选 \(B\)。〔亲算印证：轮\(2\)亲算②\(-1\)「\(\geqslant \)½通径\(=4\)，选\(B\)」与 \(3\)章件\(18-\#2\) 源答案一致〕}
\end{ansblock}

\zux{拓展·距离和最小与夹角余弦}
\tuhao{28} \tieside{难}已知抛物线 \(C\)：\(y^{2}=4x\) 的焦点为 \(F\)，准线为 \(l\)，直线 \(l'\)：\(x-y+2=0\)，动点 \(M\) 在 \(C\) 上运动，记点 \(M\) 到直线 \(l\) 与 \(l'\) 的距离分别为 \(d_{1}\)，\(d_{2}\)，\(O\) 为坐标原点，则当 \(d_{1}+d_{2}\) 最小时，\(cos\angle MFO=\)（　）
\bindopt {\fontsize{10.5pt}{\qpTJgLeadLiti}\selectfont \makebox[\dimexpr0.25\linewidth-0.25\lxhang-0.25em\relax][l]{\(A\)．\(\sqrt{2}/2\)}\makebox[\dimexpr0.25\linewidth-0.25\lxhang-0.25em\relax][l]{\(B\)．\(\sqrt{2}/3\)}\makebox[\dimexpr0.25\linewidth-0.25\lxhang-0.25em\relax][l]{\(C\)．\(\sqrt{2}/4\)}\makebox[\dimexpr0.25\linewidth-0.25\lxhang-0.25em\relax][l]{\(D\)．\(\sqrt{2}/6\)}\par}
\begin{ansblock}[2章-拓-课时16-28]
% ans:2章-拓-课时16-28
\ansitem{28}{\(A\)}
\ansnote{详解}{由定义 \(d_{1}=|MF|\)；设 \(MN\perp l'\) 于 \(N\)，则 \(d_{2}=|MN|\)，\(d_{1}+d_{2}=|MF|+|MN|\geqslant d(F,l')\)（\(F\) 到 \(l'\) 的垂线段长），当 \(M\) 落在「过 \(F\) 垂直于 \(l'\) 的直线」与抛物线的交点上（该垂线与 \(l'\) 的垂足 \(N\) 在 \(F\) 同侧射影之间，可行）时取最小。\(d(F,l')=|1-0+2|/\sqrt{2}=3\sqrt{2}/2\)。设 \(l'\) 与 \(x\) 轴交于 \(D(-2,0)\)，\(|FD|=3\)；在 \(Rt\triangle DNF\)（\(\angle DNF=90^{\circ}\)）中 \(cos\angle NFD=|FN|/|FD|=(3\sqrt{2}/2)/3=\sqrt{2}/2\)；取最小值时 \(M\) 在线段 \(FN\) 上、射线 \(FO\) 与 \(FD\) 同向，\(cos\angle MFO=\sqrt{2}/2\)，选 \(A\)。（\(M\) 坐标核：\(M(t^{2},2t)\) 且 \(FM\parallel (1,-1)\) 法向 \(\rightarrow  t=\sqrt{2}-1\)，\(M(3-2\sqrt{2},2\sqrt{2}-2)\) 在 \(FN\) 段内 。）〔亲算印证：轮\(2\)亲算②\(-1\)「\(cos\angle MFO=\sqrt{2}/2\)，选\(A\)」与 \(3\)章件\(18-\#7\) 源答案一致；装配图位：源解含图〕}
\end{ansblock}

\zux{拓展·动圆轨迹与恒过定点}
\tuhao{29} \tieside{难}已知动圆 \(M\) 过点 \((2,0)\)，被 \(y\) 轴截得的弦长为 \(4\)．
\((1)\) 求圆心 \(M\) 的轨迹方程；
\((2)\) 若 \(\triangle ABC\) 的顶点在 \(M\) 的轨迹上，且点 \(A\)、\(C\) 关于 \(x\) 轴对称，直线 \(BC\) 经过点 \(F(1,0)\)，求证：直线 \(AB\) 恒过定点．
\liubai[24mm]{解答书写区}
\begin{ansblock}[2章-拓-课时16-29]
% ans:2章-拓-课时16-29
\ansitem{29}{\((1) y^{2}=4x\)；\((2)\) 直线 \(AB\) 恒过定点 \((-1,0)\)}
\ansnote{详解}{\((1)\) 设圆心 \(M(x,y)\)，圆过 \((2,0) \Rightarrow \) 半径\(^{2}=(x-2)^{2}+y^{2}\)；被 \(y\) 轴截弦长 \(4 \Rightarrow \) 半径\(^{2}=x^{2}+2^{2}\)（半弦 \(2\)、弦心距 \(|x|\)）。两式相等：\((x-2)^{2}+y^{2}=x^{2}+4 \Rightarrow  y^{2}=4x\)。 \((2)\) 顶点在轨迹 \(y^{2}=4x\) 上。设 \(B(x_{1},y_{1})\)、\(C(x_{2},y_{2})\)，\(BC\) 过 \(F(1,0)\) 且不垂直于坐标轴，设 \(BC\)：\(x=ty+1\)：\(y^{2}-4ty-4=0 \Rightarrow  y_{1}y_{2}=-4\)。由 \(A\)、\(C\) 关于 \(x\) 轴对称，\(A(x_{2},-y_{2})\)。设 \(AB\)：\(y=kx+m\)（\(k\neq 0\)）代入 \(ky^{2}-4y+4m=0\)，其两根为 \(y_{1}\)、\(-y_{2}\)：\(y_{1}\cdot (-y_{2})=4m/k\)，即 \(4=4m/k\)、\(m=k\)。故 \(AB\)：\(y=k(x+1)\)，对一切实有 \(k\) 恒过定点 \((-1,0)\)。\(\rule{1.05ex}{1.05ex}\)（源答案栏「证明见解析」未落定点数值，定点 \((-1,0)\) 系亲算补出，登轮\(2\)亲算⑤\(-5\)。）〔亲算印证：轮\(2\)亲算②\(-1\) 参数点法同得 \((-1,0)\)；本块 \((2)\) 采源韦达路径改写（消「见解析」）〕}
\end{ansblock}

\zux{拓展·切线交点定直线}
\tuhao{30} \tieside{难}已知抛物线 \(C\)：\(x^{2}=2py\)（\(p>0\)）与圆 \(O\)：\(x^{2}+y^{2}=12\) 相交于 \(A\)，\(B\) 两点，且点 \(A\) 的横坐标为 \(2\sqrt{2}\)．\(F\) 是抛物线 \(C\) 的焦点，过焦点的直线 \(l\) 与抛物线 \(C\) 相交于不同的两点 \(M\)，\(N\)．
\((1)\) 求抛物线 \(C\) 的方程；
\((2)\) 过点 \(M\)，\(N\) 作抛物线 \(C\) 的切线 \(l_{1}\)，\(l_{2}\)，\(P(x_{0},y_{0})\) 是 \(l_{1}\)，\(l_{2}\) 的交点，求证：点 \(P\) 在定直线上．
\liubai[24mm]{解答书写区}
\begin{ansblock}[2章-拓-课时16-30]
% ans:2章-拓-课时16-30
\ansitem{30}{\((1) x^{2}=4y\)；\((2)\) 点 \(P\) 恒在定直线 \(y=-1\)（即 \(C\) 的准线）上}
\ansnote{详解}{\((1) A\) 横坐标 \(2\sqrt{2}\)、在圆上：\(y\_A=(12-8)/2=2\)，\(A(2\sqrt{2},2)\)。代入 \(x^{2}=2py\)：\(8=4p\)、\(p=2\)，抛物线为 \(x^{2}=4y\)。 \((2)\) 先给切线式并验证：点 \(M(x_{1},x_{1}^{2}/4)\) 处，直线 \(y=(x_{1}/2)x-x_{1}^{2}/4\) 与 \(x^{2}=4y\) 联立得 \(x^{2}-2x_{1}x+x_{1}^{2}=(x-x_{1})^{2}=0\)，唯一公共点（重根 \(x_{1}\)），即为过 \(M\) 的切线 \(l_{1}\)；同理 \(N(x_{2},x_{2}^{2}/4)\) 处切线 \(l_{2}\)：\(y=(x_{2}/2)x-x_{2}^{2}/4\)。两切线联立：\((x_{1}-x_{2})x/2=(x_{1}^{2}-x_{2}^{2})/4 \Rightarrow  x_{0}=(x_{1}+x_{2})/2\)，\(y_{0}=(x_{1}/2)x_{0}-x_{1}^{2}/4=x_{1}x_{2}/4\)。设 \(l\)：\(y=kx+1\)（过 \(F(0,1)\)；\(l\) 为 \(x=0\) 时与抛物线仅一交点，不合「不同两点」），代入 \(x^{2}=4y\)：\(x^{2}-4kx-4=0 \Rightarrow  x_{1}x_{2}=-4\)，故 \(y_{0}=-1\) 恒成立，点 \(P\) 恒在定直线 \(y=-1\)（准线）上。\(\rule{1.05ex}{1.05ex}\)〔前引注：源解以导数 \(y'=x/2\) 写切线（系选修\(2\) 工具），本区消前引改「重根验证」路径，切线式与结论同源一致；\((2)\) 源答案「证明见解析」，定直线 \(y=-1\) 系亲算②\(-1\) 落值（⑤\(-5\)）〕}
\end{ansblock}

\jietitle{2.8①　压轴综合一（课时17·拓区）}
% 课时17：卷④2.8.12.1-25 让位无块无键，注记随 17-03（豁免登记在案）。章末拓展 19 席（2章-拓-章末-01~19）缺料未收编——残余项。
\zux{拓展·焦点弦长选项辨析}
\tuhao{01} \tieside{难}已知斜率为 \(1\) 的直线 \(l\) 过椭圆 \(C\)：\(x^{2}/4+y^{2}=1\) 的右焦点，交椭圆于 \(A\)，\(B\) 两点，则弦 \(AB\) 的长为（　）
\lxopt{\(A\)．〔图嵌，数值未回〕}
\lxopt{\(B\)．〔图嵌，数值未回〕}
\lxopt{\(C\)．〔图嵌，数值未回〕}
\lxopt{\(D\)．\(13/5\)}
\begin{ansblock}[2章-拓-课时17-01]
% ans:2章-拓-课时17-01
\ansitem{01}{\(C\)（弦长 \(8/5\)）}
\ansnote{详解}{\(a^{2}=4\)、\(b^{2}=1\)，\(c^{2}=a^{2}-b^{2}=3\)、\(c=\sqrt{3}\)，右焦点 \((\sqrt{3},0)\)。直线 \(l\)：\(y=x-\sqrt{3}\)。代入 \(x^{2}/4+y^{2}=1\)：\(x^{2}/4+(x-\sqrt{3})^{2}=1\)，乘 \(4\) 得 \(x^{2}+4(x^{2}-2\sqrt{3}x+3)=4\)，即 \(5x^{2}-8\sqrt{3}x+8=0\)。\(\Delta =(8\sqrt{3})^{2}-4\cdot 5\cdot 8=192-160=32>0\)，\(x_{1}+x_{2}=8\sqrt{3}/5\)、\(x_{1}x_{2}=8/5\)。弦长 \(|AB|=\sqrt{1+k^{2}}\cdot |x_{1}-x_{2}|=\sqrt{2}\cdot \sqrt{\Delta }/5=\sqrt{2}\cdot \sqrt{32}/5=\sqrt{2}\cdot 4\sqrt{2}/5=8/5\)。（互核：\(|x_{1}-x_{2}|^{2}=(x_{1}+x_{2})^{2}-4x_{1}x_{2}=192/25-32/5=192/25-160/25=32/25\)，\(|AB|^{2}=2\cdot 32/25=64/25\)，\(|AB|=8/5\) 。）故选 \(C\)。〔卷④自带详解读数校订：源【分析】给 \(a^{2}=4\)、\(b^{2}=1\)、\(c=\sqrt{3}\)、\(y=x-\sqrt{3}\)、\(\Delta =32\)、\(|AB|=\sqrt{2}\cdot \sqrt{}\Delta /5=8/5\)、选\(C\)；源【详解】与【大招指引】用「硬解定理」名目（\(\Delta =4B^{2}mn(A^{2}m+B^{2}n-C^{2})\)、\(|AB|=\sqrt{A^{2}+B^{2}}\cdot \sqrt{}\Delta /(|B||A^{2}m+B^{2}n|)\)），按 §五 名目纪律改为常规联立韦达书写，读数不变；【题后反思】栏之常规法即为上引推导〕〔装配注：题面椭圆方程与选项 \(A\)、\(B\)、\(C\) 数值在原 \(dump\) 为【图】占位，椭圆 \(x^{2}/4+y^{2}=1\) 按源【详解】「由椭圆【图】得 \(a^{2}=4,b^{2}=1\)」补全，选项 \(A\)、\(B\)、\(C\) 数值未回——装配回源图或按读数 \(8/5\) 重排选项（列瑕疵⑤）〕}
\end{ansblock}

\zux{拓展·弦长定值反求斜率}
\tuhao{02} \tieside{难}已知直线 \(l: y=kx+1\) 与椭圆 \(x^{2}/2+y^{2}=1\) 交于 \(M\)，\(N\) 两点，且 \(|MN|=4\sqrt{2}/3\)，则 \(k=\)\kongda{\(\pm 1\)}．
\begin{ansblock}[2章-拓-课时17-02]
% ans:2章-拓-课时17-02
\ansitem{02}{\(\pm 1\)}
\ansnote{详解}{联立消 \(y\)：\(x^{2}/2+(kx+1)^{2}=1\)，乘 \(2\) 得 \(x^{2}+2(k^{2}x^{2}+2kx+1)=2\)，即 \((1+2k^{2})x^{2}+4kx=0\)。二次项系数 \(1+2k^{2}>0\) 恒成立。设 \(M(x_{1},y_{1})\)、\(N(x_{2},y_{2})\)，则 \(x_{1}+x_{2}=-4k/(1+2k^{2})\)、\(x_{1}x_{2}=0\)（事实上直线过定点 \((0,1)\)，即椭圆上顶点，一交点恒为它）。由 \(|MN|=4\sqrt{2}/3\) 得 \((x_{1}-x_{2})^{2}+(y_{1}-y_{2})^{2}=32/9\)，即 \((1+k^{2})(x_{1}-x_{2})^{2}=32/9\)；而 \((x_{1}-x_{2})^{2}=(x_{1}+x_{2})^{2}-4x_{1}x_{2}=16k^{2}/(1+2k^{2})^{2}\)，代入得 \((1+k^{2})\cdot 16k^{2}/(1+2k^{2})^{2}=32/9\)，约 \(16\) 并交叉相乘：\(9k^{2}(1+k^{2})=2(1+2k^{2})^{2}\)，展开 \(9k^{2}+9k^{4}=2+8k^{2}+8k^{4}\)，即 \(k^{4}+k^{2}-2=0\)，\((k^{2}+2)(k^{2}-1)=0\)，故 \(k^{2}=1\)、\(k=\pm 1\)。（回验 \(k=1\)：\((1+2)x^{2}+4x=0 \rightarrow  x=0\)、\(-4/3\)，\(|MN|=\sqrt{2}\cdot 4/3=4\sqrt{2}/3\) 。）〔卷④自带详解照录校订（源解同路，得 \(k^{4}+k^{2}-2=0\)、\(k=\pm 1\)）〕}
\end{ansblock}

\zux{拓展·无公共点求k范围}
\tuhao{03} \tieside{简}如果直线 \(y=kx-1\) 与双曲线 \(x^{2}-y^{2}=4\) 没有公共点，求 \(k\) 的取值范围．
\liubai[16mm]{解答书写区}
\begin{ansblock}[2章-拓-课时17-03]
% ans:2章-拓-课时17-03
\ansitem{03}{\(k<-\sqrt{5}/2\) 或 \(k>\sqrt{5}/2\)（\(|k|>\sqrt{5}/2\)）}
\ansnote{详解}{代入 \(x^{2}-(kx-1)^{2}=4\)，展开 \(x^{2}-k^{2}x^{2}+2kx-1=4\)，整理为 \((1-k^{2})x^{2}+2kx-5=0\)。 ① \(1-k^{2}=0\)（\(k=\pm 1\)）：降为一次方程 \(2kx-5=0\)，有解 \(x=5/(2k)\)，即有公共点，不合「没有公共点」。 ② \(1-k^{2}\neq 0\)：无公共点 \(\Leftrightarrow  \Delta <0\)，\(\Delta =(2k)^{2}-4(1-k^{2})(-5)=4k^{2}+20(1-k^{2})=20-16k^{2}\)。令 \(20-16k^{2}<0\) 得 \(k^{2}>5/4\)，即 \(|k|>\sqrt{5}/2\)。（此时 \(1-k^{2}\neq 0\) 自动满足，因 \(5/4>1\)。） 故 \(k\) 的取值范围为 \(k<-\sqrt{5}/2\) 或 \(k>\sqrt{5}/2\)。（几何校：双曲线 \(x^{2}/4-y^{2}/4=1\) 的渐近线斜率 \(\pm 1\)，直线恒过 \((0,-1)\)；\(|k|>\sqrt{5}/2\approx 1.118>1\) 属比渐近线陡、且与实轴交点在双曲线顶点内侧的情形，无交点自洽。）〔亲算印证：§一 依据格读数 \(|k|>\sqrt{5}/2\)，一致；教材答案条 \(raw\) 乱码段未取〕}
\end{ansblock}

\zux{拓展·抛物线双切圆方程}
\tuhao{04} \tieside{简}求圆心在抛物线 \(y^{2}=2x\) 上且与 \(x\) 轴和该抛物线的准线都相切的圆的方程．
\liubai[16mm]{解答书写区}
\begin{ansblock}[2章-拓-课时17-04]
% ans:2章-拓-课时17-04
\ansitem{04}{\((x-1/2)^{2}+(y-1)^{2}=1\) 或 \((x-1/2)^{2}+(y+1)^{2}=1\)（两圆）}
\ansnote{详解}{\(y^{2}=2x\) 中 \(2p=2\)、\(p=1\)，准线 \(l: x=-1/2\)。设圆心 \(C(a,b)\) 在抛物线上，则 \(b^{2}=2a\)（故 \(a\geqslant 0\)）、半径 \(r\)。与 \(x\) 轴相切 \(\Rightarrow  r=|b|\)；与准线 \(x=-1/2\) 相切 \(\Rightarrow  r=a+1/2\)（\(a\geqslant 0\) 故 \(a-(-1/2)=a+1/2\)）。两式相等：\(|b|=a+1/2\)，以 \(a=b^{2}/2\) 代得 \(|b|=b^{2}/2+1/2\)，即 \(|b|^{2}-2|b|+1=0\)，\((|b|-1)^{2}=0\)，故 \(|b|=1\)、\(b=\pm 1\)，进而 \(a=1/2\)、\(r=1\)。所以圆有两解：\((x-1/2)^{2}+(y-1)^{2}=1\) 与 \((x-1/2)^{2}+(y+1)^{2}=1\)。（回验：圆心 \((1/2,1)\) 在抛物线上 \(1=2\cdot 1/2\) ；到 \(x\) 轴距离 \(1=r\) ；到准线距离 \(1/2+1/2=1=r\) 。）〔亲算印证：双解结构与 §四 表「两圆」读数一致〕}
\end{ansblock}

\zux{拓展·焦点弦向量数量积}
\tuhao{05} \tieside{简}已知抛物线 \(y^{2}=4x\) 与过其焦点的斜率为 \(1\) 的直线交于 \(A\)，\(B\) 两点，\(O\) 为坐标原点，求 \(OA\cdot OB\)．
\liubai[16mm]{解答书写区}
\begin{ansblock}[2章-拓-课时17-05]
% ans:2章-拓-课时17-05
\ansitem{05}{\(-3\)}
\ansnote{详解}{焦点 \(F(1,0)\)，直线 \(l: y=x-1\)。代入 \((x-1)^{2}=4x\) 得 \(x^{2}-2x+1=4x\)，即 \(x^{2}-6x+1=0\)（\(\Delta =36-4=32>0\)）。设 \(A(x_{1},y_{1})\)、\(B(x_{2},y_{2})\)，则 \(x_{1}+x_{2}=6\)、\(x_{1}x_{2}=1\)。又 \(y_{1}y_{2}=(x_{1}-1)(x_{2}-1)=x_{1}x_{2}-(x_{1}+x_{2})+1=1-6+1=-4\)。故 \(OA\cdot OB=x_{1}x_{2}+y_{1}y_{2}=1+(-4)=-3\)。（互核：以 \(y\) 为主元，\(x=y+1\) 代 \(y^{2}=4(y+1)\) 得 \(y^{2}-4y-4=0\)，\(y_{1}y_{2}=-4\) ；\(x_{1}x_{2}=(y_{1}+1)(y_{2}+1)=y_{1}y_{2}+(y_{1}+y_{2})+1=-4+4+1=1\) 。）〔亲算印证：§一 依据格与 §四 表读数 \(-3\)，一致；向量数量积系选修\(1\) 已学内容，不触前引闸〕}
\end{ansblock}

\zux{拓展·焦点弦纵积定值证明}
\tuhao{06} \tieside{难}过抛物线 \(y^{2}=2px\)（\(p>0\)）的焦点的一条直线与抛物线相交，且两个交点的纵坐标为 \(y_{1}\)，\(y_{2}\)，求证：\(y_{1}y_{2}=-p^{2}\)．
\liubai[16mm]{解答书写区}
\begin{ansblock}[2章-拓-课时17-06]
% ans:2章-拓-课时17-06
\ansitem{06}{证明见详解}
\ansnote{详解}{焦点 \(F(p/2,0)\)。设过 \(F\) 的直线为 \(x=my+p/2\)：此写法已含垂直 \(x\) 轴的弦（\(m=0\)，即通径），而未含的只有 \(y=0\)（抛物线的对称轴），它与抛物线仅一个公共点，与题设「相交（有两个交点）」矛盾，故不损一般性。代入 \(y^{2}=2px\)：\(y^{2}=2p(my+p/2)=2pmy+p^{2}\)，整理为关于 \(y\) 的一元二次方程 \(y^{2}-2pmy-p^{2}=0\)。其二次项系数为 \(1\neq 0\)，且 \(\Delta =(2pm)^{2}+4p^{2}=4p^{2}(m^{2}+1)>0\)（\(p>0\)），故确两交点，与题设相符。由韦达定理，\(y_{1}y_{2}=-p^{2}/1=-p^{2}\)。\(\rule{1.05ex}{1.05ex}\)（特位核验：\(m=0\) 通径时 \(y^{2}=p^{2}\)、\(y=\pm p\)、积 \(-p^{2}\) ；\(m=1\)、\(p=2\) 时 \(y^{2}-4y-4=0\)、积 \(-4=-p^{2}\) ，与拓\(17-05\) 之 \(y_{1}y_{2}=-4\) 同。）〔亲算印证：定义式直推\(+\)韦达，与课时\(16\) 拓\(16-13/14\) 中「焦点弦两端点纵（横）坐标之积为定值」同族；本条原划批\(E\)、批\(E\) 实取复习题未取，回笼本件（§四 表注）〕}
\end{ansblock}

\zux{拓展·焦点弦长（定义法）}
\tuhao{07} \tieside{难}已知斜率为 \(2\) 的直线 \(AB\) 过抛物线 \(y^{2}=8x\) 的焦点，求弦 \(AB\) 的长．
\liubai[16mm]{解答书写区}
\begin{ansblock}[2章-拓-课时17-07]
% ans:2章-拓-课时17-07
\ansitem{07}{\(10\)}
\ansnote{详解}{\(2p=8\)、\(p=4\)，焦点 \(F(2,0)\)、准线 \(x=-2\)。直线 \(l: y=2(x-2)\)。代入 \(y^{2}=8x\) 得 \(4(x-2)^{2}=8x\)，即 \((x-2)^{2}=2x\)，展开 \(x^{2}-4x+4=2x\)，得 \(x^{2}-6x+4=0\)（\(\Delta =36-16=20>0\)）。设 \(A(x_{1},y_{1})\)、\(B(x_{2},y_{2})\)，则 \(x_{1}+x_{2}=6\)。由抛物线定义（焦半径\(=\)到准线距离，准线 \(x=-2\)）\(|AF|=x_{1}+2\)、\(|BF|=x_{2}+2\)，焦点弦 \(|AB|=|AF|+|BF|=x_{1}+x_{2}+4=6+4=10\)。（互核：以 \(y\) 为主元，\(x=y/2+2\) 代 \(y^{2}=8(y/2+2)\) 得 \(y^{2}-4y-16=0\)，\(y_{1}+y_{2}=4\)、\(y_{1}y_{2}=-16\)，\((y_{1}-y_{2})^{2}=16+64=80\)，\(|AB|=\sqrt{}(1+(1/2)^{2})\cdot |y_{1}-y_{2}|=(\sqrt{5}/2)\cdot 4\sqrt{5}=10\) 。）〔亲算印证：两主元路径互核同为 \(10\)；焦点弦公式 \(|AB|=x_{1}+x_{2}+p=6+4=10\) 同〕}
\end{ansblock}

\zux{拓展·弦长中点与垂直判定}
\tuhao{08} \tieside{难}已知直线 \(l\)：\(y=x-3\) 与抛物线 \(C\)：\(x^{2}=-8y\) 相交于 \(A\)，\(B\) 两点，且 \(O\) 为坐标原点．
\((1)\) 求弦长 \(|AB|\) 以及线段 \(AB\) 的中点坐标；
\((2)\) 判断 \(OA\perp OB\) 是否成立，并说明理由．
\liubai[24mm]{解答书写区}
\begin{ansblock}[2章-拓-课时17-08]
% ans:2章-拓-课时17-08
\ansitem{08}{\((1) |AB|=8\sqrt{5}\)，中点 \((-4,-7)\)；\((2)\) 不成立（\(OA\cdot OB=-15\neq 0\)）}
\ansnote{详解}{\((1)\) 代入 \(x^{2}=-8(x-3)\)：\(x^{2}=-8x+24\)，即 \(x^{2}+8x-24=0\)（\(\Delta =64+96=160>0\)）。设 \(A(x_{1},y_{1})\)、\(B(x_{2},y_{2})\)，则 \(x_{1}+x_{2}=-8\)、\(x_{1}x_{2}=-24\)。弦长 \(|AB|=\sqrt{1+1^{2}}\cdot |x_{1}-x_{2}|=\sqrt{2}\cdot \sqrt{\Delta }=\sqrt{2}\cdot \sqrt{160}=\sqrt{320}=8\sqrt{5}\)。中点 \(x_{0}=(x_{1}+x_{2})/2=-4\)，\(y_{0}=x_{0}-3=-7\)，即 \((-4,-7)\)。（根之实值 \(x=-4\pm 2\sqrt{10}\)，对应 \(y=-7\pm 2\sqrt{10}\)，代 \(x^{2}=-8y\) 校：\((-4\pm 2\sqrt{10})^{2}=16\mp 16\sqrt{10}+40=56\mp 16\sqrt{10}\)，\(-8(-7\pm 2\sqrt{10})=56\mp 16\sqrt{10}\) 。） \((2) y_{1}y_{2}=(x_{1}-3)(x_{2}-3)=x_{1}x_{2}-3(x_{1}+x_{2})+9=-24+24+9=9\)。故 \(OA\cdot OB=x_{1}x_{2}+y_{1}y_{2}=-24+9=-15\neq 0\)，所以 \(OA\) 与 \(OB\) 不垂直，「\(OA\perp OB\)」不成立。（若按三点式复核：取 \(A(-4+2\sqrt{10}, -7+2\sqrt{10})\)、\(B(-4-2\sqrt{10}, -7-2\sqrt{10})\)，两斜率之积 \(k\_OA\cdot k\_OB=(y_{1}y_{2})/(x_{1}x_{2})=9/(-24)=-3/8\neq -1\)，同得否。）〔亲算印证：\((1) 8\sqrt{5}\)、\((-4,-7)\)；\((2) -15\neq 0\) 不垂直，两式自验一致〕}
\end{ansblock}

\zux{拓展·平行渐近线公共点上界}
\tuhao{09} \tieside{难}已知直线 \(l\) 的斜率与双曲线 \(C\) 的渐近线的斜率相等，求证：直线 \(l\) 与双曲线 \(C\) 最多只有一个公共点．
\liubai[16mm]{解答书写区}
\begin{ansblock}[2章-拓-课时17-09]
% ans:2章-拓-课时17-09
\ansitem{09}{证明见详解}
\ansnote{详解}{双曲线分两种标准形，证法同，先取 \(C: x^{2}/a^{2}-y^{2}/b^{2}=1\)（\(a>0,b>0\)），其渐近线为 \(y=\pm (b/a)x\)。题设「\(l\) 的斜率与渐近线的斜率相等」即 \(l\) 与某条渐近线平行或重合，不妨设 \(l: y=(b/a)x+m\)（\(m\) 为任意实数，\(m=0\) 时 \(l\) 就是渐近线本身）。代入 \(C\)：\(x^{2}/a^{2}-((b/a)x+m)^{2}/b^{2}=1\)，而 \(((b/a)x+m)^{2}/b^{2}=(x/a+m/b)^{2}=x^{2}/a^{2}+2mx/(ab)+m^{2}/b^{2}\)，故二次项 \(x^{2}/a^{2}\) 相消，得 \(-2mx/(ab)-m^{2}/b^{2}=1\)，即 \(-(2m/(ab))x=1+m^{2}/b^{2}\)。 ① \(m\neq 0\)：一次项系数 \(-2m/(ab)\neq 0\)，方程有唯一解 \(x=-ab(1+m^{2}/b^{2})/(2m)\)，即 \(l\) 与 \(C\) 恰一个公共点。 ② \(m=0\)：等式左端为 \(0\)、右端为 \(1\)，矛盾，无解，即渐近线与 \(C\) 没有公共点。 两种情形公共点均不超过一个，故「最多只有一个公共点」\(\rule{1.05ex}{1.05ex}\)。 焦点在 \(y\) 轴上时（\(C: y^{2}/a^{2}-x^{2}/b^{2}=1\)，渐近线 \(y=\pm (a/b)x\)）同法：设 \(y=(a/b)x+m\) 代入并以 \(x\) 为主元，二次项 \(x^{2}(a^{2}/b^{2})/a^{2}-x^{2}/b^{2}=x^{2}/b^{2}-x^{2}/b^{2}=0\) 相消，化为一次方程或矛盾式，结论同上。\(\rule{1.05ex}{1.05ex}\)（本条与 \(17-11(a)\)、\(17-14\) 之降次情形同一事理，可作其原理说明。）〔亲算印证：消元降次直推，不引后置结论〕}
\end{ansblock}

\zux{拓展·过定点单公共点直线}
\tuhao{10} \tieside{难}求过点 \(A\)（\(0\)，\(p\)）且与抛物线 \(y^{2}=2px\)（\(p>0\)）只有一个公共点的直线的方程．
\liubai[16mm]{解答书写区}
\begin{ansblock}[2章-拓-课时17-10]
% ans:2章-拓-课时17-10
\ansitem{10}{三条：\(x=0\)；\(y=p\)；\(y=(1/2)x+p\)（即 \(x-2y+2p=0\)）}
\ansnote{详解}{分「斜率不存在」「斜率为 \(0\)」「斜率存在且非 \(0\)」三类穷举。 ① 斜率不存在：直线 \(x=0\)（即 \(y\) 轴），代 \(y^{2}=2p\cdot 0=0\) 得 \(y=0\)，唯一公共点 \((0,0)\)（顶点），符合。 ② 斜率为 \(0\)：直线 \(y=p\)（过 \(A(0,p)\) 且平行对称轴 \(x\) 轴），代 \(p^{2}=2px\) 得 \(x=p/2\)，唯一公共点 \((p/2,p)\)，符合——此即 \(17-02/17-05\) 的退化型「一公共点而不切」。 ③ 斜率存在且非 \(0\)：设 \(y=kx+p\)（\(k\neq 0\)），代 \((kx+p)^{2}=2px\)，得 \(k^{2}x^{2}+(2kp-2p)x+p^{2}=0\)（二次项系数 \(k^{2}\neq 0\)）。只有一个公共点 \(\Leftrightarrow  \Delta =0\)：\(\Delta =(2kp-2p)^{2}-4k^{2}p^{2}=4p^{2}(k-1)^{2}-4k^{2}p^{2}=4p^{2}[(k-1)^{2}-k^{2}]=4p^{2}(1-2k)=0\)，故 \(k=1/2\)，切线为 \(y=(1/2)x+p\)，唯一公共点为 \(x=(2p-2kp)/(2k^{2})=(2p-p)/(1/2)=2p\)、\(y=2p\)，即切点 \((2p,2p)\)（校：\((2p)^{2}=4p^{2}=2p\cdot 2p\) ）。 综上满足条件的直线共三条：\(x=0\)、\(y=p\)、\(y=(1/2)x+p\)。（几何校：\(A(0,p)\) 在抛物线外（\(p^{2}>0=2p\cdot 0\)），故恰一条切线；另两条分别由「与对称轴平行」「与准线同向的竖直割线」两个退化方向提供，条数与 §四 表「三条」读数合。）〔亲算印证：三类穷举\(+\Delta =0\)，读数自验；点 \(A(0,p)\) 恰在准线 \(x=-p/2\) 右侧、且在过顶点之切线（\(y\) 轴）上〕}
\end{ansblock}

\zux{拓展·弦长定值反求直线}
\tuhao{11} \tieside{难}已知斜率为 \(2\) 的直线 \(l\) 与抛物线 \(y^{2}=4x\) 相交于 \(A\)，\(B\) 两点，如果线段 \(AB\) 的长等于 \(5\)，求直线 \(l\) 的方程．
\liubai[16mm]{解答书写区}
\begin{ansblock}[2章-拓-课时17-11]
% ans:2章-拓-课时17-11
\ansitem{11}{\(y=2x-2\)}
\ansnote{详解}{设 \(l: y=2x+m\)。因抛物线以 \(x=y^{2}/4\) 表出，改以 \(y\) 为主元：由 \(y=2x+m\) 得 \(x=(y-m)/2\)，代入 \(y^{2}=4x=(4(y-m))/2=2(y-m)\)，整理为 \(y^{2}-2y+2m=0\)（二次项系数 \(1\neq 0\)）。相交于两点须 \(\Delta =(-2)^{2}-4\cdot 2m=4-8m>0\)，即 \(m<1/2\)；此时 \(y_{1}+y_{2}=2\)、\(y_{1}y_{2}=2m\)，\((y_{1}-y_{2})^{2}=(y_{1}+y_{2})^{2}-4y_{1}y_{2}=4-8m\)。 弦长用 \(y\) 的差：\(|AB|=\sqrt{}(1+(1/k)^{2})\cdot |y_{1}-y_{2}|=\sqrt{1+1/4}\cdot \sqrt{4-8m}=(\sqrt{5}/2)\sqrt{4-8m}\)。令其等于 \(5\)：\(\sqrt{4-8m}=10/\sqrt{5}=2\sqrt{5}\)，平方得 \(4-8m=20\)，故 \(m=-2\)（\(<1/2\)  合相交条件）。 所以直线 \(l\) 的方程为 \(y=2x-2\)。（回代以 \(x\) 为主元复核：\((2x-2)^{2}=4x \rightarrow  4x^{2}-8x+4=4x \rightarrow  x^{2}-3x+1=0\)，\(|x_{1}-x_{2}|=\sqrt{9-4}=\sqrt{5}\)，\(|AB|=\sqrt{1+4}\cdot \sqrt{5}=\sqrt{5}\cdot \sqrt{5}=5\) 。）〔亲算印证：两主元路径互核同得 \(m=-2\)；单解由 \(\Delta >0\) 与平方可逆性保证，无需舍根外的取舍〕}
\end{ansblock}

\zux{拓展·弦中点到准线距离}
\tuhao{12} \tieside{难}过抛物线 \(y^{2}=8x\) 的焦点 \(F\) 的一条直线与此抛物线相交于 \(A\)，\(B\) 两点，已知 \(A\)（\(8\)，\(8\)），求线段 \(AB\) 的中点到抛物线准线的距离．
\liubai[16mm]{解答书写区}
\begin{ansblock}[2章-拓-课时17-12]
% ans:2章-拓-课时17-12
\ansitem{12}{\(25/4\)}
\ansnote{详解}{\(2p=8\)、\(p=4\)，焦点 \(F(2,0)\)，准线 \(x=-2\)。由 \(A(8,8)\)、\(F(2,0)\) 得直线 \(AB\) 的斜率\(=(8-0)/(8-2)=4/3\)，方程 \(y=(4/3)(x-2)\)。代入 \(y^{2}=8x\)：\((16/9)(x-2)^{2}=8x\)，即 \(2(x-2)^{2}=9x\)，展开 \(2x^{2}-8x+8=9x\)，得 \(2x^{2}-17x+8=0\)（\(\Delta =289-64=225=15^{2}\)）。解得 \(x=(17\pm 15)/4\)，即 \(x=8\)（点 \(A\)）或 \(x=1/2\)，故 \(B(1/2, -2)\)（\(y=(4/3)(1/2-2)=(4/3)(-3/2)=-2\)）。 方法一（中点坐标）：中点横坐标 \(x_{0}=(8+1/2)/2=17/4\)，到准线 \(x=-2\) 的距离\(=17/4+2=25/4\)。 方法二（焦半径\(/\)定义）：\(|AF|=x\_A+p/2=8+2=10\)、\(|BF|=x\_B+p/2=1/2+2=5/2\)，焦点弦上「中点到准线的距离\(=(|AF|+|BF|)/2=|AB|/2\)」（梯形中位线，同 \(16\) 片拓\(16-27\) 型），得 \((10+5/2)/2=(25/2)/2=25/4\) 。 （校：\(y_{1}y_{2}=8\cdot (-2)=-16=-p^{2}\)  与 拓\(17-06\) 结论相符。）〔亲算印证：两法同为 \(25/4\)〕}
\end{ansblock}

\zux{拓展·垂直弦长反求直线}
\tuhao{13} \tieside{难}垂直于 \(x\) 轴的直线与抛物线 \(y^{2}=4x\) 交于 \(A\)，\(B\) 两点，且 \(|AB|=4\sqrt{3}\)，求直线 \(AB\) 的方程．
\liubai[16mm]{解答书写区}
\begin{ansblock}[2章-拓-课时17-13]
% ans:2章-拓-课时17-13
\ansitem{13}{\(x=3\)}
\ansnote{详解}{设直线为 \(x=x_{0}\)（\(x_{0}\geqslant 0\)，否则与抛物线无交点）。代入 \(y^{2}=4x_{0}\)，得 \(y=\pm 2\sqrt{x}_{0}\)，两交点 \(A(x_{0},2\sqrt{x}_{0})\)、\(B(x_{0},-2\sqrt{x}_{0})\)（关于 \(x\) 轴对称，故该弦被 \(x\) 轴垂直平分）。\(|AB|=2\cdot 2\sqrt{x}_{0}=4\sqrt{x}_{0}=4\sqrt{3}\)，故 \(\sqrt{x}_{0}=\sqrt{3}\)、\(x_{0}=3\)。所以直线 \(AB\) 的方程为 \(x=3\)。（回验：\(y^{2}=12\)、\(y=\pm 2\sqrt{3}\)、\(|AB|=4\sqrt{3}\) ；此即通径型垂直弦，长 \(4\sqrt{3}\) 大于通径 \(2p=4\)，点 \((3,\pm 2\sqrt{3})\) 在曲线上 。）〔亲算印证：读数 \(x=3\) 自验；§一 依据格「\(2|y|=4\sqrt{3} \Rightarrow  y=\pm 2\sqrt{3}\)、\(x=y^{2}/4=3\)」同〕}
\end{ansblock}

\zux{拓展·MQ平行对称轴证明}
\tuhao{14} \tieside{难}过抛物线的焦点的一条直线与它交于 \(P\)，\(Q\) 两点，过点 \(P\) 和此抛物线顶点的直线与抛物线的准线交于点 \(M\)，求证：直线 \(MQ\) 平行于此抛物线的对称轴．
\liubai[16mm]{解答书写区}
\begin{ansblock}[2章-拓-课时17-14]
% ans:2章-拓-课时17-14
\ansitem{14}{证明见详解}
\ansnote{详解}{取抛物线标准形 \(y^{2}=2px\)（\(p>0\)），顶点 \(O(0,0)\)、焦点 \(F(p/2,0)\)、准线 \(x=-p/2\)、对称轴为 \(x\) 轴。设 \(P(x_{1},y_{1})\)、\(Q(x_{2},y_{2})\)，\(y_{1}y_{2}\neq 0\)（焦点弦的端点不为顶点，且由 拓\(17-06\) 有 \(y_{1}y_{2}=-p^{2}\neq 0\)）。 直线 \(OP\) 的斜率\(=y_{1}/x_{1}\)，而 \(y_{1}^{2}=2px_{1}\) 故 \(x_{1}=y_{1}^{2}/(2p)\)，斜率\(=y_{1}/(y_{1}^{2}/(2p))=2p/y_{1}\)，即 \(OP: y=(2p/y_{1})x\)。 令 \(x=-p/2\)（准线），得 \(M\) 的纵坐标 \(y\_M=(2p/y_{1})(-p/2)=-p^{2}/y_{1}\)。 由 拓\(17-06\)（焦点弦两端点纵坐标之积 \(y_{1}y_{2}=-p^{2}\)）得 \(y_{2}=-p^{2}/y_{1}\)，于是 \(y\_M=y_{2}\)。 \(M(-p/2, y\_M)\) 与 \(Q(x_{2},y_{2})\) 纵坐标相同，而横坐标不同（\(x_{2}\geqslant 0>-p/2\)），故直线 \(MQ\) 为水平线，即 \(MQ\) 平行于抛物线的对称轴 \(x\) 轴。\(\rule{1.05ex}{1.05ex}\) （通径特位：\(PQ\perp x\) 轴时 \(x_{1}=x_{2}=p/2\)、\(y_{2}=-y_{1}\)，\(y\_M=-p^{2}/y_{1}=-(2p\cdot p/2)/y_{1}=-p^{2}/y_{1}\)，由 \(y_{1}^{2}=2p\cdot (p/2)=p^{2}\) 得 \(y\_M=-y_{1}=y_{2}\)  证法不变；焦点在 \(y\) 轴、开口向左等其他标准形同法（换元即可），结论为「\(MQ\) 平行于该抛物线的对称轴」。）〔亲算印证：以 拓\(17-06\) 之焦点弦纵积定值为引，定义与韦达可直推，不引后置结论〕}
\end{ansblock}

\zux{拓展·焦点弦直径圆切准线}
\tuhao{15} \tieside{难}过抛物线的焦点 \(F\) 的一条直线与此抛物线相交于 \(P_{1}\)，\(P_{2}\) 两点．求证：以 \(P_{1}P_{2}\) 为直径的圆与该抛物线的准线相切．
\liubai[16mm]{解答书写区}
\begin{ansblock}[2章-拓-课时17-15]
% ans:2章-拓-课时17-15
\ansitem{15}{证明见详解}
\ansnote{详解}{取标准形 \(y^{2}=2px\)（\(p>0\)），准线 \(l: x=-p/2\)。设 \(P_{1}(x_{1},y_{1})\)、\(P_{2}(x_{2},y_{2})\)，弦 \(P_{1}P_{2}\) 的中点为 \(M\)，则圆心为 \(M\)、半径 \(r=|P_{1}P_{2}|/2\)。 由抛物线定义，焦半径等于点到准线的距离：\(|P_{1}F|=x_{1}+p/2\)、\(|P_{2}F|=x_{2}+p/2\)。因 \(F\) 在线段 \(P_{1}P_{2}\) 上（焦点弦过焦点），\(|P_{1}P_{2}|=|P_{1}F|+|P_{2}F|=x_{1}+x_{2}+p\)，故 \(r=(x_{1}+x_{2}+p)/2\)。 又圆心 \(M\) 的横坐标为 \((x_{1}+x_{2})/2\)，\(M\) 到准线 \(x=-p/2\) 的距离 \(d=(x_{1}+x_{2})/2+p/2=(x_{1}+x_{2}+p)/2\)。 于是 \(d=r\)，即圆心到准线的距离恰等于半径，故以 \(P_{1}P_{2}\) 为直径的圆与准线相切。\(\rule{1.05ex}{1.05ex}\) （另一表述：分别过 \(P_{1}\)、\(P_{2}\)、\(M\) 向准线作垂线，垂足为 \(H_{1}\)、\(H_{2}\)、\(H\)，则 \(|MH|\) 是直角梯形 \(P_{1}H_{1}H_{2}P_{2}\) 的中位线，\(|MH|=(|P_{1}H_{1}|+|P_{2}H_{2}|)/2=(|P_{1}F|+|P_{2}F|)/2=|P_{1}P_{2}|/2=r\)，同得相切——此即 \(17\) 片 \(16-11/\)拓\(16-27\) 反复使用的「中位线\(=\)焦半径和之半」结构。）〔亲算印证：定义\(+\)中位线双表述一致；与 拓\(17-12\) 方法二同族〕}
\end{ansblock}

\zux{拓展·过弦端点切准线之圆}
\tuhao{16} \tieside{难}设抛物线 \(C\)：\(y^{2}=4x\) 的焦点为 \(F\)，过 \(F\) 且斜率为 \(1\) 的直线 \(l\) 与 \(C\) 相交于 \(A\)，\(B\) 两点，求过点 \(A\)，\(B\) 且与 \(C\) 的准线相切的圆的方程．
\liubai[16mm]{解答书写区}
\begin{ansblock}[2章-拓-课时17-16]
% ans:2章-拓-课时17-16
\ansitem{16}{\((x-3)^{2}+(y-2)^{2}=16\) 或 \((x-11)^{2}+(y+6)^{2}=144\)（两圆）}
\ansnote{详解}{焦点 \(F(1,0)\)、准线 \(x=-1\)，直线 \(l: y=x-1\)。代入 \((x-1)^{2}=4x\) 得 \(x^{2}-6x+1=0\)（\(\Delta =32>0\)），设 \(A(x_{1},y_{1})\)、\(B(x_{2},y_{2})\)，则 \(x_{1}+x_{2}=6\)、\(x_{1}x_{2}=1\)，\(x=3\pm 2\sqrt{2}\)。 先定必备数据：弦的中点 \(M=(\)（\(x_{1}+x_{2})/2, (x_{1}+x_{2})/2-1)=(3,2)\)；\(|AB|=\sqrt{2}\cdot |x_{1}-x_{2}|=\sqrt{2}\cdot \sqrt{36-4}=\sqrt{2}\cdot \sqrt{32}=8\)，故 \(|MA|=|MB|=4\)。 所求圆过 \(A\)、\(B\)，其圆心在 \(AB\) 的中垂线上：中垂线斜率 \(-1\)、过 \(M\)，即 \(y=-x+5\)，设圆心 \(C(t, 5-t)\)。 圆的半径按「与准线 \(x=-1\) 相切」为 \(r=|t+1|\)；按「过 \(A\)」为 \(r^{2}=|CA|^{2}=|CM|^{2}+|MA|^{2}=2(t-3)^{2}+16\)（因 \(|CM|^{2}=(t-3)^{2}+(5-t-2)^{2}=2(t-3)^{2}\)）。 令两式相等：\((t+1)^{2}=2(t-3)^{2}+16\)，展开 \(t^{2}+2t+1=2t^{2}-12t+18+16\)，移项得 \(t^{2}-14t+33=0\)，\((t-3)(t-11)=0\)，故 \(t=3\) 或 \(t=11\)（均使 \(r=|t+1|>0\) ，两解皆合）。 \(t=3\)：圆心 \((3,2)=M\)、\(r=4\)，即 \((x-3)^{2}+(y-2)^{2}=16\)——此圆恰为以焦点弦 \(AB\) 为直径的圆，与 拓\(17-15\)（\(B9\) 结论「以焦点弦为直径的圆与准线相切」）相符。 \(t=11\)：圆心 \((11,-6)\)、\(r=12\)，即 \((x-11)^{2}+(y+6)^{2}=144\)。 故满足条件的圆有两个：\((x-3)^{2}+(y-2)^{2}=16\) 与 \((x-11)^{2}+(y+6)^{2}=144\)。（验第二圆：取 \(A(3+2\sqrt{2}, 2+2\sqrt{2})\)，\(|CA|^{2}=(11-3-2\sqrt{2})^{2}+(-6-2-2\sqrt{2})^{2}=(8-2\sqrt{2})^{2}+(-8-2\sqrt{2})^{2}=(72-32\sqrt{2})+(72+32\sqrt{2})=144=r^{2}\) ；圆心到准线距离 \(11+1=12=r\) 。）〔亲算印证：双解 \(t=3\)、\(11\) 系本块自算（教材答案条 \(raw\) 乱码段未取，§三\(.1\)），并以 拓\(17-15\) 之定理与代点计算双向校验〕}
\end{ansblock}

\zux{拓展·单公共点直线计数}
\tuhao{17} \tieside{难}过点 \((0,3)\) 作直线 \(l\)，它与双曲线 \(x^{2}/4-y^{2}/3=1\) 只有一个公共点，这样的直线有（　）
\bindopt {\fontsize{10.5pt}{\qpTJgLeadLiti}\selectfont \makebox[\dimexpr0.25\linewidth-0.25\lxhang-0.25em\relax][l]{\(A\)．\(1\)条}\makebox[\dimexpr0.25\linewidth-0.25\lxhang-0.25em\relax][l]{\(B\)．\(2\)条}\makebox[\dimexpr0.25\linewidth-0.25\lxhang-0.25em\relax][l]{\(C\)．\(3\)条}\makebox[\dimexpr0.25\linewidth-0.25\lxhang-0.25em\relax][l]{\(D\)．\(4\)条}\par}
\begin{ansblock}[2章-拓-课时17-17]
% ans:2章-拓-课时17-17
\ansitem{17}{\(D\)（\(4\)条）}
\ansnote{详解}{① 斜率不存在：\(l\) 为 \(x=0\)（\(y\) 轴），代入得 \(-y^{2}/3=1\) 无实数解，公共点 \(0\) 个，不合。 ② 斜率存在：设 \(l: y=kx+3\)，代入 \(x^{2}/4-y^{2}/3=1\)（乘 \(12\)）：\(3x^{2}-4(kx+3)^{2}=12\)，展开得 \((3-4k^{2})x^{2}-24kx-48=0\)。 \((a) 3-4k^{2}\neq 0\)（\(l\) 不平行渐近线 \(y=\pm (\sqrt{3}/2)x\)）：一个公共点 \(\Leftrightarrow  \Delta =0\)，\(\Delta =(-24k)^{2}+4(3-4k^{2})\cdot 48=576k^{2}+192(3-4k^{2})=576-192k^{2}\)，令其为 \(0\) 得 \(k^{2}=3\)，\(k=\pm \sqrt{3}\)——两条切线，符合。 \((b) 3-4k^{2}=0\)，即 \(k=\pm \sqrt{3}/2\)：方程降为一次 \(-24kx-48=0\)，一次项系数 \(-24k\neq 0\)（\(k=\pm \sqrt{3}/2\) 时），有唯一解 \(x=-2/k\)，即 \(l\) 平行渐近线时与双曲线恰一个公共点——又两条，符合。 合计 \(2+2=4\) 条，选 \(D\)。（校：\(k=\pm \sqrt{3}\) 时 \(3-4k^{2}=3-12=-9\neq 0\)  两类不重；\(k=\pm \sqrt{3}/2\) 与 \(\pm \sqrt{3}\) 互异  四类无重复。）〔\(3\)章件\(11-\#7\) 自带详解照录校订（源解同路：\((3-4k^{2})x^{2}-24kx-48=0\)、\(\Delta =0\) 得 \(k=\pm \sqrt{3}\)、退化一次得 \(k=\pm \sqrt{3}/2\)，共\(4\)条，选\(D\)）；本块补「斜率不存在」与「一次项系数非零」两步穷举〕〔亲算印证：轮\(2\)亲算③难端档「件\(11-7 4\)条，选\(D\)」与源答案一致〕}
\end{ansblock}

\zux{拓展·交右支两点求k范围}
\tuhao{18} \tieside{难}已知直线 \(l\) 的方程为 \(y=kx-1\)，双曲线 \(C\) 的方程为 \(x^{2}-y^{2}=1\)．若直线 \(l\) 与双曲线 \(C\) 的右支相交于不同的两点，则实数 \(k\) 的取值范围是（　）
\bindopt {\fontsize{10.5pt}{\qpTJgLeadLiti}\selectfont \makebox[\dimexpr0.5\linewidth-0.5\lxhang-0.25em\relax][l]{\(A\)．\((-\sqrt{2},\sqrt{2})\)}\makebox[\dimexpr0.5\linewidth-0.5\lxhang-0.25em\relax][l]{\(B\)．\([1,\sqrt{2})\)}\par}
\bindopt {\fontsize{10.5pt}{\qpTJgLeadLiti}\selectfont \makebox[\dimexpr0.5\linewidth-0.5\lxhang-0.25em\relax][l]{\(C\)．\([-\sqrt{2},\sqrt{2}]\)}\makebox[\dimexpr0.5\linewidth-0.5\lxhang-0.25em\relax][l]{\(D\)．\((1,\sqrt{2})\)}\par}
\begin{ansblock}[2章-拓-课时17-18]
% ans:2章-拓-课时17-18
\ansitem{18}{\(D\)}
\ansnote{详解}{代入 \(x^{2}-(kx-1)^{2}=1\)，展开 \(x^{2}-k^{2}x^{2}+2kx-1=1\)，整理为 \((1-k^{2})x^{2}+2kx-2=0\)。设两交点横坐标 \(x_{1}\)、\(x_{2}\)，「同在右支且为不同两点」须同时满足： \((i) 1-k^{2}\neq 0\)（保持二次，否则一次方程只一交点）； \((ii) \Delta =(2k)^{2}-4(1-k^{2})(-2)=4k^{2}+8(1-k^{2})=8-4k^{2}>0 \Leftrightarrow  k^{2}<2\)； \((iii)\) 两根之积 \(x_{1}x_{2}=-2/(1-k^{2})>0 \Leftrightarrow  1-k^{2}<0 \Leftrightarrow  k^{2}>1\)（两根同号）； \((iv)\) 两根之和 \(x_{1}+x_{2}=-2k/(1-k^{2})=2k/(k^{2}-1)>0 \Leftrightarrow  k>0\)（结合 \((iii)\) 同为正，即在右支 \(x\geqslant 1\)）。 合 \((ii)(iii)(iv)\)：\(1<k^{2}<2\) 且 \(k>0\)，得 \(1<k<\sqrt{2}\)，即 \((1,\sqrt{2})\)，选 \(D\)。（校：此范围内 \(1-k^{2}\neq 0\) 自动满足；端点 \(k=1\) 使二次项消失（一交点）、\(k=\sqrt{2}\) 使 \(\Delta =0\)（切，一交点），故两端均开，排除 \(B\) 之左闭。）〔\(3\)章件\(11-\#9\) 自带详解照录校订（源解列 \(\{1-k^{2}\neq 0, \Delta >0,\) 和\(>0,\) 积\(>0\}\) 四条件，解得 \(1<k<\sqrt{2}\)，选\(D\)）〕〔亲算印证：轮\(2\)亲算③难端档「件\(11-9 (1,\sqrt{2})\)，选\(D\)」与源答案一致〕}
\end{ansblock}

\zux{拓展·右支公共点斜率范围}
\tuhao{19} \tieside{简}若曲线 \(2x=\sqrt{4+y^{2}}\) 与直线 \(y=m(x+1)\) 有公共点，则实数 \(m\) 的取值范围是\kongda{\((-2,2)\)}．
\begin{ansblock}[2章-拓-课时17-19]
% ans:2章-拓-课时17-19
\ansitem{19}{\((-2,2)\)}
\ansnote{详解}{曲线侧：\(2x=\sqrt{4+y^{2}}\) 要求 \(x\geqslant 0\)，两边平方得 \(4x^{2}=4+y^{2}\)，即 \(x^{2}-y^{2}/4=1\) 且 \(x\geqslant 0\)（即该双曲线的右支，含右顶点 \((1,0)\)，渐近线 \(y=\pm 2x\)）。直线 \(y=m(x+1)\) 恒过定点 \((-1,0)\)。代入 \(4x^{2}-y^{2}=4\)：\(4x^{2}-m^{2}(x+1)^{2}=4\)，即 \((4-m^{2})x^{2}-2m^{2}x-(m^{2}+4)=0\)。 ① \(|m|<2\)（\(4-m^{2}>0\)）：两根之积 \(=-(m^{2}+4)/(4-m^{2})<0\)，故有一正根；正根满足 \(x^{2}-y^{2}/4=1\) 且 \(x>0\)，即 \(x\geqslant 1\)，落在右支上 \(\Rightarrow \) 有公共点。 ② \(|m|=2\)：方程降为 \(-2m^{2}x-(m^{2}+4)=0\)，解得 \(x=-(m^{2}+4)/(2m^{2})=-8/8=-1<0\)，不在右支 \(\Rightarrow \) 无公共点。 ③ \(|m|>2\)（\(4-m^{2}<0\)）：积 \(>0\) 且和 \(=2m^{2}/(4-m^{2})<0\)，两根同为负 \(\Rightarrow \) 右支无公共点。 故有公共点 \(\Leftrightarrow  |m|<2\)，即 \(m\in (-2,2)\)。〔\(3\)章件\(11-\#10\) 自带详解照录校订（源解为数形结合：右支、渐近线 \(y=\pm 2x\)、直线过 \((-1,0)\)，斜率须夹在两渐近线之间，得 \(-2<m<2\)；本块补代数穷举①②③以去图依赖，读数同）〕〔亲算印证：轮\(2\)亲算④抽样档「件\(11-10 (-2,2)\)」与源答案一致〕}
\end{ansblock}

\zux{拓展·中点弦长求弦}
\tuhao{20} \tieside{难}已知双曲线 \(C: 2x^{2}-y^{2}=2\)，过点 \(P(1,2)\) 的直线 \(l\) 与双曲线 \(C\) 交于 \(M\)、\(N\) 两点，若 \(P\) 为线段 \(MN\) 的中点，则弦长 \(|MN|\) 等于（　）
\bindopt {\fontsize{10.5pt}{\qpTJgLeadLiti}\selectfont \makebox[\dimexpr0.5\linewidth-0.5\lxhang-0.25em\relax][l]{\(A\)．\(4\sqrt{2}/3\)}\makebox[\dimexpr0.5\linewidth-0.5\lxhang-0.25em\relax][l]{\(B\)．\(3\sqrt{3}/4\)}\par}
\bindopt {\fontsize{10.5pt}{\qpTJgLeadLiti}\selectfont \makebox[\dimexpr0.5\linewidth-0.5\lxhang-0.25em\relax][l]{\(C\)．\(4\sqrt{3}\)}\makebox[\dimexpr0.5\linewidth-0.5\lxhang-0.25em\relax][l]{\(D\)．\(4\sqrt{2}\)}\par}
\begin{ansblock}[2章-拓-课时17-20]
% ans:2章-拓-课时17-20
\ansitem{20}{\(D\)}
\ansnote{详解}{若 \(l\) 斜率不存在（\(x=1\)）：代入得 \(2-y^{2}=2\)，\(y=0\)，只一交点，不合「交于两点且 \(P\) 为中点」；故设 \(l: y=k(x-1)+2\)。代入 \(2x^{2}-y^{2}=2\)：\((2-k^{2})x^{2}+2k(k-2)x-(k^{2}-4k+6)=0\)（其中一次项系数由 \(-2k(2-k)\) 化为 \(2k(k-2)\)、常数项为 \(-(2-k)^{2}-2=-(k^{2}-4k+6)\)）。须 \(2-k^{2}\neq 0\)。 中点条件：\(x_{1}+x_{2}=2 \Rightarrow  -2k(k-2)/(2-k^{2})=2 \Rightarrow  -2k^{2}+4k=4-2k^{2} \Rightarrow  4k=4 \Rightarrow  k=1\)（合 \(2-k^{2}=1\neq 0\) ）。 此时方程为 \(x^{2}-2x-3=0\)，\(x_{1}+x_{2}=2\)、\(x_{1}x_{2}=-3\)，两根 \(3\)、\(-1\)，\(\Delta =4+12=16>0\)  两交点 \(M(3,4)\)、\(N(-1,0)\)（校：\(2\cdot 9-16=2\) ；\(2\cdot 1-0=2\) ）。 \(|MN|=\sqrt{1+k^{2}}\cdot |x_{1}-x_{2}|=\sqrt{2}\cdot \sqrt{}((x_{1}+x_{2})^{2}-4x_{1}x_{2})=\sqrt{2}\cdot \sqrt{4+12}=\sqrt{2}\cdot 4=4\sqrt{2}\)，选 \(D\)。〔\(3\)章件\(12-\#2\) 自带详解照录校订（源解同路：设 \(y-2=k(x-1)\)、\(x_{1}+x_{2}=2x\_P\)、解得 \(k=1\)、\(x_{1}x_{2}=-3\)、\(|MN|=\sqrt{2}\cdot \sqrt{16}=4\sqrt{2}\)，选\(D\)；源 \(dump\) 常数项作「\(-(k^(4)-4k+6)\)」系提取层 \(k^{2}\rightarrow k^(4)\) 讹写，本块按 \(-(k^{2}-4k+6)\) 校正并复核：\(k=1\) 时常数项\(=-3\)，与 \(x^{2}-2x-3=0\) 合 ）〕〔亲算印证：轮\(2\)亲算③难端档「件\(12-2 k=1\)、\(|MN|=4\sqrt{2}\)，选\(D\)」与源答案一致〕}
\end{ansblock}

\zux{拓展·直径圆方程求参}
\tuhao{21} \tieside{简}直线 \(l\) 与双曲线 \(x^{2}-y^{2}/2=1\) 交于 \(A\)，\(B\) 两点，以 \(AB\) 为直径的圆 \(C\) 的方程为 \(x^{2}+y^{2}+2x+4y+m=0\)，则 \(m=\)（　）
\bindopt {\fontsize{10.5pt}{\qpTJgLeadLiti}\selectfont \makebox[\dimexpr0.5\linewidth-0.5\lxhang-0.25em\relax][l]{\(A\)．\(-3\)}\makebox[\dimexpr0.5\linewidth-0.5\lxhang-0.25em\relax][l]{\(B\)．\(3\)}\par}
\bindopt {\fontsize{10.5pt}{\qpTJgLeadLiti}\selectfont \makebox[\dimexpr0.5\linewidth-0.5\lxhang-0.25em\relax][l]{\(C\)．\(5-2\sqrt{2}\)}\makebox[\dimexpr0.5\linewidth-0.5\lxhang-0.25em\relax][l]{\(D\)．\(2\sqrt{2}\)}\par}
\begin{ansblock}[2章-拓-课时17-21]
% ans:2章-拓-课时17-21
\ansitem{21}{\(A\)（\(m=-3\)）}
\ansnote{详解}{圆 \(C\) 配方法：\((x+1)^{2}+(y+2)^{2}=5-m\)（须 \(m<5\)），圆心 \((-1,-2)\)、半径 \(r=\sqrt{5-m}\)。\(AB\) 是圆 \(C\) 的直径，故圆心即 \(AB\) 的中点：中点 \((-1,-2)\)。 求弦所在直线（按常规联立韦达，不引中点弦公式之名目）：设 \(l\) 斜率为 \(k\)（若 \(l\) 为 \(x=-1\)，代入双曲线得 \(1-y^{2}/2=1\)、\(y=0\)，只一交点，不合），以 \(u=x+1\) 平移（则 \(x=u-1\)、\(y=ku-2\)，中点条件即 \(u_{1}+u_{2}=0\)）。代入 \(2x^{2}-y^{2}=2\)：\(2(u-1)^{2}-(ku-2)^{2}=2 \rightarrow  2u^{2}-4u+2-k^{2}u^{2}+4ku-4=2\)，即 \((2-k^{2})u^{2}+(4k-4)u-4=0\)。由 \(u_{1}+u_{2}=0\) 得 \(-(4k-4)/(2-k^{2})=0\)，故 \(k=1\)（\(2-k^{2}=1\neq 0\) ）。 于是 \(l: y=(x+1)-2=x-1\)，方程化为 \(u^{2}-4=0\)，\(u=\pm 2 \rightarrow  A(1,0)\)、\(B(-3,-4)\)（校：\(1-0=1\) ；\(9-16/2=1\) ），\(|AB|=\sqrt{4^{2}+4^{2}}=4\sqrt{2}\)。 由直径 \(|AB|=\sqrt{D^{2}+E^{2}-4F}\)（此处 \(D=2\)、\(E=4\)、\(F=m\)）：\(|AB|^{2}=4+16-4m=20-4m\)，令 \((4\sqrt{2})^{2}=32=20-4m\)，解得 \(m=-3\)（\(<5\)  圆确存在）。故选 \(A\)。（互核：\(r=|AB|/2=2\sqrt{2}\)、\(r^{2}=8=5-m \Rightarrow  m=-3\) ；又 \(|CA|^{2}=(1+1)^{2}+(0+2)^{2}=8=r^{2}\) 。）〔\(3\)章件\(12-\#19\) 自带详解照录校订（源解：圆心 \(C(-1,-2)\) 为 \(AB\) 中点、「根据双曲线中点差法的结论 \(k\_AB=(b^{2}/a^{2})\cdot (x_{0}/y_{0})=2\cdot (-1)/(-2)=1\)」、\(l: y=x-1\)、解得 \(A(1,0)\)、\(B(-3,-4)\)、\(|AB|=4\sqrt{2}\)、\(\sqrt{4+16-4m}=4\sqrt{2} \Rightarrow  m=-3\)，选\(A\)；名目「中点差法」按 §五 口径改以平移\(+\)韦达写出，\(k\_AB\) 读数同为 \(1\)）〕〔亲算印证：轮\(2\)亲算④抽样档「件\(12-19 m=-3\)，选\(A\)」与源答案一致；\(0913\) 收口轮补「\(3\)章」前缀（§四 表注），与 \(2\)章件\(12\)《圆与圆的位置关系》\(40\)题件无关〕}
\end{ansblock}

\tailfill
\end{multicols}
\end{document}